% GUID v6, Part I: the proposal. It states the theory and its predictions; Parts II and III test them, Part IV judges them.
\section{DMP for physicists}\label{sec:dmp}

DMP opens with undefined words, then three axioms, then a ground, G0, that is stated and not derived; every later section is a theorem
(T) or a derived constraint (C) that names the axioms it rests on~\cite{dmp}.\footnote{A Lean~4 development, core library only,
models the axioms from a byte-for-byte copy of DMP pinned by hash~\cite{lean4}. It holds 243 theorems, each checked by Lean's kernel and
by an independently written second kernel, among them that the gate of a stack of layers is the meet along its chain~\cite{denning},
that every pass was passed by a minted pair of independent witnesses~\cite{mr}, that adding a layer changes no existing
decision~\cite{ohearn}, that placement never enters a decision~\cite{gm}, and that the potential~\eqref{eq:vibe} never rises under the
system's own refresh. It concerns the axiom system, not the physics: no result of this paper rests on it, and it shows only that the
axioms are consistent in at least one model.} We state the parts this paper uses, quoting DMP's sentences
where the wording matters, and give each a formal counterpart \tagD.

% GUID v6, Primer P1 (label box:dmp): DMP and trust systems, with a worked act. Background only; it states no result of the paper.
% Place it at the start of Section sec:dmp (marker %PRIMER:dmp in part1.tex): % GUID v6, Primer P1 (label box:dmp): DMP and trust systems, with a worked act. Background only; it states no result of the paper.
% Place it at the start of Section sec:dmp (marker %PRIMER:dmp in part1.tex): % GUID v6, Primer P1 (label box:dmp): DMP and trust systems, with a worked act. Background only; it states no result of the paper.
% Place it at the start of Section sec:dmp (marker %PRIMER:dmp in part1.tex): \input{primers/p1-dmp}
\begin{primer}{DMP and trust systems}{box:dmp}
\emph{What a trust system is for.} In a computer system an access-control or trust protocol decides, for every request, who may do what
to which information, and it keeps the evidence of what was done. Its classic design rules are to grant the least privilege an act needs,
to demand two independent keys for what matters (separation of privilege), and to keep a complete audit trail~\cite{saltzer1975}. Its
formal tools include lattices that say where information may flow~\cite{denning}, noninterference, which says what one side may learn of
the other~\cite{gm}, and quorums of witnesses that fail independently~\cite{mr}. DMP~\cite{dmp} states such a system as three axioms and a
ground, G0; every later statement is a numbered theorem (T) or derived constraint (C).

\emph{Words.} A \emph{part} is ``a Duality of Trust'': a trust half and an information half. An \emph{act} is what a part does. A
\emph{layer} is where trust over acts is decided; its two \emph{witnesses} are its halves. The \emph{trust half} decides whether the act
may happen, the \emph{information half} whether it is meant and right. Two witnesses are \emph{independent} when they share no member of the
\emph{failure classes} the layer declares (host, key, author, engine, person, time source, and whatever the minting lists). A
\emph{standing} is trust held in the field of Axiom~2. It has a source (the witness that granted it), a depth (its \emph{grant chain} back
to the first pair) and a freshness that decays with time; \emph{staleness} is that decay. A \emph{record} is evidence that outlives its act.

\emph{The gate.} Each witness returns a verdict in the lattice $\textsf{refuse}<\textsf{hold}<\textsf{pass}$, and the gate of layer $L$
is their meet, $\operatorname{gate}_L(a,t)=v_{w_T}(a,t)\wedge v_{w_I}(a,t)$. The stricter answer stands, and the act is performed only on
$\textsf{pass}\wedge\textsf{pass}$.

\emph{The axioms in one line each.}
\begin{itemize}[nosep,leftmargin=*]
\item \emph{Axiom~1, the Duality.} A part acts only while both halves of its layer trust it for that act; neither overrules the other. Its
\emph{balance}: every part works towards the trust its boundary declares.
\item \emph{Axiom~2, Trust.} A witness trusts only on evidence it verified itself. ``Trust is gravitational. It is a field, never a
position.'' A strong enough boundary ``bends trustspace'' so that only entropy leaves.
\item \emph{Axiom~3, Information.} Every grant, use and loss of trust leaves a record that every witness of the act reads in full. Any other
reader sees the \emph{counts} always and the content only by a \emph{grant}; from trust, nothing about information can be derived but its
existence. Every layer is a tree rooted at its own Duality. The \emph{layout}, where an item sits in that tree, is ``the median of that
demand'', the place of fewest hops in total, recomputed by both halves and never by one.
\end{itemize}

\emph{A worked act} (Figure~\ref{fig:act}). A part $p$ asks to change one entry of a directory. The trust half checks $p$'s standing: it
was granted by a witness, lies two grants from the minted pair and is still fresh, so the verdict is \textsf{pass}. The information half
finds the change well formed but not yet confirmed: \textsf{hold}. The gate gives $\textsf{pass}\wedge\textsf{hold}=\textsf{hold}$, and
nothing is done. Once the change is confirmed, both return \textsf{pass} and the act is performed. A record is written at the act, and both
witnesses can read it in full. A reader outside the layer, say an auditor, sees the counts (here, one more act of this layer) but not the
new entry, unless it is granted to him. From the trust side alone one learns only that the entry exists.

\emph{Containment and G0.} A boundary \emph{contains} information when no part beyond it is trusted for any act on that information.
G0: ``The dual gate cannot witness its own first instance: somebody has to trust the first two witnesses.'' The first act, the
\emph{minting}, makes at least five parts the holders of the \emph{root keys}; it is stated, not derived, and every later standing
descends from it through grant chains.

\emph{The theorems this paper cites.}
\begin{itemize}[nosep,leftmargin=*]
\item \emph{T4}: every parameter is configuration. DMP fixes structure, never a value, and so derives no constant.
\item \emph{T27}: a contained item is destroyed at its scheduled date; with T31 it describes the trap.
\item \emph{T31}: contained information reaches no part by any act. What still leaves is minimal, thermal emission, the \emph{entropy
floor}, which ``carries entropy and no information''. Outside, one learns only the size, the staleness and the rotation interval.
\item \emph{T32}: the VIBE field. Each member has a depth, a mass (one plus the members deriving through it) and a shortfall; its
potential $\Phi_{\mathrm V}\ge0$ never rises under the system's own acts and vanishes in the \emph{balanced state}. Readers other than the
witnesses see one number, the boundary's temperature, its mass-weighted staleness.
\item \emph{T37}: information that no admission names is contained at the item itself.
\item \emph{T40}: a minting originates only from G0 together with a state in which every axiom and theorem provably holds.
\end{itemize}
\end{primer}

\begin{primer}{DMP and trust systems}{box:dmp}
\emph{What a trust system is for.} In a computer system an access-control or trust protocol decides, for every request, who may do what
to which information, and it keeps the evidence of what was done. Its classic design rules are to grant the least privilege an act needs,
to demand two independent keys for what matters (separation of privilege), and to keep a complete audit trail~\cite{saltzer1975}. Its
formal tools include lattices that say where information may flow~\cite{denning}, noninterference, which says what one side may learn of
the other~\cite{gm}, and quorums of witnesses that fail independently~\cite{mr}. DMP~\cite{dmp} states such a system as three axioms and a
ground, G0; every later statement is a numbered theorem (T) or derived constraint (C).

\emph{Words.} A \emph{part} is ``a Duality of Trust'': a trust half and an information half. An \emph{act} is what a part does. A
\emph{layer} is where trust over acts is decided; its two \emph{witnesses} are its halves. The \emph{trust half} decides whether the act
may happen, the \emph{information half} whether it is meant and right. Two witnesses are \emph{independent} when they share no member of the
\emph{failure classes} the layer declares (host, key, author, engine, person, time source, and whatever the minting lists). A
\emph{standing} is trust held in the field of Axiom~2. It has a source (the witness that granted it), a depth (its \emph{grant chain} back
to the first pair) and a freshness that decays with time; \emph{staleness} is that decay. A \emph{record} is evidence that outlives its act.

\emph{The gate.} Each witness returns a verdict in the lattice $\textsf{refuse}<\textsf{hold}<\textsf{pass}$, and the gate of layer $L$
is their meet, $\operatorname{gate}_L(a,t)=v_{w_T}(a,t)\wedge v_{w_I}(a,t)$. The stricter answer stands, and the act is performed only on
$\textsf{pass}\wedge\textsf{pass}$.

\emph{The axioms in one line each.}
\begin{itemize}[nosep,leftmargin=*]
\item \emph{Axiom~1, the Duality.} A part acts only while both halves of its layer trust it for that act; neither overrules the other. Its
\emph{balance}: every part works towards the trust its boundary declares.
\item \emph{Axiom~2, Trust.} A witness trusts only on evidence it verified itself. ``Trust is gravitational. It is a field, never a
position.'' A strong enough boundary ``bends trustspace'' so that only entropy leaves.
\item \emph{Axiom~3, Information.} Every grant, use and loss of trust leaves a record that every witness of the act reads in full. Any other
reader sees the \emph{counts} always and the content only by a \emph{grant}; from trust, nothing about information can be derived but its
existence. Every layer is a tree rooted at its own Duality. The \emph{layout}, where an item sits in that tree, is ``the median of that
demand'', the place of fewest hops in total, recomputed by both halves and never by one.
\end{itemize}

\emph{A worked act} (Figure~\ref{fig:act}). A part $p$ asks to change one entry of a directory. The trust half checks $p$'s standing: it
was granted by a witness, lies two grants from the minted pair and is still fresh, so the verdict is \textsf{pass}. The information half
finds the change well formed but not yet confirmed: \textsf{hold}. The gate gives $\textsf{pass}\wedge\textsf{hold}=\textsf{hold}$, and
nothing is done. Once the change is confirmed, both return \textsf{pass} and the act is performed. A record is written at the act, and both
witnesses can read it in full. A reader outside the layer, say an auditor, sees the counts (here, one more act of this layer) but not the
new entry, unless it is granted to him. From the trust side alone one learns only that the entry exists.

\emph{Containment and G0.} A boundary \emph{contains} information when no part beyond it is trusted for any act on that information.
G0: ``The dual gate cannot witness its own first instance: somebody has to trust the first two witnesses.'' The first act, the
\emph{minting}, makes at least five parts the holders of the \emph{root keys}; it is stated, not derived, and every later standing
descends from it through grant chains.

\emph{The theorems this paper cites.}
\begin{itemize}[nosep,leftmargin=*]
\item \emph{T4}: every parameter is configuration. DMP fixes structure, never a value, and so derives no constant.
\item \emph{T27}: a contained item is destroyed at its scheduled date; with T31 it describes the trap.
\item \emph{T31}: contained information reaches no part by any act. What still leaves is minimal, thermal emission, the \emph{entropy
floor}, which ``carries entropy and no information''. Outside, one learns only the size, the staleness and the rotation interval.
\item \emph{T32}: the VIBE field. Each member has a depth, a mass (one plus the members deriving through it) and a shortfall; its
potential $\Phi_{\mathrm V}\ge0$ never rises under the system's own acts and vanishes in the \emph{balanced state}. Readers other than the
witnesses see one number, the boundary's temperature, its mass-weighted staleness.
\item \emph{T37}: information that no admission names is contained at the item itself.
\item \emph{T40}: a minting originates only from G0 together with a state in which every axiom and theorem provably holds.
\end{itemize}
\end{primer}

\begin{primer}{DMP and trust systems}{box:dmp}
\emph{What a trust system is for.} In a computer system an access-control or trust protocol decides, for every request, who may do what
to which information, and it keeps the evidence of what was done. Its classic design rules are to grant the least privilege an act needs,
to demand two independent keys for what matters (separation of privilege), and to keep a complete audit trail~\cite{saltzer1975}. Its
formal tools include lattices that say where information may flow~\cite{denning}, noninterference, which says what one side may learn of
the other~\cite{gm}, and quorums of witnesses that fail independently~\cite{mr}. DMP~\cite{dmp} states such a system as three axioms and a
ground, G0; every later statement is a numbered theorem (T) or derived constraint (C).

\emph{Words.} A \emph{part} is ``a Duality of Trust'': a trust half and an information half. An \emph{act} is what a part does. A
\emph{layer} is where trust over acts is decided; its two \emph{witnesses} are its halves. The \emph{trust half} decides whether the act
may happen, the \emph{information half} whether it is meant and right. Two witnesses are \emph{independent} when they share no member of the
\emph{failure classes} the layer declares (host, key, author, engine, person, time source, and whatever the minting lists). A
\emph{standing} is trust held in the field of Axiom~2. It has a source (the witness that granted it), a depth (its \emph{grant chain} back
to the first pair) and a freshness that decays with time; \emph{staleness} is that decay. A \emph{record} is evidence that outlives its act.

\emph{The gate.} Each witness returns a verdict in the lattice $\textsf{refuse}<\textsf{hold}<\textsf{pass}$, and the gate of layer $L$
is their meet, $\operatorname{gate}_L(a,t)=v_{w_T}(a,t)\wedge v_{w_I}(a,t)$. The stricter answer stands, and the act is performed only on
$\textsf{pass}\wedge\textsf{pass}$.

\emph{The axioms in one line each.}
\begin{itemize}[nosep,leftmargin=*]
\item \emph{Axiom~1, the Duality.} A part acts only while both halves of its layer trust it for that act; neither overrules the other. Its
\emph{balance}: every part works towards the trust its boundary declares.
\item \emph{Axiom~2, Trust.} A witness trusts only on evidence it verified itself. ``Trust is gravitational. It is a field, never a
position.'' A strong enough boundary ``bends trustspace'' so that only entropy leaves.
\item \emph{Axiom~3, Information.} Every grant, use and loss of trust leaves a record that every witness of the act reads in full. Any other
reader sees the \emph{counts} always and the content only by a \emph{grant}; from trust, nothing about information can be derived but its
existence. Every layer is a tree rooted at its own Duality. The \emph{layout}, where an item sits in that tree, is ``the median of that
demand'', the place of fewest hops in total, recomputed by both halves and never by one.
\end{itemize}

\emph{A worked act} (Figure~\ref{fig:act}). A part $p$ asks to change one entry of a directory. The trust half checks $p$'s standing: it
was granted by a witness, lies two grants from the minted pair and is still fresh, so the verdict is \textsf{pass}. The information half
finds the change well formed but not yet confirmed: \textsf{hold}. The gate gives $\textsf{pass}\wedge\textsf{hold}=\textsf{hold}$, and
nothing is done. Once the change is confirmed, both return \textsf{pass} and the act is performed. A record is written at the act, and both
witnesses can read it in full. A reader outside the layer, say an auditor, sees the counts (here, one more act of this layer) but not the
new entry, unless it is granted to him. From the trust side alone one learns only that the entry exists.

\emph{Containment and G0.} A boundary \emph{contains} information when no part beyond it is trusted for any act on that information.
G0: ``The dual gate cannot witness its own first instance: somebody has to trust the first two witnesses.'' The first act, the
\emph{minting}, makes at least five parts the holders of the \emph{root keys}; it is stated, not derived, and every later standing
descends from it through grant chains.

\emph{The theorems this paper cites.}
\begin{itemize}[nosep,leftmargin=*]
\item \emph{T4}: every parameter is configuration. DMP fixes structure, never a value, and so derives no constant.
\item \emph{T27}: a contained item is destroyed at its scheduled date; with T31 it describes the trap.
\item \emph{T31}: contained information reaches no part by any act. What still leaves is minimal, thermal emission, the \emph{entropy
floor}, which ``carries entropy and no information''. Outside, one learns only the size, the staleness and the rotation interval.
\item \emph{T32}: the VIBE field. Each member has a depth, a mass (one plus the members deriving through it) and a shortfall; its
potential $\Phi_{\mathrm V}\ge0$ never rises under the system's own acts and vanishes in the \emph{balanced state}. Readers other than the
witnesses see one number, the boundary's temperature, its mass-weighted staleness.
\item \emph{T37}: information that no admission names is contained at the item itself.
\item \emph{T40}: a minting originates only from G0 together with a state in which every axiom and theorem provably holds.
\end{itemize}
\end{primer}

% GUID v6, Fig. 3 (label fig:act): one act of DMP. Input in Section sec:dmp next to Primer box:dmp: % GUID v6, Fig. 3 (label fig:act): one act of DMP. Input in Section sec:dmp next to Primer box:dmp: % GUID v6, Fig. 3 (label fig:act): one act of DMP. Input in Section sec:dmp next to Primer box:dmp: \input{figures/fig-act}
\begin{figure}[tbp]
\centering
\begin{tikzpicture}[x=1cm,y=1cm,font=\footnotesize,>={Stealth[length=1.8mm]},
  every node/.style={align=center},
  box/.style={draw,rounded corners=2pt,inner sep=3pt},
  wit/.style={box,fill=black!7,text width=3.6cm,minimum height=1.05cm},
  rdr/.style={box,font=\scriptsize,text width=2.6cm,minimum height=0.85cm},
  lbl/.style={font=\scriptsize,text=black!70},
]
% the layer and its boundary
\draw[rounded corners=6pt,dashed,black!55] (0,-2.6) rectangle (12.6,2.6);
\node[lbl,anchor=north west] at (0.1,2.5) {layer $L$};
% the part
\node[box,text width=1.75cm,minimum height=1.6cm] (p) at (1.2,-0.2) {part $p$\\[1pt]{\scriptsize asks for}\\ act $a$ at $t$};
% the two witnesses
\node[wit] (wt) at (5.1,1.2) {\textbf{trust half} $w_T$\\ {\scriptsize may the act happen?}};
\node[wit] (wi) at (5.1,-1.5) {\textbf{information half} $w_I$\\ {\scriptsize is it meant and right?}};
% the gate
\node[circle,draw,thick,fill=white,minimum size=0.8cm,inner sep=0pt,font=\normalsize] (meet) at (8.05,-0.2) {$\wedge$};
\node[lbl,anchor=east,inner sep=1pt] at (meet.west) {gate};
% the record
\node[box,fill=black!3,text width=2.6cm,minimum height=1.05cm] (rec) at (11.0,-0.2) {\textbf{record}\\ {\scriptsize written at the
  act;}\\[-1pt] {\scriptsize outlives it}};
% requests
\draw[->] ([yshift=0.3cm]p.east) |- (wt.west);
\draw[->] ([yshift=-0.3cm]p.east) |- (wi.west);
% verdicts into the meet
\draw[->] (wt.east) -| node[pos=0.3,above,lbl] {$v_{w_T}$} (meet.north);
\draw[->] (wi.east) -| node[pos=0.3,below,lbl] {$v_{w_I}$} (meet.south);
% pass: the act is performed and recorded
\draw[->,thick] (meet.east) -- node[above,lbl] {\textsf{pass}} (rec.west);
% read in full by both witnesses
\draw[->,black!70] (rec.north) -- (rec.north |- 0,2.15) -- node[pos=0.5,above,lbl] {read in full by both witnesses}
  (wt.north |- 0,2.15) -- (wt.north);
\draw[->,black!70] (rec.south) -- (rec.south |- 0,-2.35) -- (wi.south |- 0,-2.35) -- (wi.south);
% outside readers
\node[rdr] (oth) at (14.55,0.75) {\textbf{any other reader}\\ sees the counts};
\node[rdr,dashed] (gra) at (14.55,-1.15) {\textbf{a granted part}\\ gets the content};
\draw[->,densely dashed] ([yshift=0.25cm]rec.east) to[out=0,in=180] (oth.west);
\draw[->,dotted,thick] ([yshift=-0.25cm]rec.east) to[out=0,in=180] (gra.west);
% verdict lattice
\node[lbl,anchor=north west,align=left,text width=15.6cm,inner sep=0pt] at (0.1,-2.75) {dashed frame: the layer $L$,
  whose boundary contains what no part beyond it is trusted for. Verdicts: \textsf{refuse} $<$ \textsf{hold} $<$ \textsf{pass}; the gate is
  their meet $v_{w_T}\wedge v_{w_I}$, so the stricter answer stands: \textsf{pass} $\wedge$ \textsf{hold} $=$ \textsf{hold}, and nothing
  is done};
\end{tikzpicture}
\caption{One act in DMP (Section~\ref{sec:dmp}; Primer~\ref{box:dmp}). A part $p$ asks for an act $a$. The two halves of its layer, two
independent witnesses, each return a verdict; the gate is their meet, so the act is performed only if both say \textsf{pass}. The act
writes a record that both witnesses read in full. Every other reader sees only the counts (Axiom~3); the content moves only to a part
granted its class. From trust alone nothing about the content can be derived but its existence. Section~\ref{sec:fields} reads the record
as the value of the frame and the counts as energy and momentum.}
\label{fig:act}
\end{figure}

\begin{figure}[tbp]
\centering
\begin{tikzpicture}[x=1cm,y=1cm,font=\footnotesize,>={Stealth[length=1.8mm]},
  every node/.style={align=center},
  box/.style={draw,rounded corners=2pt,inner sep=3pt},
  wit/.style={box,fill=black!7,text width=3.6cm,minimum height=1.05cm},
  rdr/.style={box,font=\scriptsize,text width=2.6cm,minimum height=0.85cm},
  lbl/.style={font=\scriptsize,text=black!70},
]
% the layer and its boundary
\draw[rounded corners=6pt,dashed,black!55] (0,-2.6) rectangle (12.6,2.6);
\node[lbl,anchor=north west] at (0.1,2.5) {layer $L$};
% the part
\node[box,text width=1.75cm,minimum height=1.6cm] (p) at (1.2,-0.2) {part $p$\\[1pt]{\scriptsize asks for}\\ act $a$ at $t$};
% the two witnesses
\node[wit] (wt) at (5.1,1.2) {\textbf{trust half} $w_T$\\ {\scriptsize may the act happen?}};
\node[wit] (wi) at (5.1,-1.5) {\textbf{information half} $w_I$\\ {\scriptsize is it meant and right?}};
% the gate
\node[circle,draw,thick,fill=white,minimum size=0.8cm,inner sep=0pt,font=\normalsize] (meet) at (8.05,-0.2) {$\wedge$};
\node[lbl,anchor=east,inner sep=1pt] at (meet.west) {gate};
% the record
\node[box,fill=black!3,text width=2.6cm,minimum height=1.05cm] (rec) at (11.0,-0.2) {\textbf{record}\\ {\scriptsize written at the
  act;}\\[-1pt] {\scriptsize outlives it}};
% requests
\draw[->] ([yshift=0.3cm]p.east) |- (wt.west);
\draw[->] ([yshift=-0.3cm]p.east) |- (wi.west);
% verdicts into the meet
\draw[->] (wt.east) -| node[pos=0.3,above,lbl] {$v_{w_T}$} (meet.north);
\draw[->] (wi.east) -| node[pos=0.3,below,lbl] {$v_{w_I}$} (meet.south);
% pass: the act is performed and recorded
\draw[->,thick] (meet.east) -- node[above,lbl] {\textsf{pass}} (rec.west);
% read in full by both witnesses
\draw[->,black!70] (rec.north) -- (rec.north |- 0,2.15) -- node[pos=0.5,above,lbl] {read in full by both witnesses}
  (wt.north |- 0,2.15) -- (wt.north);
\draw[->,black!70] (rec.south) -- (rec.south |- 0,-2.35) -- (wi.south |- 0,-2.35) -- (wi.south);
% outside readers
\node[rdr] (oth) at (14.55,0.75) {\textbf{any other reader}\\ sees the counts};
\node[rdr,dashed] (gra) at (14.55,-1.15) {\textbf{a granted part}\\ gets the content};
\draw[->,densely dashed] ([yshift=0.25cm]rec.east) to[out=0,in=180] (oth.west);
\draw[->,dotted,thick] ([yshift=-0.25cm]rec.east) to[out=0,in=180] (gra.west);
% verdict lattice
\node[lbl,anchor=north west,align=left,text width=15.6cm,inner sep=0pt] at (0.1,-2.75) {dashed frame: the layer $L$,
  whose boundary contains what no part beyond it is trusted for. Verdicts: \textsf{refuse} $<$ \textsf{hold} $<$ \textsf{pass}; the gate is
  their meet $v_{w_T}\wedge v_{w_I}$, so the stricter answer stands: \textsf{pass} $\wedge$ \textsf{hold} $=$ \textsf{hold}, and nothing
  is done};
\end{tikzpicture}
\caption{One act in DMP (Section~\ref{sec:dmp}; Primer~\ref{box:dmp}). A part $p$ asks for an act $a$. The two halves of its layer, two
independent witnesses, each return a verdict; the gate is their meet, so the act is performed only if both say \textsf{pass}. The act
writes a record that both witnesses read in full. Every other reader sees only the counts (Axiom~3); the content moves only to a part
granted its class. From trust alone nothing about the content can be derived but its existence. Section~\ref{sec:fields} reads the record
as the value of the frame and the counts as energy and momentum.}
\label{fig:act}
\end{figure}

\begin{figure}[tbp]
\centering
\begin{tikzpicture}[x=1cm,y=1cm,font=\footnotesize,>={Stealth[length=1.8mm]},
  every node/.style={align=center},
  box/.style={draw,rounded corners=2pt,inner sep=3pt},
  wit/.style={box,fill=black!7,text width=3.6cm,minimum height=1.05cm},
  rdr/.style={box,font=\scriptsize,text width=2.6cm,minimum height=0.85cm},
  lbl/.style={font=\scriptsize,text=black!70},
]
% the layer and its boundary
\draw[rounded corners=6pt,dashed,black!55] (0,-2.6) rectangle (12.6,2.6);
\node[lbl,anchor=north west] at (0.1,2.5) {layer $L$};
% the part
\node[box,text width=1.75cm,minimum height=1.6cm] (p) at (1.2,-0.2) {part $p$\\[1pt]{\scriptsize asks for}\\ act $a$ at $t$};
% the two witnesses
\node[wit] (wt) at (5.1,1.2) {\textbf{trust half} $w_T$\\ {\scriptsize may the act happen?}};
\node[wit] (wi) at (5.1,-1.5) {\textbf{information half} $w_I$\\ {\scriptsize is it meant and right?}};
% the gate
\node[circle,draw,thick,fill=white,minimum size=0.8cm,inner sep=0pt,font=\normalsize] (meet) at (8.05,-0.2) {$\wedge$};
\node[lbl,anchor=east,inner sep=1pt] at (meet.west) {gate};
% the record
\node[box,fill=black!3,text width=2.6cm,minimum height=1.05cm] (rec) at (11.0,-0.2) {\textbf{record}\\ {\scriptsize written at the
  act;}\\[-1pt] {\scriptsize outlives it}};
% requests
\draw[->] ([yshift=0.3cm]p.east) |- (wt.west);
\draw[->] ([yshift=-0.3cm]p.east) |- (wi.west);
% verdicts into the meet
\draw[->] (wt.east) -| node[pos=0.3,above,lbl] {$v_{w_T}$} (meet.north);
\draw[->] (wi.east) -| node[pos=0.3,below,lbl] {$v_{w_I}$} (meet.south);
% pass: the act is performed and recorded
\draw[->,thick] (meet.east) -- node[above,lbl] {\textsf{pass}} (rec.west);
% read in full by both witnesses
\draw[->,black!70] (rec.north) -- (rec.north |- 0,2.15) -- node[pos=0.5,above,lbl] {read in full by both witnesses}
  (wt.north |- 0,2.15) -- (wt.north);
\draw[->,black!70] (rec.south) -- (rec.south |- 0,-2.35) -- (wi.south |- 0,-2.35) -- (wi.south);
% outside readers
\node[rdr] (oth) at (14.55,0.75) {\textbf{any other reader}\\ sees the counts};
\node[rdr,dashed] (gra) at (14.55,-1.15) {\textbf{a granted part}\\ gets the content};
\draw[->,densely dashed] ([yshift=0.25cm]rec.east) to[out=0,in=180] (oth.west);
\draw[->,dotted,thick] ([yshift=-0.25cm]rec.east) to[out=0,in=180] (gra.west);
% verdict lattice
\node[lbl,anchor=north west,align=left,text width=15.6cm,inner sep=0pt] at (0.1,-2.75) {dashed frame: the layer $L$,
  whose boundary contains what no part beyond it is trusted for. Verdicts: \textsf{refuse} $<$ \textsf{hold} $<$ \textsf{pass}; the gate is
  their meet $v_{w_T}\wedge v_{w_I}$, so the stricter answer stands: \textsf{pass} $\wedge$ \textsf{hold} $=$ \textsf{hold}, and nothing
  is done};
\end{tikzpicture}
\caption{One act in DMP (Section~\ref{sec:dmp}; Primer~\ref{box:dmp}). A part $p$ asks for an act $a$. The two halves of its layer, two
independent witnesses, each return a verdict; the gate is their meet, so the act is performed only if both say \textsf{pass}. The act
writes a record that both witnesses read in full. Every other reader sees only the counts (Axiom~3); the content moves only to a part
granted its class. From trust alone nothing about the content can be derived but its existence. Section~\ref{sec:fields} reads the record
as the value of the frame and the counts as energy and momentum.}
\label{fig:act}
\end{figure}

\subsection{Words}

The undefined words are \emph{part}, \emph{act}, \emph{witness}, \emph{evidence} and \emph{time}. A part is ``a Duality of Trust: two halves,
one owning trust and one owning information; nothing further is said about what a part is made of''. Trust is a relation $T(w,p,a,t)$:
witness $w$ trusts part $p$ for act $a$ at time $t$. Two witnesses are \emph{independent} when they share no member of the failure classes
declared for their layer (host, key, author, engine, person, time source, and whatever else the minting lists). A \emph{layer} is a place
where trust is decided; its two witnesses are its \emph{halves}, the \emph{trust half} (whether an act may happen) and the
\emph{information half} (whether it is meant and right). A \emph{record} is evidence that outlives its act. A boundary \emph{contains}
information when no part beyond it is trusted for any act on that information, so that no act carries it across.

\begin{definition}[gate]\label{def:gate}
Let $\mathcal V=\{\textsf{refuse}<\textsf{hold}<\textsf{pass}\}$ be the verdict lattice with meet $\wedge$. A layer $L$ with halves $w_T$
(trust) and $w_I$ (information) decides an act $a$ at time $t$ by
\[
\operatorname{gate}_L(a,t)=v_{w_T}(a,t)\wedge v_{w_I}(a,t),
\]
and the act is performed exactly when $\operatorname{gate}_L(a,t)=\textsf{pass}$.
\end{definition}

\subsection{The three axioms}

\paragraph{Axiom 1 (the Duality).} A part acts only while the two halves of its layer, two independent witnesses, ``each with full reach
over the act and gated only by the other'', trust it for that act. ``One half owns trust and the other owns information; neither overrules
the other, and the two keep each other constantly in check and balance.'' Where they disagree the stricter answer stands; every input
gets exactly one answer. \emph{Formally}: Definition~\ref{def:gate}, with $w_T,w_I$ independent; the meet is ``the stricter answer
stands''. The axiom also carries the \emph{balance}: every part works towards the trust its boundary declares, restoring by evidence what
was lost without evidence against it.

\paragraph{Axiom 2 (Trust).} A witness trusts only on evidence it verified itself at the point of use; no information is ever evidence of
itself; the absence of evidence moves nothing. ``Trust is gravitational. It is a field, never a position: every standing has a source
(the witness that granted it), a depth (its grant chain back to the pair) and a freshness that decays.'' A boundary strong enough ``bends
trustspace so that no act carries information across it, and what still leaves is entropy and nothing else''. \emph{Formally}: each
verdict is a function $v_w(a,t)=J_w(E_w(a,t))$ of the evidence $E_w$ that $w$ itself verified at $t$, and a standing verified at $t_v$
with validity $\tau$ has freshness $\varphi(t)=\max\{0,1-(t-t_v)/\tau\}$.

\paragraph{Axiom 3 (Information).} ``Every grant, every use and every loss of trust leaves a record, written at the act, that every witness
of that act can read in full and that outlives the act. Any other reader sees the counts always and the content only by a grant. From trust
nothing can be derived about information but its existence.'' Content moves only by a grant to a part trusted for its class. Information
is hierarchical: every layer has its own Duality and is a tree rooted there; the layout is ``the median of that demand, recomputed by both
halves of the layer and never by one''. The axiom carries the mission ``while maintaining absolute Observability''. \emph{Formally}, the
existence clause is a noninterference condition~\cite{gm}: for every observable $O$ of the trust side and every two contents $x,x'$ that
agree on existence, $O(x)=O(x')$.

\subsection{The ground, containment and the field}

\paragraph{G0.} ``The dual gate cannot witness its own first instance: somebody has to trust the first two witnesses.'' That first act is
the \emph{minting}, in which at least five parts become the holders of the root keys; it is ``stated, not derived'', and every later
standing descends from it. Theorem T40: a minting originates only from G0 together with a state in which every axiom and theorem
provably holds.

\paragraph{Containment and the floor.} T37: information that no admission names is contained at the item itself; ``in the absence of any
external trust, trustspace collapses in on the item''. T31: contained information reaches no part by any act; the substrate still
radiates, and the design makes that emission minimal and \emph{thermal}, ``so it carries entropy and no information'' (the \emph{entropy
floor}); an outside observer learns of the inside ``only its size, its staleness and its rotation interval''.

\paragraph{The VIBE field.} T32 makes Axiom~2's field concrete. Each member $i$ of the system has a \emph{depth} $d_i$ (its shortest grant
chain back to the pair), a \emph{mass} $m_i$ (one plus the number of members deriving through it) and a \emph{shortfall} $s_i$ (what the
boundary declares for it minus what it holds, weighted by freshness). The VIBE field (Variable Influence Balancing Equation) and its
potential are
\begin{equation}\label{eq:vibe}
V_i=m_i\,s_i,\qquad \Phi_{\mathrm V}=\sum_iV_i\ \ge 0,
\end{equation}
with $\Phi_{\mathrm V}=0$ ``exactly when the boundary holds what it declares''. Under the axioms ``the system's own acts never raise the potential;
only losses and ageing do''. Every reader other than the two witnesses reads one number, the boundary's \emph{temperature}, its
mass-weighted staleness. We call $\Phi_{\mathrm V}=0$ the \emph{balanced state}.

\section{The physical reading: two fields}\label{sec:fields}

\subsection{Hypotheses}

Everything physical in this paper rests on steps that no derivation supplies: identifications of a DMP word with a physical quantity,
and a few bridge principles. Table~\ref{tab:hyp} lists all of them with where they are stated and what uses them; every result below
names the ones it uses. This section states the first five and the refinement~1$'$; the others are stated where they are first needed.

\begin{table}[tb]\centering\small
\caption{Every step of the paper that no derivation supplies. ``Adopted'' marks the reading the paper uses where two are given.}
\label{tab:hyp}
\begin{tabularx}{\linewidth}{@{}lLlL@{}}
\toprule
& Statement & Stated in & Used in \\
\midrule
H1 & trust is existence (first reading) & Sec.~\ref{sec:fields} & superseded by 1$'$ \\
H1$'$ & the trust field is the gravitational frame, shaped by the counts (adopted) & Sec.~\ref{sec:fields} & Sections~\ref{sec:fix}--\ref{sec:exp} \\
H2 & what the frame couples to is $T_{\mu\nu}$ & Sec.~\ref{sec:fields} & free fall, Sec.~\ref{sec:fix} \\
H3 & the field is a frame of four vectors (a tetrad) & Sec.~\ref{sec:fields} & no preferred frame, Sec.~\ref{sec:fields} \\
H4 & the information half carries quantum content and a scalar that sets $\alpha$ & Sec.~\ref{sec:fields} & Sections~\ref{sec:fix}, \ref{sec:matter} \\
H5 & a containing boundary is a horizon & Sec.~\ref{sec:fields} & Sections~\ref{sec:bh}, \ref{sec:darkmatter} \\
H6 & the start: inflation with its oscillon decay is the minting & Sec.~\ref{sec:start} & classicality, cosmology \\
H6$'$ & causality begins with the strong split & Sec.~\ref{sec:split} & cosmology, proton decay \\
H7 & remnants of interactions are trapped in the frame & Sec.~\ref{sec:darkmatter} & dark matter \\
H8 & hadrons are binary trees of colour-conjugate halves & Sec.~\ref{sec:trees} & tetraquarks, pentaquarks \\
H9 & a pentaquark is a minting (hidden-flavour pair) & Sec.~\ref{sec:trees} & pentaquarks \\
Balance & DMP's balance holds the trade-off at equality & Sec.~\ref{sec:bridges} & Bridge~B \\
Bridge A & the outermost horizon sets the remnant's scale & Sec.~\ref{sec:darkmatter} & dark matter (refuted) \\
Bridge B & one record per Planck time sets the coupling & Sec.~\ref{sec:bridges} & the Di\'osi--Penrose rate \\
Reach & ``full reach over the act'' means reach over its existence & Sec.~\ref{sec:holo} & Prop.~\ref{th:host}(ii) \\
$B-L$ & $B-L$ is a count, changed only by a recorded act & Sec.~\ref{sec:matter} & Dirac neutrinos \\
\bottomrule
\end{tabularx}
\end{table}

\begin{hypothesis}[trust is existence]\label{h:exist}
DMP's trust half is \emph{existence}: the field of Axiom~2 is the gravitational field, and a standing in it is the existence of what it
stands for. \tagH
\end{hypothesis}

\begin{hypothesis}[existence is measured by energy-momentum]\label{h:energy}
The existence of a body is measured by its energy-momentum tensor $T_{\mu\nu}$: two bodies with the same $T_{\mu\nu}$ have the same
existence, whatever they are made of. \tagH
\end{hypothesis}

\begin{hypothesis}[the gravitational field is a frame]\label{h:frame}
The gravitational field spans every direction of spacetime, three of space and one of time: it is a frame of four vector fields
$e_a{}^{\mu}$ (a tetrad), not a single vector. \tagH
\end{hypothesis}

\begin{hypothesis}[information carries quantum content]\label{h:info}
DMP's information half carries the quantum content of the world, the states of matter and radiation, together with one scalar field
$\psi$ that sets the strength of the electromagnetic coupling, $\alpha=\alpha(\psi)$. \tagH
\end{hypothesis}

\begin{hypothesis}[containment is a horizon]\label{h:horizon}
A containing boundary is an event horizon; the counts every outside reader sees are its asymptotic charges (mass, angular momentum,
electric charge); the entropy floor is the horizon's thermal radiation. \tagH
\end{hypothesis}

\begin{hypothesisprime}[the trust field is the frame; refines Hypothesis~\ref{h:exist}]
The trust field of Axiom~2 (``Trust is gravitational. It is a field'') is the gravitational field, the \emph{frame} of
Hypothesis~\ref{h:frame}. It is separate from the existence and the information it contains and is shaped by them only through the counts
every reader sees (Axiom~3), their energy and momentum; it is never either of them. Its link to G0 is the causal chain of whatever exists.
Existence is what is granted, the rest mass the Higgs vacuum gives (Section~\ref{sec:higgsplit}); information is the content. \tagH
\end{hypothesisprime}

\noindent\emph{How the hypotheses fit together.} Hypotheses~\ref{h:exist} and~\ref{h:energy} are the first reading: gravity is the
existence half itself, and ``existence'' must then mean all energy and momentum, bound as well as granted, because free fall weighs all of
it (Section~\ref{sec:fix}, \chk{Q6}). Hypothesis~1$'$ is the reading this paper adopts. It keeps ``existence'' for what is granted, the
rest mass the Higgs vacuum gives, and makes gravity a separate frame shaped by the counts of everything it contains; Hypothesis~\ref{h:energy}
then becomes a statement about the counts: the frame couples to the energy-momentum tensor $T_{\mu\nu}$ of existence and information
together. Both readings give the same frame, the tetrad of Hypothesis~\ref{h:frame}, and the same propositions: every result of
Section~\ref{sec:fix} keeps its proof with ``the counts $T_{\mu\nu}$'' in place of ``existence'', and two proofs become shorter (the remarks
after Proposition~\ref{pr:wep} and Lemma~\ref{lem:classical}). Under Hypothesis~1$'$ the Higgs field, which grants existence in the narrow
sense, is not the frame. The refinement
keeps the classical-gravity prediction only because the frame is a DMP structure, shaped by the counts; were it outside DMP's structures
altogether, DMP would say nothing about whether gravity is quantum, and the test of Section~\ref{sec:exp} would be lost.

% GUID v6, Primer P3 (label box:gr): general relativity and the tetrad. Background only; it states no result of the paper.
% Place it before "The frame" (marker %PRIMER:gr in part1.tex): % GUID v6, Primer P3 (label box:gr): general relativity and the tetrad. Background only; it states no result of the paper.
% Place it before "The frame" (marker %PRIMER:gr in part1.tex): % GUID v6, Primer P3 (label box:gr): general relativity and the tetrad. Background only; it states no result of the paper.
% Place it before "The frame" (marker %PRIMER:gr in part1.tex): \input{primers/p3-gr}
\begin{primer}{General relativity and the tetrad}{box:gr}
\setlength{\parskip}{3pt plus 1pt}%
\emph{The metric.} General relativity describes gravity as the geometry of spacetime. A metric $g_{\mu\nu}$ gives the interval
$\dd s^2=g_{\mu\nu}\dd x^\mu\dd x^\nu$, with signature $(-,+,+,+)$, and Einstein's equation
$G_{\mu\nu}+\Lambda g_{\mu\nu}=(8\pi G/c^4)\,T_{\mu\nu}$ ties its curvature to the energy-momentum tensor $T_{\mu\nu}$ of matter and
radiation~\cite{wald1984}.

\emph{The tetrad.} A tetrad, or frame, is a set of four orthonormal one-forms $e^a=e^a{}_\mu\dd x^\mu$ ($a=0,\dots,3$) with
$g_{\mu\nu}=\eta_{ab}e^a{}_\mu e^b{}_\nu$ and $\eta=\operatorname{diag}(-1,1,1,1)$. At each point it can be rotated and boosted,
${e^a\to\Lambda^a{}_b(x)\,e^b}$, without changing $g_{\mu\nu}$: this is the \emph{local Lorentz freedom}. Of the tetrad's 16 components, 6 are
that freedom and 10 are the metric. General relativity can be written with the tetrad as its only field~\cite{aldrovandi}. A single unit
timelike vector field added to the metric is different: it picks out a rest frame at every point. That is Einstein--aether
theory~\cite{aether}, and its preferred-frame effects are tightly bounded, for example by the equal speeds of gravitational waves and
light~\cite{gw170817}.

\emph{Free fall.} Matter couples \emph{minimally} when it feels gravity only through $g_{\mu\nu}$, with partial derivatives replaced by
covariant ones. Then $\nabla^\mu T_{\mu\nu}=0$, and a small body follows a geodesic,
$\ddot x^\mu+\Gamma^\mu{}_{\alpha\beta}\dot x^\alpha\dot x^\beta=0$ (dots: derivatives by proper time), whatever it is made of. That is the
weak equivalence principle. Experiments compare the accelerations $a_1$, $a_2$ of two materials through the E\"otv\"os ratio
${\eta=2|a_1-a_2|/|a_1+a_2|}$~\cite{microscope}; with indices, $\eta$ is the Minkowski metric.

\emph{PPN parameters.} In weak fields, metric theories are compared through the parametrised post-Newtonian (PPN) expansion~\cite{will}.
With $U>0$ the Newtonian potential,
\[
g_{00}=-1+\frac{2U}{c^2}-2\beta\frac{U^2}{c^4},\qquad g_{ij}=\Big(1+2\gamma\frac{U}{c^2}\Big)\delta_{ij}.
\]
$\gamma$ says how much space curvature a unit mass produces; light bending and the Shapiro delay scale as $(1+\gamma)/2$. $\beta$ measures
the nonlinearity of gravity and enters the perihelion advance with $\gamma$. General relativity has $\gamma=\beta=1$.

\emph{The lapse.} Split into space and time, spacetime has a \emph{lapse} $N$: the proper time that a static clock records per unit of
coordinate time. In a weak potential $\Phi=-U$ it is $N=\sqrt{1+2\Phi/c^2}$, outside a mass $M$ it is $N=\sqrt{1-2GM/rc^2}$, and it
vanishes at a horizon.

\emph{Black holes.} A stationary black hole is fixed by three numbers: its mass $M$, angular momentum $J$ and charge
$Q$~\cite{israel,kerr}. The rotating hole (Kerr) has its horizon at $r_+=(GM/c^2)\big(1+\sqrt{1-a_*^2}\big)$, with spin
$a_*=cJ/GM^2\le1$~\cite{kerr}. Its \emph{surface gravity} $\kappa$ is the acceleration, as measured from far away, needed to hold a body
at the horizon; for a hole without spin $\kappa=c^4/4GM$. Quantum fields make the horizon radiate at $k_BT_H=\hbar\kappa/2\pi c$~\cite{hawking75}.
The \emph{area theorem}: with positive energy, the horizon area of classical general relativity never decreases~\cite{hawking1971}. It is
the second law of black-hole mechanics~\cite{bch}, with entropy $S=k_Bc^3A/4G\hbar$~\cite{bekenstein73}.
\end{primer}

\begin{primer}{General relativity and the tetrad}{box:gr}
\setlength{\parskip}{3pt plus 1pt}%
\emph{The metric.} General relativity describes gravity as the geometry of spacetime. A metric $g_{\mu\nu}$ gives the interval
$\dd s^2=g_{\mu\nu}\dd x^\mu\dd x^\nu$, with signature $(-,+,+,+)$, and Einstein's equation
$G_{\mu\nu}+\Lambda g_{\mu\nu}=(8\pi G/c^4)\,T_{\mu\nu}$ ties its curvature to the energy-momentum tensor $T_{\mu\nu}$ of matter and
radiation~\cite{wald1984}.

\emph{The tetrad.} A tetrad, or frame, is a set of four orthonormal one-forms $e^a=e^a{}_\mu\dd x^\mu$ ($a=0,\dots,3$) with
$g_{\mu\nu}=\eta_{ab}e^a{}_\mu e^b{}_\nu$ and $\eta=\operatorname{diag}(-1,1,1,1)$. At each point it can be rotated and boosted,
${e^a\to\Lambda^a{}_b(x)\,e^b}$, without changing $g_{\mu\nu}$: this is the \emph{local Lorentz freedom}. Of the tetrad's 16 components, 6 are
that freedom and 10 are the metric. General relativity can be written with the tetrad as its only field~\cite{aldrovandi}. A single unit
timelike vector field added to the metric is different: it picks out a rest frame at every point. That is Einstein--aether
theory~\cite{aether}, and its preferred-frame effects are tightly bounded, for example by the equal speeds of gravitational waves and
light~\cite{gw170817}.

\emph{Free fall.} Matter couples \emph{minimally} when it feels gravity only through $g_{\mu\nu}$, with partial derivatives replaced by
covariant ones. Then $\nabla^\mu T_{\mu\nu}=0$, and a small body follows a geodesic,
$\ddot x^\mu+\Gamma^\mu{}_{\alpha\beta}\dot x^\alpha\dot x^\beta=0$ (dots: derivatives by proper time), whatever it is made of. That is the
weak equivalence principle. Experiments compare the accelerations $a_1$, $a_2$ of two materials through the E\"otv\"os ratio
${\eta=2|a_1-a_2|/|a_1+a_2|}$~\cite{microscope}; with indices, $\eta$ is the Minkowski metric.

\emph{PPN parameters.} In weak fields, metric theories are compared through the parametrised post-Newtonian (PPN) expansion~\cite{will}.
With $U>0$ the Newtonian potential,
\[
g_{00}=-1+\frac{2U}{c^2}-2\beta\frac{U^2}{c^4},\qquad g_{ij}=\Big(1+2\gamma\frac{U}{c^2}\Big)\delta_{ij}.
\]
$\gamma$ says how much space curvature a unit mass produces; light bending and the Shapiro delay scale as $(1+\gamma)/2$. $\beta$ measures
the nonlinearity of gravity and enters the perihelion advance with $\gamma$. General relativity has $\gamma=\beta=1$.

\emph{The lapse.} Split into space and time, spacetime has a \emph{lapse} $N$: the proper time that a static clock records per unit of
coordinate time. In a weak potential $\Phi=-U$ it is $N=\sqrt{1+2\Phi/c^2}$, outside a mass $M$ it is $N=\sqrt{1-2GM/rc^2}$, and it
vanishes at a horizon.

\emph{Black holes.} A stationary black hole is fixed by three numbers: its mass $M$, angular momentum $J$ and charge
$Q$~\cite{israel,kerr}. The rotating hole (Kerr) has its horizon at $r_+=(GM/c^2)\big(1+\sqrt{1-a_*^2}\big)$, with spin
$a_*=cJ/GM^2\le1$~\cite{kerr}. Its \emph{surface gravity} $\kappa$ is the acceleration, as measured from far away, needed to hold a body
at the horizon; for a hole without spin $\kappa=c^4/4GM$. Quantum fields make the horizon radiate at $k_BT_H=\hbar\kappa/2\pi c$~\cite{hawking75}.
The \emph{area theorem}: with positive energy, the horizon area of classical general relativity never decreases~\cite{hawking1971}. It is
the second law of black-hole mechanics~\cite{bch}, with entropy $S=k_Bc^3A/4G\hbar$~\cite{bekenstein73}.
\end{primer}

\begin{primer}{General relativity and the tetrad}{box:gr}
\setlength{\parskip}{3pt plus 1pt}%
\emph{The metric.} General relativity describes gravity as the geometry of spacetime. A metric $g_{\mu\nu}$ gives the interval
$\dd s^2=g_{\mu\nu}\dd x^\mu\dd x^\nu$, with signature $(-,+,+,+)$, and Einstein's equation
$G_{\mu\nu}+\Lambda g_{\mu\nu}=(8\pi G/c^4)\,T_{\mu\nu}$ ties its curvature to the energy-momentum tensor $T_{\mu\nu}$ of matter and
radiation~\cite{wald1984}.

\emph{The tetrad.} A tetrad, or frame, is a set of four orthonormal one-forms $e^a=e^a{}_\mu\dd x^\mu$ ($a=0,\dots,3$) with
$g_{\mu\nu}=\eta_{ab}e^a{}_\mu e^b{}_\nu$ and $\eta=\operatorname{diag}(-1,1,1,1)$. At each point it can be rotated and boosted,
${e^a\to\Lambda^a{}_b(x)\,e^b}$, without changing $g_{\mu\nu}$: this is the \emph{local Lorentz freedom}. Of the tetrad's 16 components, 6 are
that freedom and 10 are the metric. General relativity can be written with the tetrad as its only field~\cite{aldrovandi}. A single unit
timelike vector field added to the metric is different: it picks out a rest frame at every point. That is Einstein--aether
theory~\cite{aether}, and its preferred-frame effects are tightly bounded, for example by the equal speeds of gravitational waves and
light~\cite{gw170817}.

\emph{Free fall.} Matter couples \emph{minimally} when it feels gravity only through $g_{\mu\nu}$, with partial derivatives replaced by
covariant ones. Then $\nabla^\mu T_{\mu\nu}=0$, and a small body follows a geodesic,
$\ddot x^\mu+\Gamma^\mu{}_{\alpha\beta}\dot x^\alpha\dot x^\beta=0$ (dots: derivatives by proper time), whatever it is made of. That is the
weak equivalence principle. Experiments compare the accelerations $a_1$, $a_2$ of two materials through the E\"otv\"os ratio
${\eta=2|a_1-a_2|/|a_1+a_2|}$~\cite{microscope}; with indices, $\eta$ is the Minkowski metric.

\emph{PPN parameters.} In weak fields, metric theories are compared through the parametrised post-Newtonian (PPN) expansion~\cite{will}.
With $U>0$ the Newtonian potential,
\[
g_{00}=-1+\frac{2U}{c^2}-2\beta\frac{U^2}{c^4},\qquad g_{ij}=\Big(1+2\gamma\frac{U}{c^2}\Big)\delta_{ij}.
\]
$\gamma$ says how much space curvature a unit mass produces; light bending and the Shapiro delay scale as $(1+\gamma)/2$. $\beta$ measures
the nonlinearity of gravity and enters the perihelion advance with $\gamma$. General relativity has $\gamma=\beta=1$.

\emph{The lapse.} Split into space and time, spacetime has a \emph{lapse} $N$: the proper time that a static clock records per unit of
coordinate time. In a weak potential $\Phi=-U$ it is $N=\sqrt{1+2\Phi/c^2}$, outside a mass $M$ it is $N=\sqrt{1-2GM/rc^2}$, and it
vanishes at a horizon.

\emph{Black holes.} A stationary black hole is fixed by three numbers: its mass $M$, angular momentum $J$ and charge
$Q$~\cite{israel,kerr}. The rotating hole (Kerr) has its horizon at $r_+=(GM/c^2)\big(1+\sqrt{1-a_*^2}\big)$, with spin
$a_*=cJ/GM^2\le1$~\cite{kerr}. Its \emph{surface gravity} $\kappa$ is the acceleration, as measured from far away, needed to hold a body
at the horizon; for a hole without spin $\kappa=c^4/4GM$. Quantum fields make the horizon radiate at $k_BT_H=\hbar\kappa/2\pi c$~\cite{hawking75}.
The \emph{area theorem}: with positive energy, the horizon area of classical general relativity never decreases~\cite{hawking1971}. It is
the second law of black-hole mechanics~\cite{bch}, with entropy $S=k_Bc^3A/4G\hbar$~\cite{bekenstein73}.
\end{primer}

\subsection{The frame}

In the notation of Primer~\ref{box:gr}, the tetrad $e^a=e^a{}_\mu\,\dd x^\mu$ defines the metric,
\begin{equation}\label{eq:tetrad}
g_{\mu\nu}=\eta_{ab}\,e^a{}_\mu\,e^b{}_\nu,\qquad \eta=\operatorname{diag}(-1,1,1,1),
\end{equation}
and the metric fixes the tetrad only up to a local Lorentz transformation. General relativity can be written with the tetrad as its only
field, in the Einstein--Cartan or the teleparallel form, with exactly Einstein's predictions~\cite{aldrovandi} \tagP. Hypothesis~\ref{h:frame}
matters because a single timelike unit vector added to the metric would be Einstein--aether theory~\cite{aether}, whose preferred frame
GW170817 bounds: $-3\times10^{-15}\le c_{\text{gw}}/c-1\le7\times10^{-16}$~\cite{gw170817} \chk{Q17}. A frame of four has no preferred
direction, provided that the records hold it up to local Lorentz transformations, that is, the metric: a record of one particular tetrad
would single out its timelike leg. We read Hypothesis~\ref{h:frame} that way.

\subsection{The information field and the content}

The information half holds every quantum field of matter and radiation, which this paper calls \emph{content}, and one scalar $\psi$ whose
value sets the fine-structure constant. A scalar that multiplies the Maxwell term and leaves the gravitational term alone is Bekenstein's
theory of a varying $\alpha$~\cite{bekenstein82}; with the scalar's energy fed back into cosmology it is the theory of Sandvik, Barrow
and Magueijo~\cite{bsbm} \tagP. The action whose shape the hypotheses fix is
\begin{equation}\label{eq:action}
S=\int\dd^4x\,|e|\left[\frac{c^4}{16\pi G}R(e)-\frac{c^4}{16\pi G}\,\partial_\mu\phi\,\partial^\mu\phi-\frac{1}{4\mu_0}\,B(\phi)\,
F_{\mu\nu}F^{\mu\nu}\right]+S_{\text{content}}[e,\phi,\Psi],
\end{equation}
where $|e|=|\det e^a{}_\mu|=\sqrt{-g}$, $R(e)$ is the curvature scalar of the metric~\eqref{eq:tetrad}, $\phi$ is $\psi$ measured in units of
the reduced Planck mass (so that its kinetic term is normalised like the gravitational one), $\alpha\propto1/B(\phi)$, and $\Psi$ stands for
the quantum fields of matter. We write
\begin{equation}\label{eq:lambda}
\lambda\equiv\frac{\dd\ln\alpha}{\dd\phi}
\end{equation}
for the scalar's coupling. The action is our construction; what the axioms fix about it is the subject of the next section.

\subsection{The dictionary}

\begin{center}\small
\begin{tabularx}{\linewidth}{@{}lLL@{}}
\toprule
DMP & GUID (physics) & Hypothesis \\
\midrule
trust half; the trust field (Axiom 2), first reading & the frame $e^a{}_\mu$; gravity & \ref{h:exist}, \ref{h:frame} \\
information half & quantum content $\Psi$ and the scalar $\phi$ that sets $\alpha$ & \ref{h:info} \\
a standing; mass of a member (T32) & what the frame couples to: $T_{\mu\nu}$ & \ref{h:energy} \\
record, read in full by every witness (Axiom 3) & the value of the frame, held by every system it acts on & \ref{h:exist} \\
counts, seen by every reader (Axiom 3) & asymptotic charges $M$, $J$, $Q$ & \ref{h:horizon} \\
containing boundary; trap (T27, T31) & event horizon & \ref{h:horizon} \\
entropy floor; temperature of a boundary (T31, T32) & thermal radiation of the horizon; Hawking temperature & \ref{h:horizon} \\
G0, the minting & the start: inflation with the decay of its oscillons; the classical era begins as it completes & \ref{h:minting} \\
balanced state $\Phi_{\mathrm V}=0$ (T32) & the vacuum at the minimum of its potential & Section~\ref{sec:matter} \\
the trust field (Axiom 2), adopted reading & the frame, separate from what it contains and shaped by
its counts & 1$'$, \ref{h:frame} \\
existence, in the narrow sense & the rest mass the Higgs vacuum grants & 1$'$; Section~\ref{sec:higgsplit} \\
\bottomrule
\end{tabularx}
\end{center}

\subsection{Names}

Three words carry the account, and we use them in one sense each from here on. \emph{The frame} is the gravitational field, the tetrad
of Hypothesis~\ref{h:frame}, which the trust field of Axiom~2 becomes under Hypothesis~1$'$. \emph{Content} is the quantum state of matter
and radiation, carried by the information half (Hypothesis~\ref{h:info}). \emph{The grant} is the rest mass an elementary particle has by
its coupling to the Higgs vacuum, the narrow sense of ``existence'' in Hypothesis~1$'$. A \emph{record} is evidence that outlives its act
and that the act's witnesses read in full; the \emph{counts} are what every other reader sees of it, which for the frame are energy and
momentum. Where earlier versions spoke of the ``existence frame'' or the ``gravitational plane'', this paper says \emph{the frame}.
\section{What the axioms fix}\label{sec:fix}

\subsection{Free fall is universal}

\begin{proposition}[weak equivalence principle]\label{pr:wep}
Under Hypotheses~\ref{h:exist} and~\ref{h:energy}, two bodies with the same energy-momentum fall alike in the frame, whatever
they are made of.
\end{proposition}

\begin{proof}
Axiom~3 \tagD: from trust nothing can be derived about information but its existence. By Hypothesis~\ref{h:exist} the frame is
the trust half, so no response of the frame to a body can depend on what the body carries (composition, internal state) except through
the body's existence. By Hypothesis~\ref{h:energy} existence is $T_{\mu\nu}$. So the frame couples to a body through $T_{\mu\nu}$ alone:
the coupling is minimal, and then $\nabla^\mu T_{\mu\nu}=0$ makes every small body, whatever it is made of, follow a geodesic of the
metric the frame defines. Bodies that differ in composition but share their initial position and velocity therefore share their
worldline: the weak equivalence principle \tagP.
\end{proof}

\noindent \emph{Under Hypothesis~1$'$} the proof needs no reading of existence: the frame is shaped only by the counts every reader
sees (Axiom~3), which are energy and momentum whatever carries them; with the same geodesic step free fall is universal, and the frame
weighs granted and bound energy alike, as MICROSCOPE requires.

The principle is among the best-tested in physics. The MICROSCOPE satellite compared the free fall of titanium and platinum and found the
E\"otv\"os ratio $\eta=[-1.5\pm2.3\,(\text{stat})\pm1.5\,(\text{syst})]\times10^{-15}$~\cite{microscope}. Two consequences are sharp.

\emph{Gravity weighs all energy.} About 9\% of a nucleon's mass comes from the Higgs-given quark masses; the rest is the energy of quark
motion and of the gluon field~\cite{yang}. A gravity that weighed only the Higgs-given mass would make titanium and platinum fall
differently by $\eta\approx1\times10^{-3}$ \chk{Q6}, more than $10^{11}$ times what MICROSCOPE allows. The frame weighs the binding in full.

\emph{The information scalar may barely pull.} A scalar that sets $\alpha$ changes the electrostatic energy of every nucleus, so each body
of mass $m_A$ carries a scalar charge $\partial\ln m_A/\partial\phi=\lambda f_A$, where $f_A$ is the electrostatic share of its mass. With the
normalisation of~\eqref{eq:action} the exchange of $\phi$ adds to the Newtonian force between $A$ and $B$ a fraction
\begin{equation}\label{eq:fifth}
\frac{F_\phi}{F_N}=2\lambda^2f_Af_B ,
\end{equation}
so two bodies falling towards the Earth (share $f_\oplus$) differ by $\eta=2\lambda^2f_\oplus(f_1-f_2)$ (Dvali and Zaldarriaga~\cite{dz}).
With the Coulomb term of the semi-empirical mass formula, $f_{\text{Ti}}=2.0\times10^{-3}$ and $f_{\text{Pt}}=4.1\times10^{-3}$. The
profile likelihood of MICROSCOPE together with the E\"ot-Wash beryllium-titanium balance~\cite{schlamminger2008} then bounds the coupling
at 95\% confidence by
\begin{equation}\label{eq:lambdabound}
|\lambda|<3.1\times10^{-5}\qquad\chk{Q6},\ \chk{data:D1.1}.
\end{equation}
\emph{Verdict: derived} (from Axiom~3 with Hypotheses~\ref{h:exist},~\ref{h:energy}); \emph{holds} against
MICROSCOPE, E\"ot-Wash and the lunar and planetary tests of Section~\ref{sec:data}.

\subsection{The scalar stays out of gravity: \texorpdfstring{$\gamma=1$}{gamma = 1}}

\begin{proposition}[light bending]\label{pr:gamma}
If the information half cannot change how the frame responds to existence, the parametrised post-Newtonian parameter of the
metric is $\gamma=1$, as in general relativity.
\end{proposition}

\begin{proof}
Axiom~1 \tagD: neither half overrules the other. A scalar that multiplies the curvature term, $F(\phi)R$, sets the local strength of gravity
and so overrules the existence half; Axiom~1 excludes it. What remains is minimal coupling, as in~\eqref{eq:action}: the scalar's energy is
one more source of the frame, negligible in the solar system, and the post-Newtonian expansion is that of general relativity,
$\gamma=1$~\cite{will} \tagP.
\end{proof}

In Brans--Dicke gravity~\cite{bd}, where the scalar does multiply the curvature, $\gamma=(1+\omega)/(2+\omega)$, and the Cassini
measurement $\gamma-1=(2.1\pm2.3)\times10^{-5}$~\cite{cassini} forces $\omega>4\times10^4$ at two standard deviations \chk{Q7}: the
scalar must decouple almost completely. Axiom~1 asks for that decoupling from the start. Two qualifications. The scalar's
composition-dependent force~\eqref{eq:fifth} shifts the $\gamma$ inferred from light bending and ranging by a term of order
$\lambda^2f_\odot f_{\text{body}}<10^{-14}$, far below any measurement. And whether a term $F(\phi)R$ is present depends on the conformal
frame in which the action is written, as Section~\ref{sec:higgsplit} discusses for Higgs inflation; the physical content of the
proposition is the bound~\eqref{eq:lambdabound} together with $\gamma=1$ for the metric, not the form of the action. \emph{Verdict:
derived}; \emph{holds}.

% GUID v6, Primer P2 (label box:qi): quantum information. Background only; it states no result of the paper.
% Place it before "The frame is classical" (marker %PRIMER:qi in part1.tex): % GUID v6, Primer P2 (label box:qi): quantum information. Background only; it states no result of the paper.
% Place it before "The frame is classical" (marker %PRIMER:qi in part1.tex): % GUID v6, Primer P2 (label box:qi): quantum information. Background only; it states no result of the paper.
% Place it before "The frame is classical" (marker %PRIMER:qi in part1.tex): \input{primers/p2-qi}
\begin{primer}{Quantum information}{box:qi}
\emph{States and processes.} A quantum state is a \emph{density matrix} $\rho$: Hermitian, positive ($\rho\succeq0$) and of unit trace.
It is pure when $\rho=\ket\psi\bra\psi$ and mixed otherwise. The state of one part $A$ of a system $AB$ is the partial trace
$\rho_A=\Tr_B\rho_{AB}$~\cite{nielsen2010}. Every physical process is a \emph{completely positive, trace-preserving} (CPTP) map,
$\rho\mapsto\sum_kK_k\rho K_k^\dagger$ with $\sum_kK_k^\dagger K_k=\mathbb 1$; ``completely'' means that the map stays positive when it
acts on one part of a larger system. A measurement with outcomes $k$ returns $k$ with probability $p_k=\Tr(K_k\rho K_k^\dagger)$ and leaves
the state $K_k\rho K_k^\dagger/p_k$.

\emph{Dephasing is reading.} Dephasing in a basis $\{\ket k\}$, $\rho\mapsto\sum_k\ket k\bra k\rho\ket k\bra k$, erases the off-diagonal
elements (the coherences) and keeps the diagonal. It is what a measurement in that basis does to a system when the outcome is written
somewhere else. A system that is dephased in one fixed basis after every interaction behaves as a classical variable.

\emph{Entanglement.} A state of $AB$ is \emph{separable} if it is a mixture of products, $\rho_{AB}=\sum_ip_i\,\rho_A^i\otimes\rho_B^i$,
and \emph{entangled} otherwise~\cite{horodecki2009}. Two numbers measure it for two qubits.
\begin{itemize}[nosep,leftmargin=*]
\item The \emph{negativity} $\mathcal N$ is the magnitude of the sum of the negative eigenvalues of the partial transpose $\rho^{T_B}$
(the transpose taken on $B$ alone)~\cite{vw}. For two qubits $\mathcal N>0$ exactly when the state is entangled~\cite{peres1996,horodecki1996}.
\item The \emph{concurrence} of a pure two-qubit state $a\ket{00}+b\ket{01}+c\ket{10}+d\ket{11}$ is $C=2|ad-bc|$, from 0 for a product to 1
for a maximally entangled state~\cite{wootters1998}. For pure two-qubit states $\mathcal N=C/2$.
\end{itemize}
An \emph{entanglement witness} is an observable $W$ with $\Tr(W\sigma)\ge0$ for every separable $\sigma$; a measured $\Tr(W\rho)<0$
certifies entanglement~\cite{guhne2009}. A given witness misses every entangled state for which $\Tr(W\rho)\ge0$.

\emph{LOCC.} In local operations and classical communication each party acts on its own system and the parties exchange only measurement
outcomes. LOCC maps take separable states to separable states, and no entanglement measure grows under them~\cite{vw,horodecki2009}.
Creating entanglement between two systems therefore needs a carrier that is itself quantum.

\emph{No-cloning and no-broadcasting.} No process copies two pure states that are neither equal nor orthogonal~\cite{wz,dieks}. A family
of states can be \emph{broadcast}, so that each of two readers ends up with the original, only if its members commute pairwise~\cite{barnum}.
A quantity that several readers hold in full is therefore diagonal in one basis, which is the core of quantum Darwinism~\cite{zurek}.

\emph{Worked example: a dephased mediator.} Qubits $A$ and $B$ start in $\ket+=(\ket0+\ket1)/\sqrt2$ and a mediator qubit $M$ in $\ket0$.
They interact only through $M$: a CNOT from $A$ to $M$, a controlled-$Z$ between $M$ and $B$, and the same CNOT again.
\begin{itemize}[nosep,leftmargin=*]
\item \emph{$M$ coherent.} The first CNOT gives $(\ket{0_A0_M}+\ket{1_A1_M})/\sqrt2$. The controlled-$Z$ turns $B$ into
$\ket-=(\ket0-\ket1)/\sqrt2$ in the second branch, and the second CNOT returns $M$ to $\ket0$. What remains is
$(\ket{0}_A\ket{+}_B+\ket{1}_A\ket{-}_B)/\sqrt2$, a maximally entangled state: $C=1$, $\mathcal N=1/2$.
\item \emph{$M$ dephased} in $\{\ket0,\ket1\}$ after the first CNOT. The state becomes the mixture
$\tfrac12(\ket{00}\bra{00}+\ket{11}\bra{11})_{AM}\otimes\ket+\bra+_B$, and the same two steps leave
$\tfrac12(\ket0\bra0_A\otimes\ket+\bra+_B+\ket1\bra1_A\otimes\ket-\bra-_B)$. This is separable: $\mathcal N=0$.
\end{itemize}
The dephased mediator carried one classical bit from $A$ to $B$; that is LOCC, and it correlates the two qubits without entangling them.
\end{primer}

\begin{primer}{Quantum information}{box:qi}
\emph{States and processes.} A quantum state is a \emph{density matrix} $\rho$: Hermitian, positive ($\rho\succeq0$) and of unit trace.
It is pure when $\rho=\ket\psi\bra\psi$ and mixed otherwise. The state of one part $A$ of a system $AB$ is the partial trace
$\rho_A=\Tr_B\rho_{AB}$~\cite{nielsen2010}. Every physical process is a \emph{completely positive, trace-preserving} (CPTP) map,
$\rho\mapsto\sum_kK_k\rho K_k^\dagger$ with $\sum_kK_k^\dagger K_k=\mathbb 1$; ``completely'' means that the map stays positive when it
acts on one part of a larger system. A measurement with outcomes $k$ returns $k$ with probability $p_k=\Tr(K_k\rho K_k^\dagger)$ and leaves
the state $K_k\rho K_k^\dagger/p_k$.

\emph{Dephasing is reading.} Dephasing in a basis $\{\ket k\}$, $\rho\mapsto\sum_k\ket k\bra k\rho\ket k\bra k$, erases the off-diagonal
elements (the coherences) and keeps the diagonal. It is what a measurement in that basis does to a system when the outcome is written
somewhere else. A system that is dephased in one fixed basis after every interaction behaves as a classical variable.

\emph{Entanglement.} A state of $AB$ is \emph{separable} if it is a mixture of products, $\rho_{AB}=\sum_ip_i\,\rho_A^i\otimes\rho_B^i$,
and \emph{entangled} otherwise~\cite{horodecki2009}. Two numbers measure it for two qubits.
\begin{itemize}[nosep,leftmargin=*]
\item The \emph{negativity} $\mathcal N$ is the magnitude of the sum of the negative eigenvalues of the partial transpose $\rho^{T_B}$
(the transpose taken on $B$ alone)~\cite{vw}. For two qubits $\mathcal N>0$ exactly when the state is entangled~\cite{peres1996,horodecki1996}.
\item The \emph{concurrence} of a pure two-qubit state $a\ket{00}+b\ket{01}+c\ket{10}+d\ket{11}$ is $C=2|ad-bc|$, from 0 for a product to 1
for a maximally entangled state~\cite{wootters1998}. For pure two-qubit states $\mathcal N=C/2$.
\end{itemize}
An \emph{entanglement witness} is an observable $W$ with $\Tr(W\sigma)\ge0$ for every separable $\sigma$; a measured $\Tr(W\rho)<0$
certifies entanglement~\cite{guhne2009}. A given witness misses every entangled state for which $\Tr(W\rho)\ge0$.

\emph{LOCC.} In local operations and classical communication each party acts on its own system and the parties exchange only measurement
outcomes. LOCC maps take separable states to separable states, and no entanglement measure grows under them~\cite{vw,horodecki2009}.
Creating entanglement between two systems therefore needs a carrier that is itself quantum.

\emph{No-cloning and no-broadcasting.} No process copies two pure states that are neither equal nor orthogonal~\cite{wz,dieks}. A family
of states can be \emph{broadcast}, so that each of two readers ends up with the original, only if its members commute pairwise~\cite{barnum}.
A quantity that several readers hold in full is therefore diagonal in one basis, which is the core of quantum Darwinism~\cite{zurek}.

\emph{Worked example: a dephased mediator.} Qubits $A$ and $B$ start in $\ket+=(\ket0+\ket1)/\sqrt2$ and a mediator qubit $M$ in $\ket0$.
They interact only through $M$: a CNOT from $A$ to $M$, a controlled-$Z$ between $M$ and $B$, and the same CNOT again.
\begin{itemize}[nosep,leftmargin=*]
\item \emph{$M$ coherent.} The first CNOT gives $(\ket{0_A0_M}+\ket{1_A1_M})/\sqrt2$. The controlled-$Z$ turns $B$ into
$\ket-=(\ket0-\ket1)/\sqrt2$ in the second branch, and the second CNOT returns $M$ to $\ket0$. What remains is
$(\ket{0}_A\ket{+}_B+\ket{1}_A\ket{-}_B)/\sqrt2$, a maximally entangled state: $C=1$, $\mathcal N=1/2$.
\item \emph{$M$ dephased} in $\{\ket0,\ket1\}$ after the first CNOT. The state becomes the mixture
$\tfrac12(\ket{00}\bra{00}+\ket{11}\bra{11})_{AM}\otimes\ket+\bra+_B$, and the same two steps leave
$\tfrac12(\ket0\bra0_A\otimes\ket+\bra+_B+\ket1\bra1_A\otimes\ket-\bra-_B)$. This is separable: $\mathcal N=0$.
\end{itemize}
The dephased mediator carried one classical bit from $A$ to $B$; that is LOCC, and it correlates the two qubits without entangling them.
\end{primer}

\begin{primer}{Quantum information}{box:qi}
\emph{States and processes.} A quantum state is a \emph{density matrix} $\rho$: Hermitian, positive ($\rho\succeq0$) and of unit trace.
It is pure when $\rho=\ket\psi\bra\psi$ and mixed otherwise. The state of one part $A$ of a system $AB$ is the partial trace
$\rho_A=\Tr_B\rho_{AB}$~\cite{nielsen2010}. Every physical process is a \emph{completely positive, trace-preserving} (CPTP) map,
$\rho\mapsto\sum_kK_k\rho K_k^\dagger$ with $\sum_kK_k^\dagger K_k=\mathbb 1$; ``completely'' means that the map stays positive when it
acts on one part of a larger system. A measurement with outcomes $k$ returns $k$ with probability $p_k=\Tr(K_k\rho K_k^\dagger)$ and leaves
the state $K_k\rho K_k^\dagger/p_k$.

\emph{Dephasing is reading.} Dephasing in a basis $\{\ket k\}$, $\rho\mapsto\sum_k\ket k\bra k\rho\ket k\bra k$, erases the off-diagonal
elements (the coherences) and keeps the diagonal. It is what a measurement in that basis does to a system when the outcome is written
somewhere else. A system that is dephased in one fixed basis after every interaction behaves as a classical variable.

\emph{Entanglement.} A state of $AB$ is \emph{separable} if it is a mixture of products, $\rho_{AB}=\sum_ip_i\,\rho_A^i\otimes\rho_B^i$,
and \emph{entangled} otherwise~\cite{horodecki2009}. Two numbers measure it for two qubits.
\begin{itemize}[nosep,leftmargin=*]
\item The \emph{negativity} $\mathcal N$ is the magnitude of the sum of the negative eigenvalues of the partial transpose $\rho^{T_B}$
(the transpose taken on $B$ alone)~\cite{vw}. For two qubits $\mathcal N>0$ exactly when the state is entangled~\cite{peres1996,horodecki1996}.
\item The \emph{concurrence} of a pure two-qubit state $a\ket{00}+b\ket{01}+c\ket{10}+d\ket{11}$ is $C=2|ad-bc|$, from 0 for a product to 1
for a maximally entangled state~\cite{wootters1998}. For pure two-qubit states $\mathcal N=C/2$.
\end{itemize}
An \emph{entanglement witness} is an observable $W$ with $\Tr(W\sigma)\ge0$ for every separable $\sigma$; a measured $\Tr(W\rho)<0$
certifies entanglement~\cite{guhne2009}. A given witness misses every entangled state for which $\Tr(W\rho)\ge0$.

\emph{LOCC.} In local operations and classical communication each party acts on its own system and the parties exchange only measurement
outcomes. LOCC maps take separable states to separable states, and no entanglement measure grows under them~\cite{vw,horodecki2009}.
Creating entanglement between two systems therefore needs a carrier that is itself quantum.

\emph{No-cloning and no-broadcasting.} No process copies two pure states that are neither equal nor orthogonal~\cite{wz,dieks}. A family
of states can be \emph{broadcast}, so that each of two readers ends up with the original, only if its members commute pairwise~\cite{barnum}.
A quantity that several readers hold in full is therefore diagonal in one basis, which is the core of quantum Darwinism~\cite{zurek}.

\emph{Worked example: a dephased mediator.} Qubits $A$ and $B$ start in $\ket+=(\ket0+\ket1)/\sqrt2$ and a mediator qubit $M$ in $\ket0$.
They interact only through $M$: a CNOT from $A$ to $M$, a controlled-$Z$ between $M$ and $B$, and the same CNOT again.
\begin{itemize}[nosep,leftmargin=*]
\item \emph{$M$ coherent.} The first CNOT gives $(\ket{0_A0_M}+\ket{1_A1_M})/\sqrt2$. The controlled-$Z$ turns $B$ into
$\ket-=(\ket0-\ket1)/\sqrt2$ in the second branch, and the second CNOT returns $M$ to $\ket0$. What remains is
$(\ket{0}_A\ket{+}_B+\ket{1}_A\ket{-}_B)/\sqrt2$, a maximally entangled state: $C=1$, $\mathcal N=1/2$.
\item \emph{$M$ dephased} in $\{\ket0,\ket1\}$ after the first CNOT. The state becomes the mixture
$\tfrac12(\ket{00}\bra{00}+\ket{11}\bra{11})_{AM}\otimes\ket+\bra+_B$, and the same two steps leave
$\tfrac12(\ket0\bra0_A\otimes\ket+\bra+_B+\ket1\bra1_A\otimes\ket-\bra-_B)$. This is separable: $\mathcal N=0$.
\end{itemize}
The dephased mediator carried one classical bit from $A$ to $B$; that is LOCC, and it correlates the two qubits without entangling them.
\end{primer}

\subsection{The frame is classical}

We need two theorems of quantum information.

\begin{definition}[broadcasting]
A family of states $\{\rho_s\}$ on a system $S$ can be \emph{broadcast} if there is one physical process (a completely positive,
trace-preserving map) $\mathcal B:S\to S_1\otimes S_2$ such that both reduced states of $\mathcal B(\rho_s)$ equal $\rho_s$, for every $s$.
For pure states and a product output this is \emph{cloning}.
\end{definition}

\begin{lemma}[no-cloning; no-broadcasting]\label{lem:nobroadcast}
(i) A unitary that copies two pure states $\ket s,\ket t$, $U\ket s\ket0=\ket s\ket s$ and $U\ket t\ket0=\ket t\ket t$, exists only if
$\langle s|t\rangle\in\{0,1\}$ (Wootters and Zurek~\cite{wz}; Dieks~\cite{dieks}). (ii) A family of states can be broadcast if and only if
its members commute pairwise (Barnum, Caves, Fuchs, Jozsa and Schumacher~\cite{barnum}). \tagP
\end{lemma}

\begin{proof}[Proof of (i)]
Unitaries preserve inner products, so $\langle s|t\rangle=\langle s|t\rangle^2$, whose only solutions are $0$ and $1$. For (ii) see~\cite{barnum};
the ``if'' direction is the copier in the common eigenbasis, $\ket i\to\ket i\ket i$, which \chk{Q3} runs: a diagonal state is broadcast,
$\ket+$ is not (each reader is left with $I/2$).
\end{proof}

\begin{lemma}[the frame is classical]\label{lem:classical}
Under Hypothesis~\ref{h:exist}, every value the frame takes after the first recorded act belongs to a pairwise commuting family:
the frame admits no superposition of two of its values.
\end{lemma}

\begin{proof}
Axiom~3 \tagD: every grant, use and loss of trust leaves a record that every witness of the act can read in full, and any other reader sees
the counts always. Axiom~2 \tagD: every standing of the trust field is evidence held in such records. Axiom~1 \tagD: there are at least two
independent witnesses. So every value of the trust field is held in full by at least two independent readers: it is broadcast. By
Lemma~\ref{lem:nobroadcast}(ii) the values commute. By Hypothesis~\ref{h:exist} they are the values of the frame. Writing the
value into records that are read in full also dephases the frame in the record basis after every act, which is the hypothesis of
Lemma~\ref{lem:locc} below.
\end{proof}

\begin{remark}
The lemma is close to definitional: DMP demands absolute observability of its shared structure, and the no-broadcasting theorem states the
price. It is not a discovery about nature; it is the statement of which side of an open question the axioms take. It has a precise physical
counterpart in quantum Darwinism: a quantity of which the environment holds many redundant records is objective and classical, and generic
quantum dynamics makes exactly the pointer observables so recorded~\cite{zurek,bph}. DMP's ``every witness reads in full'' is that
redundancy, imposed as an axiom. Read in the other direction, the lemma is a design statement for information systems: a control structure
that every party can audit in full is classical, while the data it governs need not be.
\end{remark}

\noindent \emph{Under Hypothesis~1$'$} the lemma applies to the frame directly: the frame depends on the parts only through the
counts, which every reader holds in full, so from the parts' point of view it is a measure-and-prepare channel (Lemma~\ref{lem:locc}) and
classical.

\subsection{A classical mediator cannot entangle}

\begin{lemma}[classical mediation is LOCC]\label{lem:locc}
Let $A$ and $B$ be quantum systems that interact only through a mediator $M$, by alternating local interactions $A$--$M$ and $M$--$B$. If
$M$ is classical in the sense that it is dephased in a fixed basis $\{\ket k\}$ after each interaction, then the induced map on $AB$ is a
local operation with classical communication, and it maps every separable state of $AB$ to a separable state. (A diagonal marginal of $M$
alone is not enough: a coherent mediator entangled with $A$ can have one.)
\end{lemma}

\begin{proof}
Dephasing $M$ in its basis is equivalent to reading it there. An $A$--$M$ interaction followed by that reading is an instrument on $A$: it
takes $\rho_A$ to the ensemble $\{p_k,\rho_A^{(k)}\}$ with the outcome $k$ written in $M$. An $M$--$B$ interaction with $M$ in $\ket k$ is a
channel $\mathcal F_k$ on $B$. One round therefore maps $\rho_{AB}$ to $\sum_k(\mathcal E_k\otimes\mathcal F_k)(\rho_{AB})$ with
$\mathcal E_k$ completely positive and $\sum_k\mathcal E_k$ trace preserving: a one-way LOCC map. Further rounds compose LOCC maps, which
remain LOCC. LOCC maps take separable states to separable states, and no entanglement measure increases under them~\cite{vw} \tagP.
\end{proof}

Galley, Giacomini and Selby prove the same statement in the generality of probabilistic theories: a mediator that entangles two quantum
systems by local interaction cannot be classical~\cite{ggs}. Check \chk{Q4} runs Lemma~\ref{lem:locc} on qubits: three random local
couplings $A$--$M$, $M$--$B$, $A$--$M$ leave $A$ and $B$ entangled in most of 300 trials when $M$ is coherent (median negativity $0.12$)
and never when $M$ is dephased after each coupling (largest negativity below $10^{-12}$).

\subsection{Where classicality begins: the start}\label{sec:start}

G0 is stated, not derived, and the minting is the first act: before it there is no act and therefore no record \tagD.
Lemma~\ref{lem:classical} rests on records, so it applies only once acts exist. GUID describes the beginning as one event, the
\emph{start}. For anything to begin to exist it must be possible for it to exist, and the first step of existing is to inflate, through the
decay of the oscillons that coalesce once the initial energy has been emitted. Inflation and that decay cannot be separated or ordered,
because there is no causality yet in which to order them; together they are the start.

\begin{hypothesis}[the start]\label{h:minting}
Inflation and the decay of the oscillons it ends in are one event, the start, which is the minting. It has no internal order the axioms can
refer to, and the classical era of Lemma~\ref{lem:classical}, with causality in DMP's sense (acts between existing parts), begins as it
completes. \tagH
\end{hypothesis}

The physics of the start is standard \tagP. Inflation ends as the inflaton oscillates about the minimum of its potential; in a broad class
of models self-resonance fragments it into oscillons, long-lived localised lumps of the field, which dominate the universe for a while
(Amin et al.~\cite{amin}); their decay into the particles of the Standard Model reheats the universe (Kofman, Linde and
Starobinsky~\cite{kls}). In physics these are stages of one field's history; GUID's statement that they form one event agrees with that,
since the axioms, which need acts to order anything, have nothing by which to order them. Within the start the axioms place no restriction
on the geometry, which may be in superposition; after it, every geometry the axioms define is recorded and classical. This matches the
standard semiclassical picture, in which the geometry of the universe becomes classical through decoherence at its very beginning
(Halliwell~\cite{halliwell}; Kiefer~\cite{kiefer87}), and the no-boundary proposal, in which there is no ``before'' (Hartle and
Hawking~\cite{hartlehawking}). In GUID the start is therefore the boundary of quantum gravity in time: whatever is quantum about the
geometry itself lies within it, where DMP by its own words makes no statement.

\subsection{DMP can host the proposed account}\label{sec:host}

\begin{definition}[base configuration, host]\label{def:model}
The \emph{base configuration currently assumed for our universe} is what physics supplies and DMP does not derive, the configuration in
DMP's sense (T4): the constants $G$, $\hbar$, $c$, $\Lambda$ and $\alpha$, the fields and parameters of the Standard Model, and the initial
data of the $\Lambda$CDM model. A configuration that obeys DMP \emph{hosts} a theory of gravity if its own shared structures (the trust
field, the layout, the order of acts, the counts) carry the theory's classical degrees of freedom and its information half carries the
theory's quantum content, so that nothing in the axioms prevents the configuration from reproducing what the theory predicts.
\end{definition}

\begin{proposition}[DMP can host the proposed account of gravity]\label{th:host}
Under Hypotheses~\ref{h:exist}--\ref{h:info}, and supplied with the base configuration currently assumed for our universe, DMP hosts the
theory of this paper: a classical frame over quantum content. In it, (i) the frame, carried by DMP's shared structures, mediates
no entanglement between the systems it acts on; (ii) the information half can carry any quantum content.
\end{proposition}

\begin{proof}
(i) DMP's shared structures are broadcast records (Axiom~3), so by Lemma~\ref{lem:classical} the field they carry is classical, whichever
structure carries it, and by Lemma~\ref{lem:locc} a classical mediator creates no entanglement. (ii) Axiom~3 restricts what is
\emph{derived from trust} and how content moves (by grant); it does not require content to be broadcast, so nothing in the axioms forbids
non-commuting states in the information half \tagD. A grant moves content without copying it, as no-cloning requires. The base
configuration supplies the constants and the content that DMP does not derive. Part~(ii) also uses the reading of ``full reach over
the act'' as reach over the act's existence (Section~\ref{sec:holo}), a hypothesis \tagH.
\end{proof}

The proposition is a statement of compatibility: it shows that nothing in the axioms obstructs the account, not that the axioms produce
its dynamics. Two remarks place it. First, the constants: DMP derives none of them, and the proposition needs them supplied, as every physical
theory needs its parameters (Section~\ref{sec:break}). Second, universality: whatever can run a universal quantum computer as content can
also run a theory in which the geometry itself is quantum, by the Church--Turing--Deutsch principle~\cite{deutsch} (see below);
in such a theory the quantum geometry would sit in the information half and would not be DMP's frame. Part~(i) is therefore the proposed
theory's characteristic, falsifiable prediction, and Section~\ref{sec:exp} names the experiments that test it.

\paragraph{The Existence Turing machine.}

Proposition~\ref{th:host}(ii) says that the information half can carry any quantum content. Assembled from DMP parts, a universal
quantum computer can then run as content, and by the Church--Turing--Deutsch principle it approximates every finitely realisable
physical system to any finite precision~\cite{deutsch}; GUID calls such an assembly the \emph{Existence Turing machine}. On a classical
host this costs memory and time of order $2^n$ for $n$ qubits of generic content, so the general route needs a quantum information half;
Section~\ref{sec:hosted} hosts four known quantum configurations classically and recovers their known results. Two consequences matter
for gravity. Every candidate theory of quantum gravity whose predictions are finitely realisable can also run on the machine, as content:
the axioms do not forbid quantum gravity in nature, but its geometry would not be DMP's frame. And the machine's own structure, its gates,
records and layout, is classical: in the proposed theory it is the frame itself (Section~\ref{sec:gravsim}).

\subsection{Gravitational waves and the graviton}

The frame's linearised excitations, $g_{\mu\nu}=\eta_{\mu\nu}+h_{\mu\nu}$, are a massless spin-2 field. In $D$ spacetime dimensions such a
field has $D(D-3)/2$ propagating polarisations, two in 3+1 \chk{Q8}. Weinberg's low-energy theorem~\cite{weinberg64} and Deser's
self-coupling argument~\cite{deser70} show that a massless spin-2 field must couple equally to all energy and, made consistent, is general
relativity \tagP: the spin-2 wave is the frame's own, and its universal coupling is Proposition~\ref{pr:wep} again. Its speed, its
mass and its polarisations are tested in Section~\ref{sec:data}.

By Lemma~\ref{lem:classical}, once the start is complete (Section~\ref{sec:cosmo}) these waves are classical waves of the frame: GUID has no graviton as a quantum
of the frame \tagD. That is consistent with every observation, for reasons the literature has made precise. Dyson showed that detectors of
the LIGO type cannot register single gravitons, because mirrors heavy enough would collapse into a black hole~\cite{dyson}; Rothman and
Boughn found that even a Jupiter-mass absorber near a neutron star would be swamped by neutrinos, and a shield against them would collapse
too~\cite{rb}. Every gravitational wave measured is deep in the classical regime: the energy GW150914 radiated corresponds to more than $10^{78}$
quanta at its frequency \chk{Q9}. Localising any quantum inside its own Schwarzschild radius needs a wavelength below $\sqrt{4\pi}\,L_P$ and
an energy above $\sqrt\pi\,E_P$ \chk{Q9}, the Planck scale at which measurement and gravity cannot be separated~\cite{mead}; a LIGO-band
graviton is $2\times10^{-41}$ of that. The proposals to register single quanta of a gravitational wave as jumps of a cooled resonant
bar~\cite{tobar} are a gravitational photoelectric effect, and like the photoelectric effect they can be explained by a classical wave acting
on a quantised detector~\cite{cdr}. So the absence of the graviton in GUID is not testable by counting gravitons; it is testable only
through the mediation experiments of Section~\ref{sec:exp} and, with Bridge~B, through the decoherence test of
Section~\ref{sec:bridges}. \emph{Verdict: derived}; not testable by counting gravitons, only through those experiments.

\section{How a classical frame can carry quantum matter}\label{sec:cq}

Proposition~\ref{th:host} places GUID among the theories in which spacetime is classical and matter quantum. Such theories have a known
difficulty, and its resolution is a prediction.

\subsection{The semiclassical equation fails}

The simplest coupling, $G_{\mu\nu}=(8\pi G/c^4)\langle\hat T_{\mu\nu}\rangle$ (M{\o}ller~\cite{moller}, Rosenfeld~\cite{rosenfeld63}),
sources gravity by the quantum expectation value. A mass in an equal superposition of two places then attracts a probe towards the midpoint,
where the mass never is. Page and Geilker tested a macroscopic version, with the position of lead masses set by a quantum random event, and
found the probe following the actual masses, not the average~\cite{pagegeilker} \tagP. The equation is also nonlinear in the quantum state,
which in general permits signalling. So a classical frame cannot simply be sourced by expectation values.

% GUID v6, Primer P4 (label box:cq): classical-quantum dynamics and collapse. Background only; it states no result of the paper.
% Place it before "Consistent classical-quantum dynamics" (marker %PRIMER:cq in part1.tex): % GUID v6, Primer P4 (label box:cq): classical-quantum dynamics and collapse. Background only; it states no result of the paper.
% Place it before "Consistent classical-quantum dynamics" (marker %PRIMER:cq in part1.tex): % GUID v6, Primer P4 (label box:cq): classical-quantum dynamics and collapse. Background only; it states no result of the paper.
% Place it before "Consistent classical-quantum dynamics" (marker %PRIMER:cq in part1.tex): \input{primers/p4-cq}
\begin{primer}{Classical-quantum dynamics and collapse}{box:cq}
\setlength{\parskip}{3pt plus 1pt}%
\emph{The classical-quantum (CQ) master equation.} A consistent coupling keeps the joint description linear, positive and
normalised, as probability theory demands. The joint state is a family $\varrho(z)$ of positive operators, one for each value
$z$ of the classical variables (here the frame), normalised as $\int\dd z\,\Tr\varrho(z)=1$~\cite{oppenheim}. For one classical variable
$z$ driven by one quantum observable $L$ (for gravity, the mass density), continuous Markovian dynamics has the form~\cite{osswd}
\[
\partial_t\varrho=-\frac i\hbar[H,\varrho]+D_0\Big(L\varrho L-\tfrac12\{L^2,\varrho\}\Big)-\frac{D_1}2\,\partial_z\{L,\varrho\}
+\frac{D_2}2\,\partial_z^2\varrho ,
\]
with three coefficients beyond quantum mechanics:
\begin{itemize}[nosep,leftmargin=*]
\item $D_0$, \emph{decoherence}: the quantum system loses coherence between states that source the classical field differently (for
gravity, different mass distributions);
\item $D_1$, \emph{back-reaction}: a drift of $z$ driven by the quantum state, that is, how matter pushes the frame;
\item $D_2$, \emph{diffusion}: random kicks of $z$, that is, noise in the frame.
\end{itemize}
Complete positivity requires $D_1^2\le4D_0D_2$ here, and $D_0>0$ wherever $D_1\ne0$; Oppenheim et al.\ give the general matrix form of
this trade-off~\cite[Eq.~(23)]{osswd}, which is Eq.~\eqref{eq:tradeoff} of this paper. In words: \emph{a push implies noise}. If quantum matter pushes a
classical field at all, the matter decoheres and the field diffuses, and the product of the two is bounded below by the square of the push,
$D_0D_2\ge D_1^2/4$.

\emph{Measurement and feedback.} An older family of consistent models makes gravity a channel. Each mass's position, or its mass density,
is weakly and continuously measured, and a classical force computed from the record acts back on the other masses (Kafri, Taylor and
Milburn, KTM~\cite{ktm}; Tilloy and Di\'osi~\cite{tilloydiosi}). Newton's law then holds on average, the measurement decoheres
superpositions, and the feedback adds noise. The masses talk only through a classical record, so the channel is LOCC
(Primer~\ref{box:qi}) and cannot entangle them. A weaker measurement decoheres less but gives a noisier record and so more heating; KTM find
both effects at least of the order of the gravitational coupling. The least decoherence of two masses over all measurement strengths is
what this paper calls the \emph{KTM minimum}, $\omega/2$, where $\omega$ is the rate at which a coherent field would build up the
entangling phase.

\emph{Gravity-related collapse.} Di\'osi and Penrose proposed that a superposition of two mass distributions $\rho_1$, $\rho_2$ decays at
the rate $\lambda=E_\Delta/\hbar$, where
\[
E_\Delta=\frac G2\int\!\!\int\dd^3x\,\dd^3y\,\frac{[\rho_1(\mathbf x)-\rho_2(\mathbf x)][\rho_1(\mathbf y)-\rho_2(\mathbf y)]}
{|\mathbf x-\mathbf y|}
\]
is the gravitational self-energy of their difference~\cite{diosi87,penrose96,bassi2013}. For a body of mass $M$ and size $R$ displaced by
more than its size, $\lambda\approx GM^2/\hbar R$: the \emph{Di\'osi--Penrose (DP) rate}. For point-like nuclei $E_\Delta$ diverges, so the
mass density must be smeared over a length $R_0$. The same noise makes charged matter radiate. An underground germanium detector sees no
such radiation, which excludes smearing at nuclear size and requires $R_0$ of order $10^{-10}$\,m or more~\cite{donadi,figurato}. In the
Newtonian limit of CQ gravity the DP kernel is one admissible choice of the decoherence term~\cite{osswd}.
\end{primer}

\begin{primer}{Classical-quantum dynamics and collapse}{box:cq}
\setlength{\parskip}{3pt plus 1pt}%
\emph{The classical-quantum (CQ) master equation.} A consistent coupling keeps the joint description linear, positive and
normalised, as probability theory demands. The joint state is a family $\varrho(z)$ of positive operators, one for each value
$z$ of the classical variables (here the frame), normalised as $\int\dd z\,\Tr\varrho(z)=1$~\cite{oppenheim}. For one classical variable
$z$ driven by one quantum observable $L$ (for gravity, the mass density), continuous Markovian dynamics has the form~\cite{osswd}
\[
\partial_t\varrho=-\frac i\hbar[H,\varrho]+D_0\Big(L\varrho L-\tfrac12\{L^2,\varrho\}\Big)-\frac{D_1}2\,\partial_z\{L,\varrho\}
+\frac{D_2}2\,\partial_z^2\varrho ,
\]
with three coefficients beyond quantum mechanics:
\begin{itemize}[nosep,leftmargin=*]
\item $D_0$, \emph{decoherence}: the quantum system loses coherence between states that source the classical field differently (for
gravity, different mass distributions);
\item $D_1$, \emph{back-reaction}: a drift of $z$ driven by the quantum state, that is, how matter pushes the frame;
\item $D_2$, \emph{diffusion}: random kicks of $z$, that is, noise in the frame.
\end{itemize}
Complete positivity requires $D_1^2\le4D_0D_2$ here, and $D_0>0$ wherever $D_1\ne0$; Oppenheim et al.\ give the general matrix form of
this trade-off~\cite[Eq.~(23)]{osswd}, which is Eq.~\eqref{eq:tradeoff} of this paper. In words: \emph{a push implies noise}. If quantum matter pushes a
classical field at all, the matter decoheres and the field diffuses, and the product of the two is bounded below by the square of the push,
$D_0D_2\ge D_1^2/4$.

\emph{Measurement and feedback.} An older family of consistent models makes gravity a channel. Each mass's position, or its mass density,
is weakly and continuously measured, and a classical force computed from the record acts back on the other masses (Kafri, Taylor and
Milburn, KTM~\cite{ktm}; Tilloy and Di\'osi~\cite{tilloydiosi}). Newton's law then holds on average, the measurement decoheres
superpositions, and the feedback adds noise. The masses talk only through a classical record, so the channel is LOCC
(Primer~\ref{box:qi}) and cannot entangle them. A weaker measurement decoheres less but gives a noisier record and so more heating; KTM find
both effects at least of the order of the gravitational coupling. The least decoherence of two masses over all measurement strengths is
what this paper calls the \emph{KTM minimum}, $\omega/2$, where $\omega$ is the rate at which a coherent field would build up the
entangling phase.

\emph{Gravity-related collapse.} Di\'osi and Penrose proposed that a superposition of two mass distributions $\rho_1$, $\rho_2$ decays at
the rate $\lambda=E_\Delta/\hbar$, where
\[
E_\Delta=\frac G2\int\!\!\int\dd^3x\,\dd^3y\,\frac{[\rho_1(\mathbf x)-\rho_2(\mathbf x)][\rho_1(\mathbf y)-\rho_2(\mathbf y)]}
{|\mathbf x-\mathbf y|}
\]
is the gravitational self-energy of their difference~\cite{diosi87,penrose96,bassi2013}. For a body of mass $M$ and size $R$ displaced by
more than its size, $\lambda\approx GM^2/\hbar R$: the \emph{Di\'osi--Penrose (DP) rate}. For point-like nuclei $E_\Delta$ diverges, so the
mass density must be smeared over a length $R_0$. The same noise makes charged matter radiate. An underground germanium detector sees no
such radiation, which excludes smearing at nuclear size and requires $R_0$ of order $10^{-10}$\,m or more~\cite{donadi,figurato}. In the
Newtonian limit of CQ gravity the DP kernel is one admissible choice of the decoherence term~\cite{osswd}.
\end{primer}

\begin{primer}{Classical-quantum dynamics and collapse}{box:cq}
\setlength{\parskip}{3pt plus 1pt}%
\emph{The classical-quantum (CQ) master equation.} A consistent coupling keeps the joint description linear, positive and
normalised, as probability theory demands. The joint state is a family $\varrho(z)$ of positive operators, one for each value
$z$ of the classical variables (here the frame), normalised as $\int\dd z\,\Tr\varrho(z)=1$~\cite{oppenheim}. For one classical variable
$z$ driven by one quantum observable $L$ (for gravity, the mass density), continuous Markovian dynamics has the form~\cite{osswd}
\[
\partial_t\varrho=-\frac i\hbar[H,\varrho]+D_0\Big(L\varrho L-\tfrac12\{L^2,\varrho\}\Big)-\frac{D_1}2\,\partial_z\{L,\varrho\}
+\frac{D_2}2\,\partial_z^2\varrho ,
\]
with three coefficients beyond quantum mechanics:
\begin{itemize}[nosep,leftmargin=*]
\item $D_0$, \emph{decoherence}: the quantum system loses coherence between states that source the classical field differently (for
gravity, different mass distributions);
\item $D_1$, \emph{back-reaction}: a drift of $z$ driven by the quantum state, that is, how matter pushes the frame;
\item $D_2$, \emph{diffusion}: random kicks of $z$, that is, noise in the frame.
\end{itemize}
Complete positivity requires $D_1^2\le4D_0D_2$ here, and $D_0>0$ wherever $D_1\ne0$; Oppenheim et al.\ give the general matrix form of
this trade-off~\cite[Eq.~(23)]{osswd}, which is Eq.~\eqref{eq:tradeoff} of this paper. In words: \emph{a push implies noise}. If quantum matter pushes a
classical field at all, the matter decoheres and the field diffuses, and the product of the two is bounded below by the square of the push,
$D_0D_2\ge D_1^2/4$.

\emph{Measurement and feedback.} An older family of consistent models makes gravity a channel. Each mass's position, or its mass density,
is weakly and continuously measured, and a classical force computed from the record acts back on the other masses (Kafri, Taylor and
Milburn, KTM~\cite{ktm}; Tilloy and Di\'osi~\cite{tilloydiosi}). Newton's law then holds on average, the measurement decoheres
superpositions, and the feedback adds noise. The masses talk only through a classical record, so the channel is LOCC
(Primer~\ref{box:qi}) and cannot entangle them. A weaker measurement decoheres less but gives a noisier record and so more heating; KTM find
both effects at least of the order of the gravitational coupling. The least decoherence of two masses over all measurement strengths is
what this paper calls the \emph{KTM minimum}, $\omega/2$, where $\omega$ is the rate at which a coherent field would build up the
entangling phase.

\emph{Gravity-related collapse.} Di\'osi and Penrose proposed that a superposition of two mass distributions $\rho_1$, $\rho_2$ decays at
the rate $\lambda=E_\Delta/\hbar$, where
\[
E_\Delta=\frac G2\int\!\!\int\dd^3x\,\dd^3y\,\frac{[\rho_1(\mathbf x)-\rho_2(\mathbf x)][\rho_1(\mathbf y)-\rho_2(\mathbf y)]}
{|\mathbf x-\mathbf y|}
\]
is the gravitational self-energy of their difference~\cite{diosi87,penrose96,bassi2013}. For a body of mass $M$ and size $R$ displaced by
more than its size, $\lambda\approx GM^2/\hbar R$: the \emph{Di\'osi--Penrose (DP) rate}. For point-like nuclei $E_\Delta$ diverges, so the
mass density must be smeared over a length $R_0$. The same noise makes charged matter radiate. An underground germanium detector sees no
such radiation, which excludes smearing at nuclear size and requires $R_0$ of order $10^{-10}$\,m or more~\cite{donadi,figurato}. In the
Newtonian limit of CQ gravity the DP kernel is one admissible choice of the decoherence term~\cite{osswd}.
\end{primer}

\subsection{Consistent classical-quantum dynamics}

The consistent alternative evolves a joint state $\varrho(z)$ of the classical frame $z$ and the quantum content by a map that is linear,
completely positive and normalised. Oppenheim showed that such a dynamics exists and can couple a classical metric to quantum
fields~\cite{oppenheim}; Primer~\ref{box:cq} gives its master equation with the decoherence $D_0$, the back-reaction $D_1$ and the diffusion
$D_2$. With matrices of coefficients, complete positivity forces the trade-off~\cite[Eq.~(23)]{osswd}
\begin{equation}\label{eq:tradeoff}
D_1\,D_0^{-1}\,D_1^{\dagger}\;\preceq\;2D_2 ,
\end{equation}
with $D_0^{-1}$ the generalised inverse, and in addition $D_0$ cannot vanish where $D_1$ does not \tagP. Long coherence times of massive superpositions demand strong random fluctuations of the
gravitational field, and weak fluctuations demand short coherence times. For a qubit coupled through $\sigma_z$ to a classical momentum
$p$, with the generator
\begin{equation}\label{eq:cqqubit}
\partial_t\varrho=D_0\,(\sigma_z\varrho\sigma_z-\varrho)-\frac{D_1}{2}\,\partial_p\{\sigma_z,\varrho\}+\frac{D_2}{2}\,\partial_p^2\varrho,
\end{equation}
the bound reads $D_1^2\le4D_0D_2$. Oppenheim et al.\ write the moments of the classical jump with factors $1/n!$, so their $D_1$ and $D_2$
are half of ours, and in their normalisation the same bound is~\eqref{eq:tradeoff}. At equality the dynamics is exactly a measurement of
$\sigma_z$ with feedback, in which $p$ is a perfect record and every block of the state stays pure (Section~\ref{sec:gravsim},
\chk{gravsim:G3}).

\begin{corollary}[GUID diffuses]\label{co:diffuse}
If the frame (Lemma~\ref{lem:classical}) is sourced by quantum content (Hypothesis~\ref{h:info}) through a dynamics that is linear,
completely positive and Markovian, then mass superpositions decohere at a non-zero rate and the gravitational field fluctuates randomly,
with~\eqref{eq:tradeoff} bounding the product from below.
\end{corollary}

The corollary turns the classicality of the frame from a philosophical position into a quantitative one: in the Newtonian limit, the
decoherence rates of massive superpositions measured in interferometry and the noise in precision measurements of gravitational forces
together squeeze the remaining parameter space from both sides~\cite{osswd}.

\subsection{The record is a measurement, the layout is feedback}

DMP's architecture coincides with the second, older family of consistent models, measurement and feedback (Primer~\ref{box:cq}): Kafri,
Taylor and Milburn measure each mass's position and apply the recorded force to the other~\cite{ktm}, and Tilloy and Di\'osi source a
Newtonian potential from a continuously measured mass density~\cite{tilloydiosi}. Both decohere and add noise, as
Corollary~\ref{co:diffuse} requires, and for two masses the least decoherence is the KTM minimum $\omega/2$, reached where the heating
is also least (Section~\ref{sec:gravsim}, \chk{gravsim:G4}).

Read through the dictionary, DMP's Axiom~3 has exactly this structure \tagD+\tagH. Every act leaves a record that every witness reads in
full: a measurement of the act's existence. Any other reader sees the counts: the measured mass density, coarse-grained. The layout of the
system is ``the median of that demand, recomputed by both halves of the layer and never by one'': a classical function of the records,
applied back to the system, which is feedback. From the parts' point of view the whole loop is a measure-and-prepare channel, which is why
it cannot entangle (Lemma~\ref{lem:locc}).

\begin{proposition}[DMP is a measurement-and-feedback gravity]\label{pr:ktm}
Under Hypotheses~\ref{h:exist}--\ref{h:info}, the dynamics that DMP's axioms prescribe for the trust field is a continuous measurement of
existence (energy-momentum) followed by classical feedback on the geometry, of the Kafri--Taylor--Milburn and Tilloy--Di\'osi type.
\end{proposition}

The proposition fixes the architecture, not the constants. A measurement-and-feedback model has a free parameter, the measurement strength
or equivalently the smearing length of the measured mass density, which sets the decoherence rate and the noise. DMP derives no constant
(Section~\ref{sec:break}), and the most economical choice, the parameter-free Di\'osi--Penrose collapse~\cite{diosi87,penrose96}, is
excluded by the absence of the spontaneous radiation it predicts in an underground germanium detector~\cite{donadi} \tagP. GUID must
therefore carry a free strength, fixed by experiment or by a principle that ties DMP to physical units. \emph{Verdict:} an identification \tagH{} of
the architecture, not a derivation of the dynamics; Section~\ref{sec:bridges} computes the constant from the two principles GUID
suggests.

\subsection{Fixing the free constant}\label{sec:bridges}

Proposition~\ref{pr:ktm} fixes the architecture, not the constant. DMP cannot supply the constant by itself: its text makes every
parameter configuration (T4, T32), and it has no unit of length, time or action. A number therefore needs one bridge principle that ties a
DMP structure to a physical scale. We compute the two that GUID's own dictionary suggests. Neither is derived from the axioms; both are
hypotheses \tagH, stated so that they can fail.

\paragraph{The balance rule.} Both bridges use one structural rule. Axiom~1's balance moves every part towards its declared state, and the
system's own acts never raise the VIBE potential (T32); read as dynamics, the frame adds no more noise than consistency requires. GUID
therefore holds the trade-off~\eqref{eq:tradeoff} at equality,
\begin{equation}\label{eq:saturate}
2D_2=D_1\,D_0^{-1}\,D_1^{\dagger},
\end{equation}
so that a measured gravitational decoherence would fix the spacetime diffusion, and conversely. Other classical-quantum theories leave the
two free above the bound; \eqref{eq:saturate} is GUID's own rule.

\paragraph{Bridge B: one record per Planck time.} The second bridge reads DMP's tick as the Planck time: the frame records existence at
Planck strength, so the coupling takes its value in Planck units, and the balance rule~\eqref{eq:saturate} holds. In the Newtonian limit
Oppenheim et al.\ bound the coupling from both sides, from above by precision measurements of gravitational forces and from below by the
observed coherence of massive molecules~\cite{osswd}. At Planck strength the discrete version of the dynamics has
$(l_P^3/m_P)D_2=m_Pt_P\approx10^{-51}$\,kg\,s, below the allowed window $10^{-25}$ to $10^{-1}$\,kg\,s: \emph{excluded}. The continuous
version with the Di\'osi--Penrose kernel has $l_P^2D_2=m_P^2/c=1.6\times10^{-24}$\,kg$^2$\,s/m, inside the allowed window $10^{-40}$ to
$10^{-9}$\,kg$^2$\,s/m: \emph{allowed} \chk{Q19}. With the trade-off at equality, a mass $M$ of size $R$ in a superposition of two
well-separated positions decoheres at the rate
\begin{equation}\label{eq:dprate}
\lambda=\frac{M^2}{R\,l_P^2D_2}=\frac{GM^2}{\hbar R},
\end{equation}
since $c/m_P^2=G/\hbar$: the Di\'osi--Penrose rate~\cite{diosi87,penrose96}, up to an order-one factor (for a uniform sphere the
exact self-energy gives $6/5-R/\Delta x$ in place of 1). Here $D_2$ is the strength of the diffusion kernel of the Newtonian potential in
Oppenheim et al.'s Newtonian-limit model, in their normalisation, and the two windows are theirs~\cite[Eqs.~(44), (46), (47)]{osswd}. For a fullerene it is $6\times10^{-16}$ per second, far below what interferometry allows. For the diamond microsphere of the Bose
et al.\ proposal ($10^{-14}$\,kg, radius 0.88\,\textmu m) it is 72 per second: coherence is lost in about 14\,ms, while the entangling phase
needs 8 to 25\,s, so only $2\times10^{-3}$\,rad builds up \chk{Q19}. Of the two versions of the bridge, the discrete one is excluded by
the windows and the continuous one allowed; adopting the continuous one is a choice made with the windows in view. One length remains that
the bridge does not fix: the mass density of each nucleus must be smeared over a length $R_0$, or the noise makes matter radiate more than
an underground germanium detector sees; we adopt the bound of current analyses, $R_0\ge4\times10^{-10}$\,m~\cite{donadi,figurato}, at
which the bulk term in~\eqref{eq:dprate} still dominates \chk{Q19}. \emph{Verdict: live} and sharp: with Bridge~B, mesoscopic
superpositions decohere at the Di\'osi--Penrose rate. Section~\ref{sec:data} shows how far today's superpositions are from testing
it.

\paragraph{What the bridges show.} Fixed by its own bridge hypotheses, GUID lands on known physics: Bridge~A (Section~\ref{sec:darkmatter})
is refuted, and Bridge~B is the Di\'osi--Penrose model, which predates this work, with its smearing length set by experiment. What GUID
adds is a reading of that model, the balance rule~\eqref{eq:saturate} with a record at Planck strength, not a derivation: at Planck
strength the only Newtonian decoherence rate is $G/\hbar$ times a self-energy, so landing on $GM^2/\hbar R$ is close to dimensional
analysis. A number of its own would need a prediction in which GUID differs from
both the classical-quantum theory and Di\'osi--Penrose. The most promising lead is DMP's two independent witnesses, which may imply two
independent noise sources with a statistical signature that a single source lacks; we have not worked it out.

\section{The experiments that decide}\label{sec:exp}

\subsection{Gravitationally mediated entanglement}

Bose et al.~\cite{bose} and Marletto and Vedral~\cite{mv} proposed to place two masses, each in a superposition of two positions, close
enough that only gravity couples them, and to test whether they end entangled. Let each mass $m$ be in $(\ket L+\ket R)/\sqrt2$ with the two
positions $\Delta x$ apart and the centres $d$ apart on one line. A coherent Newtonian field gives each branch $(i,j)$ the phase
$\varphi_{ij}=Gm^2t/(\hbar d_{ij})$, where $d_{ij}$ is the distance between the masses in that branch, so after time $t$
\begin{equation}\label{eq:bmv}
\ket{\Psi(t)}=\tfrac12\sum_{i,j\in\{L,R\}}e^{i\varphi_{ij}(t)}\ket i_A\ket j_B .
\end{equation}
Up to local phases this is $\tfrac12(\ket{00}+\ket{01}+\ket{10}+e^{i\Delta\varphi}\ket{11})$ with
$\Delta\varphi=\varphi_{RL}+\varphi_{LR}-\varphi_{LL}-\varphi_{RR}$; its concurrence is $|\sin(\Delta\varphi/2)|$ and its negativity
\begin{equation}\label{eq:neg}
\mathcal N=\tfrac12\,|\sin(\Delta\varphi/2)|.
\end{equation}
For the parameters of~\cite{bose}, $m=10^{-14}$\,kg, $\Delta x=250$\,\textmu m, $d=450$\,\textmu m, the entangling phase grows at
$0.126$\,rad/s: one radian after $8.0$\,s, maximal entanglement ($\mathcal N=1/2$, $\Delta\varphi=\pi$) after $25$\,s \chk{Q5}. A recorded,
classical field (Lemma~\ref{lem:classical}) in its projective extreme leaves each branch's phase but removes the coherence between
branches; as a measurement-and-feedback channel it reproduces the branch phases on average and decoheres each mass at a rate of at least
$\omega/2$, the classical-channel minimum. Either way the masses stay separable at all times, $\mathcal N=0$ \chk{Q5}, \chk{gravsim:G4}.

\textbf{GUID predicts $\mathcal N=0$.} With Bridge~B of Section~\ref{sec:bridges} it predicts more: the masses decohere at the
Di\'osi--Penrose rate, about 72 per second for the proposed microspheres, long before an entangling phase can build up.

\emph{What a result would show.} If decoherence at the Di\'osi--Penrose rate is real, a quantised field would not entangle two solid
spheres either: the ratio of the entangling to the decoherence rate is fixed by the geometry and stays below what entanglement needs
(Section~\ref{sec:bmv}). The first experiment to decide is therefore whether that decoherence exists, which tests Bridge~B; only in its
absence does the entanglement experiment separate classical from quantum gravity. A positive result would exclude every classical
mediator. It would refute Hypothesis~1$'$, which puts the field that pulls laboratory masses into the frame; it would not refute DMP,
which hosts either reading of that field (Section~\ref{sec:dissent}): a quantum geometry would then run on DMP as content
(Proposition~\ref{th:host}(ii)). Section~\ref{sec:bmv} also checks the spin witness proposed to certify the entanglement. Two
caveats are under debate. Aziz and Howl argue that classical gravity coupled to quantum matter fields can produce entanglement through virtual matter
processes, at a rate far below the quantum-gravity rate~\cite{azizhowl}; Marletto, Oppenheim, Vedral and Wilson~\cite{movw},
Di\'osi~\cite{diosi25}, Schneider, Huggett and Linnemann~\cite{shl} and Gundhi, Infantino and Bassi~\cite{gib} dispute it. And
Christodoulou et al.\ showed that the entanglement follows from linearised quantum gravity on a fixed background with local
mediation~\cite{christodoulou23}, which is the content-side mechanism of the dissent in Section~\ref{sec:dissent}. Neither changes the prediction:
in GUID the frame is recorded, and Lemma~\ref{lem:locc} applies to it.

\subsection{Quantumness without entanglement}

Lami, Pedernales and Plenio showed that the quantum nature of gravity can be tested without creating entanglement: an LOCC inequality bounds
how well any local classical channel can reproduce a gravitational evolution of coherent states of massive oscillators, such as torsion
pendula, and a measured violation would exclude a classical mediator~\cite{lpp}. The test needs no large spatial superpositions. GUID
predicts that the inequality holds.

\subsection{Diffusion against decoherence}

Corollary~\ref{co:diffuse} gives a two-sided test. Interferometry with ever heavier molecules and levitated particles bounds the decoherence
that gravity may cause from above; precision measurements of the gravitational force, which has been measured between masses as small as
90\,mg~\cite{westphal}, bound the random fluctuations of the field from above. Inequality~\eqref{eq:tradeoff} says the two cannot both be
made arbitrarily small. GUID predicts that they cannot.

\subsection{What existing experiments already show}

Quantum matter responds to gravity as the framework requires: neutrons in a gravitational potential acquire the phase $mgA/\hbar v$
predicted for quantum matter in a classical field (Colella, Overhauser and Werner~\cite{cow}), and ultracold neutrons above a mirror occupy
discrete gravitational quantum states~\cite{nesvizhevsky}. In both, the matter is quantum and the field is a given potential; neither tests
whether the field itself can be in superposition. The current state of the field is reviewed by Carney, Stamp and Taylor~\cite{cst}.
\emph{Verdict on GUID's classical frame: live, falsifiable.}

\subsection{A dissenting reading}\label{sec:dissent}

The same facts can be scored the other way. Split the geometry, as the
background-field method does: the background is a DMP structure and classical; the linearised fluctuation, which carries the Newtonian
pull between laboratory masses, is content between parts, and linearised quantum gravity on a fixed background already predicts the
entanglement of the masses~\cite{christodoulou23}. On that reading DMP's classical structures are the part of quantum gravity that is
classical in quantum gravity too, the classical centre of a holographic code (Section~\ref{sec:holo}), and the linearised field is quantum content on that classical
centre. The two readings agree on every fact: the shared structures are classical, quantum content including a linearised field can
move between parts by grant and stay entangled while every act is recorded, the axioms derive no constant, and they forbid coherent
superpositions of backgrounds while their records exist. They differ in one place: whether the linearised field that pulls laboratory
masses belongs to the frame (the proposed theory, Hypothesis~1$'$, which predicts $\mathcal N=0$) or to the content (the dissent,
which predicts $\mathcal N>0$ as long as the decoherence stays below $\omega/2$). DMP hosts either. If Di\'osi--Penrose decoherence is
real, neither reading gives entanglement for two spheres, and the decoherence measurement comes first; only in its absence does the
entanglement experiment decide between them.

\section{Black holes}\label{sec:bh}

\subsection{The horizon as containment}

Hypothesis~\ref{h:horizon} maps DMP's containment onto the event horizon. The match is close. No act carries information across a containing
boundary (T31); no causal curve leaves the horizon \tagP. The counts every outside reader sees (Axiom~3) are the hair a stationary black hole
keeps: mass $M$, angular momentum $J$, charge $Q$ (Israel~\cite{israel}; Kerr~\cite{kerr}). Inside a Schwarzschild horizon the radial
coordinate is timelike and every future-directed path ends at the singularity: everything inside derives its causal future from it, a
local ground in the sense of G0. Whether crossing is survivable is answered by the equivalence principle: at the horizon of a
supermassive hole the tidal stretch across a body is small; for a stellar-mass hole it is lethal before the horizon.

\subsection{What a black hole gives up: Hawking--Penrose--Information radiation}

The first law of black-hole mechanics (Bardeen, Carter and Hawking~\cite{bch}) reads, with $G=c=1$,
\begin{equation}\label{eq:firstlaw}
\dd M=\frac{\kappa}{8\pi}\,\dd A+\Omega_H\,\dd J+\Phi_H\,\dd Q ,
\end{equation}
with $\Omega_H$ the angular velocity and $\Phi_H$ the electric potential of the horizon. It separates the two ways a black hole loses mass.
The rotational and electric terms can be extracted from outside at fixed or growing area: Penrose's process with particles~\cite{penrose69},
superradiant scattering with waves, and the Blandford--Znajek process with magnetic fields~\cite{bz}, the leading model of the jets of
active galaxies. All of them work in the ergosphere, which reaches $r=2GM/c^2$ at the equator and lies outside the horizon
$r_+=(GM/c^2)(1+\sqrt{1-a^2})$ for every non-zero spin \chk{Q10}. The area term is released only quantum mechanically, as Hawking's thermal
radiation, which lowers $A$.

\begin{definition}[Hawking--Penrose--Information radiation]
The \emph{Hawking--Penrose--Information (HPI) radiation} of a black hole is everything it gives up: the extraction of its rotational and
electric energy from the ergoregion (the $\Omega_H\dd J$ and $\Phi_H\dd Q$ terms of~\eqref{eq:firstlaw}) and its thermal Hawking radiation (the
$\kappa\,\dd A/8\pi$ term).
\end{definition}

In GUID terms the extraction is the frame being dragged along by the spinning hole, outside the horizon, so no act crosses it; the
Hawking part is the floor of T31. What leaves carries the counts and no content, as Axiom~3 requires. Spinning holes are more efficient
engines: the binding energy radiated by matter falling to the innermost stable orbit is $1-\sqrt{8/9}=5.7\%$ of its rest energy without
spin, $32\%$ at Thorne's limit $a=0.998$, and approaches $1-1/\sqrt3=42.3\%$ at maximal spin (Bardeen, Press and Teukolsky~\cite{bpt};
\chk{Q10}). At fixed mass every additional count (spin, charge) shrinks the area, so the more an outside reader knows of a boundary, the
less entropy it holds and the more of its energy can be extracted \tagP. The Hawking part is unobservable for astrophysical holes: a
solar-mass hole is at $6.2\times10^{-8}$\,K, colder than the $2.7$\,K microwave sky, so it absorbs more than it emits; only holes lighter
than about $0.6$ lunar masses are hotter than the sky, and a solar-mass hole in empty space would take more than $10^{66}$ years to
evaporate \chk{Q11}. \emph{Verdict: holds}; the Hawking part is \emph{untested in nature}.

\subsection{The information problem}

If evaporation is unitary, the entropy of the radiation first rises and then falls after the Page time~\cite{page93}; the island formula,
derived from the gravitational path integral with replica wormholes, reproduces that curve~\cite{penington,aemm}. Those derivations sum
over geometries, which is quantum gravity in the sense GUID's frame excludes. On the other side, a classical-quantum dynamics decoheres the
quantum system (Corollary~\ref{co:diffuse}), so it is not unitary on the quantum system alone and does not by itself produce a Page curve.
DMP's T31 as written, ``what still leaves is entropy and nothing else'', and its T27, which destroys contained items at their scheduled date,
side with the thermal picture. In GUID the question is therefore open in a specific form: whether the content held inside a horizon (which
may be quantum, Proposition~\ref{th:host}(ii)) returns in the radiation as content, while the frame stays classical. No observation reaches it.
\emph{Verdict: open}, in physics and in GUID.

\section{Holography and quantum error correction}\label{sec:holo}
% GUID v6, Primer P5 (label box:holo): holography and quantum error correction. Background only; it states no result of the paper.
% Place it at the start of Section sec:holo (marker %PRIMER:holo in part1.tex): % GUID v6, Primer P5 (label box:holo): holography and quantum error correction. Background only; it states no result of the paper.
% Place it at the start of Section sec:holo (marker %PRIMER:holo in part1.tex): % GUID v6, Primer P5 (label box:holo): holography and quantum error correction. Background only; it states no result of the paper.
% Place it at the start of Section sec:holo (marker %PRIMER:holo in part1.tex): \input{primers/p5-holo}
\begin{primer}{Holography and quantum error correction}{box:holo}
\emph{Stabiliser codes.} A quantum error-correcting code stores $k$ logical qubits in $n$ physical ones so that some errors can be
undone. In a \emph{stabiliser code} the code space is the joint $+1$ eigenspace of $n-k$ commuting products of Pauli matrices, the
\emph{stabilisers}~\cite{gottesman1997}. Measuring them gives the \emph{syndrome}. An error that anticommutes with a stabiliser flips its
sign and so shows in the syndrome. The stabilisers commute with the logical operators, so the syndrome reveals nothing about the stored
content. A code of distance $d$, written $[[n,k,d]]$, corrects any $\lfloor(d-1)/2\rfloor$ unknown errors, or any $d-1$
\emph{erasures} (losses at known places).

\emph{The five-qubit code.} The smallest code that corrects every single-qubit error is $[[5,1,3]]$, with the stabiliser $XZZXI$ and its
cyclic shifts~\cite{laflamme,bennett1996}. It corrects any two erasures, so its logical qubit can be recovered from any three of the five
qubits and from no two: a $((3,5))$ quantum threshold scheme~\cite{cgl}.

\emph{Perfect tensors and the HaPPY code.} A tensor with $2m$ legs is \emph{perfect} if every split of its legs into $m$ and $m$ is, up to
normalisation, a unitary map. The encoding map of the five-qubit code, with its logical leg as a sixth, is such a tensor. Pastawski,
Yoshida, Harlow and Preskill tile the hyperbolic disc with pentagons, place one perfect tensor on each, contract neighbouring legs, and
leave one leg per tile in the bulk and the outermost legs on the boundary~\cite{happy}. The network is an isometry from the bulk to the
boundary, a toy model of the bulk-boundary correspondence of AdS/CFT~\cite{maldacena}.

\emph{Ryu--Takayanagi as a minimal cut.} In AdS/CFT the entropy of a boundary region $A$ is $S(A)=\mathrm{Area}(\gamma_A)/4G\hbar$ (with
$c=k_B=1$), where $\gamma_A$ is the bulk surface of least area anchored on the edge of $A$~\cite{rt}. In a network of perfect tensors the
entropy of $A$, in bits, is for many regions the number of qubit legs cut by the smallest cut that separates $A$ from the rest of the
boundary~\cite{happy}.

\emph{Reconstruction.} The bulk region between $A$ and $\gamma_A$ is the \emph{entanglement wedge} of $A$. Every bulk operator inside it can
be represented by an operator acting on $A$ alone~\cite{adh,dong2016}: the bulk is encoded redundantly, like a logical qubit. A code has
\emph{complementary recovery} when what cannot be reconstructed on $A$ can be reconstructed on the rest of the boundary.

\emph{The area operator.} Harlow showed that every code with complementary recovery obeys a quantum Ryu--Takayanagi formula
(Eq.~\eqref{eq:harlow} in Section~\ref{sec:holo}). Its area term is the expectation value of an \emph{area operator} that lies in the centre
of the algebra reconstructable on $A$: it commutes with everything that either side can do~\cite{harlow}. A \emph{fixed-area state} is an
eigenstate of this operator, and such states have a flat entanglement spectrum~\cite{dhm}. A general state of quantum gravity is a
superposition of areas.
\end{primer}

\begin{primer}{Holography and quantum error correction}{box:holo}
\emph{Stabiliser codes.} A quantum error-correcting code stores $k$ logical qubits in $n$ physical ones so that some errors can be
undone. In a \emph{stabiliser code} the code space is the joint $+1$ eigenspace of $n-k$ commuting products of Pauli matrices, the
\emph{stabilisers}~\cite{gottesman1997}. Measuring them gives the \emph{syndrome}. An error that anticommutes with a stabiliser flips its
sign and so shows in the syndrome. The stabilisers commute with the logical operators, so the syndrome reveals nothing about the stored
content. A code of distance $d$, written $[[n,k,d]]$, corrects any $\lfloor(d-1)/2\rfloor$ unknown errors, or any $d-1$
\emph{erasures} (losses at known places).

\emph{The five-qubit code.} The smallest code that corrects every single-qubit error is $[[5,1,3]]$, with the stabiliser $XZZXI$ and its
cyclic shifts~\cite{laflamme,bennett1996}. It corrects any two erasures, so its logical qubit can be recovered from any three of the five
qubits and from no two: a $((3,5))$ quantum threshold scheme~\cite{cgl}.

\emph{Perfect tensors and the HaPPY code.} A tensor with $2m$ legs is \emph{perfect} if every split of its legs into $m$ and $m$ is, up to
normalisation, a unitary map. The encoding map of the five-qubit code, with its logical leg as a sixth, is such a tensor. Pastawski,
Yoshida, Harlow and Preskill tile the hyperbolic disc with pentagons, place one perfect tensor on each, contract neighbouring legs, and
leave one leg per tile in the bulk and the outermost legs on the boundary~\cite{happy}. The network is an isometry from the bulk to the
boundary, a toy model of the bulk-boundary correspondence of AdS/CFT~\cite{maldacena}.

\emph{Ryu--Takayanagi as a minimal cut.} In AdS/CFT the entropy of a boundary region $A$ is $S(A)=\mathrm{Area}(\gamma_A)/4G\hbar$ (with
$c=k_B=1$), where $\gamma_A$ is the bulk surface of least area anchored on the edge of $A$~\cite{rt}. In a network of perfect tensors the
entropy of $A$, in bits, is for many regions the number of qubit legs cut by the smallest cut that separates $A$ from the rest of the
boundary~\cite{happy}.

\emph{Reconstruction.} The bulk region between $A$ and $\gamma_A$ is the \emph{entanglement wedge} of $A$. Every bulk operator inside it can
be represented by an operator acting on $A$ alone~\cite{adh,dong2016}: the bulk is encoded redundantly, like a logical qubit. A code has
\emph{complementary recovery} when what cannot be reconstructed on $A$ can be reconstructed on the rest of the boundary.

\emph{The area operator.} Harlow showed that every code with complementary recovery obeys a quantum Ryu--Takayanagi formula
(Eq.~\eqref{eq:harlow} in Section~\ref{sec:holo}). Its area term is the expectation value of an \emph{area operator} that lies in the centre
of the algebra reconstructable on $A$: it commutes with everything that either side can do~\cite{harlow}. A \emph{fixed-area state} is an
eigenstate of this operator, and such states have a flat entanglement spectrum~\cite{dhm}. A general state of quantum gravity is a
superposition of areas.
\end{primer}

\begin{primer}{Holography and quantum error correction}{box:holo}
\emph{Stabiliser codes.} A quantum error-correcting code stores $k$ logical qubits in $n$ physical ones so that some errors can be
undone. In a \emph{stabiliser code} the code space is the joint $+1$ eigenspace of $n-k$ commuting products of Pauli matrices, the
\emph{stabilisers}~\cite{gottesman1997}. Measuring them gives the \emph{syndrome}. An error that anticommutes with a stabiliser flips its
sign and so shows in the syndrome. The stabilisers commute with the logical operators, so the syndrome reveals nothing about the stored
content. A code of distance $d$, written $[[n,k,d]]$, corrects any $\lfloor(d-1)/2\rfloor$ unknown errors, or any $d-1$
\emph{erasures} (losses at known places).

\emph{The five-qubit code.} The smallest code that corrects every single-qubit error is $[[5,1,3]]$, with the stabiliser $XZZXI$ and its
cyclic shifts~\cite{laflamme,bennett1996}. It corrects any two erasures, so its logical qubit can be recovered from any three of the five
qubits and from no two: a $((3,5))$ quantum threshold scheme~\cite{cgl}.

\emph{Perfect tensors and the HaPPY code.} A tensor with $2m$ legs is \emph{perfect} if every split of its legs into $m$ and $m$ is, up to
normalisation, a unitary map. The encoding map of the five-qubit code, with its logical leg as a sixth, is such a tensor. Pastawski,
Yoshida, Harlow and Preskill tile the hyperbolic disc with pentagons, place one perfect tensor on each, contract neighbouring legs, and
leave one leg per tile in the bulk and the outermost legs on the boundary~\cite{happy}. The network is an isometry from the bulk to the
boundary, a toy model of the bulk-boundary correspondence of AdS/CFT~\cite{maldacena}.

\emph{Ryu--Takayanagi as a minimal cut.} In AdS/CFT the entropy of a boundary region $A$ is $S(A)=\mathrm{Area}(\gamma_A)/4G\hbar$ (with
$c=k_B=1$), where $\gamma_A$ is the bulk surface of least area anchored on the edge of $A$~\cite{rt}. In a network of perfect tensors the
entropy of $A$, in bits, is for many regions the number of qubit legs cut by the smallest cut that separates $A$ from the rest of the
boundary~\cite{happy}.

\emph{Reconstruction.} The bulk region between $A$ and $\gamma_A$ is the \emph{entanglement wedge} of $A$. Every bulk operator inside it can
be represented by an operator acting on $A$ alone~\cite{adh,dong2016}: the bulk is encoded redundantly, like a logical qubit. A code has
\emph{complementary recovery} when what cannot be reconstructed on $A$ can be reconstructed on the rest of the boundary.

\emph{The area operator.} Harlow showed that every code with complementary recovery obeys a quantum Ryu--Takayanagi formula
(Eq.~\eqref{eq:harlow} in Section~\ref{sec:holo}). Its area term is the expectation value of an \emph{area operator} that lies in the centre
of the algebra reconstructable on $A$: it commutes with everything that either side can do~\cite{harlow}. A \emph{fixed-area state} is an
eigenstate of this operator, and such states have a flat entanglement spectrum~\cite{dhm}. A general state of quantum gravity is a
superposition of areas.
\end{primer}


\subsection{The strongest assembly}

The axioms can be assembled as a holographic quantum error-correcting code on a fixed background, the kinematics of the HaPPY
model~\cite{happy}. A part (``no further statements about its constituents'') becomes one block of the five-qubit code~\cite{laflamme},
which encodes one logical qubit in five physical ones with distance 3. Its encoding map is a \emph{perfect tensor} with six legs: every
bipartition of the six legs into three and three is a unitary map, so every three-leg cut carries three qubits of maximal entanglement
\chk{A1}. A layer, which Axiom~3 makes a tree rooted at its Duality, becomes a rooted tree of such tiles, one root and five children with
twenty boundary legs. In the orientation we use, a child tile's logical leg carries the edge to its parent and one of its physical legs
carries the information sitting there (the bulk leg); both tile maps are isometries \chk{A2}.

\paragraph{Minimal cuts.} For a contiguous boundary region $R$, the entropy of the reduced boundary state equals the number of legs cut
by the minimal cut through the tree that separates $R$ from its complement \chk{A2}. This is the Ryu--Takayanagi formula~\cite{rt},
$S(R)=|\gamma_R|/4G\hbar$ with $\gamma_R$ the minimal bulk surface anchored on $\partial R$, on a tree. DMP's ``fewest hops in total''
selects the same minimal surfaces. The equality holds for every region, contiguous or not, when each bulk leg is left open (entangled with
a reference); with the bulk legs fixed it fails for some non-contiguous regions, where a tile whose legs the cut splits two against two
lets one classical bit past, and in the opposite orientation of the tiles it fails even for intervals (Section~\ref{sec:holocode},
\chk{holo:H2}).

\paragraph{Reconstruction.} The five-qubit code corrects any two erasures, so the logical qubit is recoverable from any three of the five
physical qubits and from no two: a $((3,5))$ quantum threshold scheme~\cite{cgl}. In mutual information with a reference $R$ purifying the
root, $I(R:\text{pair})=0$ for all ten pairs and $I(R:\text{triple})=2$ bits for all ten triples \chk{A3}. This is entanglement-wedge
reconstruction~\cite{adh}: a bulk operator is reconstructable on any boundary region whose wedge contains it.

\paragraph{The two witnesses.} The trust half and the information half each measure the stabiliser syndrome of the code. Stabilisers commute
with the logical operators, so the syndrome says that content is present and in the code, never what it is: this is Axiom~3's existence
clause, exactly. Both halves read the same syndrome; the gate passes only on two trivial syndromes; each of the 15 single-qubit faults gives
a distinct syndrome on which both refuse, and correction restores the content with fidelity 1. That every single-qubit fault is caught is a
property of this code's stabilisers, not of existence checks in general: a witness that reads only the existence sector detects faults
that leave the sector and passes, by construction, a fault that changes the content inside it. The witnesses certify existence; they do
not protect content (Section~\ref{sec:hosted} shows both behaviours on sectors of up to 2.7 million states). If the two halves instead read the content
itself, the content is disturbed (fidelity exactly $1/2$ for every content, \chk{holo:H5}), so ``full reach over the act'' must mean full reach over its
existence \chk{A4}.

\paragraph{Two layers in one object.} Two parts whose content is a logical Bell pair reach the Tsirelson value $2\sqrt2\approx2.83$ of the
CHSH combination on their content, while the gates' verdicts, deterministic syndrome readings, give exactly the classical value 2
\chk{A5}. The gate is classical and the content quantum, at the same time and in the same assembly.

\subsection{Where the assembly stops: the classical centre}

In any code with complementary recovery, Harlow proved the quantum Ryu--Takayanagi formula
\begin{equation}\label{eq:harlow}
S(\rho_R)=\Tr(\rho\,\mathcal L_R)+S_{\text{bulk}}(\rho,\mathcal M_R),
\end{equation}
where $\mathcal M_R$ is the algebra of bulk operators reconstructable on $R$ and the \emph{area operator} $\mathcal L_R$ lies in the centre of
$\mathcal M_R$: it commutes with everything either side can do~\cite{harlow}. DMP's structures recomputed by both halves are such central
elements. Lemma~\ref{lem:classical} then says that in GUID the area is a number, never an operator with a spread of values: the axioms
allow one fixed-area sector at a time. In holography fixed-area states are a special, though useful, class (Dong, Harlow and
Marolf~\cite{dhm}); the general state of quantum gravity superposes areas. So the assembly reproduces the kinematics of holographic codes on
a fixed background (\chk{A1}--\chk{A5}) and stops exactly at superpositions of geometry. Section~\ref{sec:holocode} runs the assembly at scale.

\paragraph{The tree is not our space.} Read as space, a tree's boundary is totally disconnected (the $p$-adic line of $p$-adic
AdS/CFT~\cite{hmss,gubser}), and no tree approximates a plane: an $n$-cycle needs distortion at least $n/3-1$ in any tree~\cite{rr}. So a
literal tree cannot be our space. Keeping the hierarchy as the routing tree of each request over a graph with loops removes this objection
and touches none of the results above. Read as the tree of histories, each branch one consistent record (decoherent
histories~\cite{gellmannhartle}), the tree owes no spacelike structure at all.

\section{Lower dimensions}\label{sec:21}

Why 3+1? In $D$ spacetime dimensions a massless spin-2 field has $D(D-3)/2$ polarisations and the Weyl tensor, the part of the curvature
not fixed by local matter, has $D^2(D^2-1)/12-D(D+1)/2$ components: two and ten in 3+1, none and none in 2+1 \chk{Q8}. In 2+1 a mass does not
attract another at all; it cuts a wedge out of space (Deser, Jackiw and 't~Hooft~\cite{djt}). Gravity there is, perturbatively, a
Chern--Simons gauge theory~\cite{witten88}, with a holographic boundary algebra for negative cosmological constant~\cite{bh}; whether pure
quantum gravity in 2+1 exists as a consistent theory is open~\cite{mw}. Its degrees of freedom are global: the holonomies of the flat
connection around particles and handles, which do not commute once quantised~\cite{nr}, and point particles scatter through their conical
deficits~\cite{thooft88}.

So 2+1 is not an exception to Proposition~\ref{th:host}. With no local field there is no local mediator to forbid, but the global holonomies are
quantum, and a DMP structure that carries them is broadcast and so can carry one classical sector of them, never their superposition; with
matter, the same rule decoheres the particles' relative positions. In GUID's terms, the frame in 2+1 would have nothing local
that could be held as a classical record, and what it does carry is global and quantum. DMP can run 2+1 quantum gravity as content, but
its frame cannot be 2+1 gravity;
and our universe is 3+1, whose local degrees of freedom, gravitational waves, are measured~\cite{gw150914}. \emph{Verdict: holds as a reading}, with
the dimension count as \emph{holds}; a classical 2+1 frame with classical holonomies would also be consistent.

\section{Cosmology}\label{sec:cosmo}

% GUID v6, Primer P6 (label box:cosmo): cosmology. Background only; it states no result of the paper.
% Place it at the start of Section sec:cosmo (marker %PRIMER:cosmo in part1.tex): % GUID v6, Primer P6 (label box:cosmo): cosmology. Background only; it states no result of the paper.
% Place it at the start of Section sec:cosmo (marker %PRIMER:cosmo in part1.tex): % GUID v6, Primer P6 (label box:cosmo): cosmology. Background only; it states no result of the paper.
% Place it at the start of Section sec:cosmo (marker %PRIMER:cosmo in part1.tex): \input{primers/p6-cosmo}
\begin{primer}{Cosmology}{box:cosmo}
\emph{The expanding universe.} On large scales the universe is homogeneous and isotropic. Distances grow with a scale factor $a(t)$, light
from the past is redshifted by $1+z=a_0/a$, and the expansion rate $H=\dot a/a$ obeys, for a flat universe,
$H^2=8\pi G\rho/3+\Lambda c^2/3$. Each component's share is $\Omega_i=\rho_i/\rho_c$ with the critical density $\rho_c=3H^2/8\pi G$; the
value today is $H_0$.

\emph{$\Lambda$CDM.} The standard model of cosmology adds to the known particles a cold, collisionless \emph{dark matter} and a
cosmological constant $\Lambda$. Six parameters fix it: the densities of baryons and dark matter ($\Omega_bh^2$, $\Omega_ch^2$), the angular
size of the sound horizon, the optical depth to reionisation, and the amplitude $A_s$ and tilt $n_s$ of the primordial
fluctuations~\cite{planck18}.

\emph{Inflation.} An early epoch of nearly exponential expansion, $a\propto e^{Ht}$, flattens and smooths the universe and stretches
quantum fluctuations into the seeds of galaxies~\cite{guth}. It is driven by a scalar field, the \emph{inflaton}, with potential $V$, and
its length is counted in \emph{e-folds}, $N=\ln(a_{\rm end}/a_{\rm start})$. In \emph{slow roll} the field creeps down a flat potential;
with the reduced Planck mass $M_{\rm pl}$, the parameters $\epsilon=\tfrac12M_{\rm pl}^2(V'/V)^2$ and $\eta_V=M_{\rm pl}^2V''/V$ are small. The
fluctuations then have a nearly scale-invariant spectrum with tilt $n_s\approx1-6\epsilon+2\eta_V$, and gravitational waves with a
tensor-to-scalar ratio $r\approx16\epsilon$~\cite{planck2020x}.

\emph{Reheating and oscillons.} After inflation the inflaton oscillates about the minimum of its potential and decays into ordinary
particles, which \emph{reheats} the universe~\cite{kls}. In a broad class of models it first fragments into \emph{oscillons}, long-lived
localised lumps of the field~\cite{amin}.

\emph{The probes.} The cosmic microwave background (CMB) was released at recombination, $z\approx1100$. Its acoustic peaks are the imprint of
sound waves in the early plasma: their spacing fixes the sound horizon and their heights the densities~\cite{planck18}. The same sound
horizon, about 150\,Mpc today, is imprinted on the clustering of galaxies as \emph{baryon acoustic oscillations} (BAO), a standard ruler
for distance against redshift~\cite{desi2}.

\emph{Dark energy.} A cosmological constant has the equation of state $w=p/\rho c^2=-1$ at all times. DESI fits the
Chevallier--Polarski--Linder (CPL) form $w(a)=w_0+w_a(1-a)$~\cite{chevallier2001,linder2003}. A \emph{thawing} dark energy has $w_0>-1$
and $w_a<0$: it was close to $-1$ early and weakens today.

\emph{Data names and tensions.} P-ACT combines the CMB spectra of Planck and of the Atacama Cosmology Telescope (ACT, data release 6);
P-ACT-LB adds CMB lensing (L) and DESI BAO (B)~\cite{louis2025}. The \emph{$H_0$ tension} is the disagreement between the Hubble constant
inferred from the CMB with $\Lambda$CDM, near 67\,km\,s$^{-1}$\,Mpc$^{-1}$, and the value from a distance ladder of Cepheids and supernovae,
near 73~\cite{planck18,riess2022}. A ladder calibrated on the tip of the red-giant branch lies in between~\cite{freedman2025}.
\end{primer}

\begin{primer}{Cosmology}{box:cosmo}
\emph{The expanding universe.} On large scales the universe is homogeneous and isotropic. Distances grow with a scale factor $a(t)$, light
from the past is redshifted by $1+z=a_0/a$, and the expansion rate $H=\dot a/a$ obeys, for a flat universe,
$H^2=8\pi G\rho/3+\Lambda c^2/3$. Each component's share is $\Omega_i=\rho_i/\rho_c$ with the critical density $\rho_c=3H^2/8\pi G$; the
value today is $H_0$.

\emph{$\Lambda$CDM.} The standard model of cosmology adds to the known particles a cold, collisionless \emph{dark matter} and a
cosmological constant $\Lambda$. Six parameters fix it: the densities of baryons and dark matter ($\Omega_bh^2$, $\Omega_ch^2$), the angular
size of the sound horizon, the optical depth to reionisation, and the amplitude $A_s$ and tilt $n_s$ of the primordial
fluctuations~\cite{planck18}.

\emph{Inflation.} An early epoch of nearly exponential expansion, $a\propto e^{Ht}$, flattens and smooths the universe and stretches
quantum fluctuations into the seeds of galaxies~\cite{guth}. It is driven by a scalar field, the \emph{inflaton}, with potential $V$, and
its length is counted in \emph{e-folds}, $N=\ln(a_{\rm end}/a_{\rm start})$. In \emph{slow roll} the field creeps down a flat potential;
with the reduced Planck mass $M_{\rm pl}$, the parameters $\epsilon=\tfrac12M_{\rm pl}^2(V'/V)^2$ and $\eta_V=M_{\rm pl}^2V''/V$ are small. The
fluctuations then have a nearly scale-invariant spectrum with tilt $n_s\approx1-6\epsilon+2\eta_V$, and gravitational waves with a
tensor-to-scalar ratio $r\approx16\epsilon$~\cite{planck2020x}.

\emph{Reheating and oscillons.} After inflation the inflaton oscillates about the minimum of its potential and decays into ordinary
particles, which \emph{reheats} the universe~\cite{kls}. In a broad class of models it first fragments into \emph{oscillons}, long-lived
localised lumps of the field~\cite{amin}.

\emph{The probes.} The cosmic microwave background (CMB) was released at recombination, $z\approx1100$. Its acoustic peaks are the imprint of
sound waves in the early plasma: their spacing fixes the sound horizon and their heights the densities~\cite{planck18}. The same sound
horizon, about 150\,Mpc today, is imprinted on the clustering of galaxies as \emph{baryon acoustic oscillations} (BAO), a standard ruler
for distance against redshift~\cite{desi2}.

\emph{Dark energy.} A cosmological constant has the equation of state $w=p/\rho c^2=-1$ at all times. DESI fits the
Chevallier--Polarski--Linder (CPL) form $w(a)=w_0+w_a(1-a)$~\cite{chevallier2001,linder2003}. A \emph{thawing} dark energy has $w_0>-1$
and $w_a<0$: it was close to $-1$ early and weakens today.

\emph{Data names and tensions.} P-ACT combines the CMB spectra of Planck and of the Atacama Cosmology Telescope (ACT, data release 6);
P-ACT-LB adds CMB lensing (L) and DESI BAO (B)~\cite{louis2025}. The \emph{$H_0$ tension} is the disagreement between the Hubble constant
inferred from the CMB with $\Lambda$CDM, near 67\,km\,s$^{-1}$\,Mpc$^{-1}$, and the value from a distance ladder of Cepheids and supernovae,
near 73~\cite{planck18,riess2022}. A ladder calibrated on the tip of the red-giant branch lies in between~\cite{freedman2025}.
\end{primer}

\begin{primer}{Cosmology}{box:cosmo}
\emph{The expanding universe.} On large scales the universe is homogeneous and isotropic. Distances grow with a scale factor $a(t)$, light
from the past is redshifted by $1+z=a_0/a$, and the expansion rate $H=\dot a/a$ obeys, for a flat universe,
$H^2=8\pi G\rho/3+\Lambda c^2/3$. Each component's share is $\Omega_i=\rho_i/\rho_c$ with the critical density $\rho_c=3H^2/8\pi G$; the
value today is $H_0$.

\emph{$\Lambda$CDM.} The standard model of cosmology adds to the known particles a cold, collisionless \emph{dark matter} and a
cosmological constant $\Lambda$. Six parameters fix it: the densities of baryons and dark matter ($\Omega_bh^2$, $\Omega_ch^2$), the angular
size of the sound horizon, the optical depth to reionisation, and the amplitude $A_s$ and tilt $n_s$ of the primordial
fluctuations~\cite{planck18}.

\emph{Inflation.} An early epoch of nearly exponential expansion, $a\propto e^{Ht}$, flattens and smooths the universe and stretches
quantum fluctuations into the seeds of galaxies~\cite{guth}. It is driven by a scalar field, the \emph{inflaton}, with potential $V$, and
its length is counted in \emph{e-folds}, $N=\ln(a_{\rm end}/a_{\rm start})$. In \emph{slow roll} the field creeps down a flat potential;
with the reduced Planck mass $M_{\rm pl}$, the parameters $\epsilon=\tfrac12M_{\rm pl}^2(V'/V)^2$ and $\eta_V=M_{\rm pl}^2V''/V$ are small. The
fluctuations then have a nearly scale-invariant spectrum with tilt $n_s\approx1-6\epsilon+2\eta_V$, and gravitational waves with a
tensor-to-scalar ratio $r\approx16\epsilon$~\cite{planck2020x}.

\emph{Reheating and oscillons.} After inflation the inflaton oscillates about the minimum of its potential and decays into ordinary
particles, which \emph{reheats} the universe~\cite{kls}. In a broad class of models it first fragments into \emph{oscillons}, long-lived
localised lumps of the field~\cite{amin}.

\emph{The probes.} The cosmic microwave background (CMB) was released at recombination, $z\approx1100$. Its acoustic peaks are the imprint of
sound waves in the early plasma: their spacing fixes the sound horizon and their heights the densities~\cite{planck18}. The same sound
horizon, about 150\,Mpc today, is imprinted on the clustering of galaxies as \emph{baryon acoustic oscillations} (BAO), a standard ruler
for distance against redshift~\cite{desi2}.

\emph{Dark energy.} A cosmological constant has the equation of state $w=p/\rho c^2=-1$ at all times. DESI fits the
Chevallier--Polarski--Linder (CPL) form $w(a)=w_0+w_a(1-a)$~\cite{chevallier2001,linder2003}. A \emph{thawing} dark energy has $w_0>-1$
and $w_a<0$: it was close to $-1$ early and weakens today.

\emph{Data names and tensions.} P-ACT combines the CMB spectra of Planck and of the Atacama Cosmology Telescope (ACT, data release 6);
P-ACT-LB adds CMB lensing (L) and DESI BAO (B)~\cite{louis2025}. The \emph{$H_0$ tension} is the disagreement between the Hubble constant
inferred from the CMB with $\Lambda$CDM, near 67\,km\,s$^{-1}$\,Mpc$^{-1}$, and the value from a distance ladder of Cepheids and supernovae,
near 73~\cite{planck18,riess2022}. A ladder calibrated on the tip of the red-giant branch lies in between~\cite{freedman2025}.
\end{primer}

\subsection{The strong split and the timeline}\label{sec:split}

The start (Hypothesis~\ref{h:minting}, Section~\ref{sec:start}) contains one more event that orders everything after it.

\begin{hypothesisprime6}[the strong split; refines Hypothesis~\ref{h:minting}]
The start also contains the split of the strong force from the electroweak force, and causality begins with that split. If the two
were ever one force, the unification lies before time: within causality the two were never one. \tagH
\end{hypothesisprime6}

\noindent The Standard Model agrees within causality: in it the strong and the electroweak force are separate gauge interactions
and never meet; run with the Standard Model's content alone, the three couplings do not meet at one point~\cite{amaldi}
(Section~\ref{sec:sm}). A unification
would leave three traces in the causal era, and they make the hypothesis testable in part \tagP. (1) \emph{Magnetic monopoles.} The split of
a unified force generically creates very heavy magnetic monopoles, about one per causal volume (Preskill~\cite{preskill}); had the split
happened after inflation, the universe would be full of them. None has been seen: MACRO bounds their flux below
$1.4\times10^{-16}$\,cm$^{-2}$\,s$^{-1}$\,sr$^{-1}$~\cite{macro}. The split must therefore happen before or during inflation, and reheating
must stay below the unification energy, which was one of Guth's reasons for inflation~\cite{guth}. Hypothesis~6$'$ puts it exactly there:
the ordering is not only allowed, it is required. (2) \emph{Proton decay.} The heavy gauge bosons of a unified force, of mass $M_X$, still
let protons decay after the split, with a lifetime of order $\tau\sim M_X^4/(\alpha_G^2m_p^5)$ to within a factor of about ten: about
$5\times10^{35}$ years for $M_X=10^{16}$\,GeV and $\alpha_G=1/40$ \chk{Q20}. Super-Kamiokande finds $\tau(p\to e^+\pi^0)>2.4\times10^{34}$
years~\cite{superk}, which requires $M_X\gtrsim5\times10^{15}$\,GeV in this estimate and excludes the simplest unified theory of Georgi and
Glashow~\cite{georgiglashow}; Hyper-Kamiokande will reach much further~\cite{hyperk} (Section~\ref{sec:data}, \chk{data:D8.1},
\chk{data:D8.2}). (3) \emph{The couplings.} A real unification needs
particles beyond the Standard Model that make the couplings meet. The bound on inflation, $V^{1/4}<1.4\times10^{16}$\,GeV
(Section~\ref{sec:tensor}), allows a split at the unification scale within the start \chk{Q20}. \emph{Verdict: holds as a reading}; the
ordering is required by the absence of monopoles; the unification itself is \emph{untested}, with proton decay as its live test.

\paragraph{The timeline.} Read together, Hypotheses~6, 6$'$ and Section~\ref{sec:higgsplit} give one sequence:
\begin{enumerate}[nosep]
\item \emph{The start.} The strong force splits from the electroweak force; inflation dilutes the monopoles of the split; the decay of the
oscillons reheats the universe to below the unification energy. Causality begins.
\item \emph{About $10^{-11}$\,s, at 160\,GeV.} The symmetric electroweak state becomes unstable; the Higgs vacuum takes over granting rest
mass to the W, the Z and the matter particles, and the information field, the photon and with it $\alpha\approx1/128$, splits out of the
electroweak field.
\item \emph{Afterwards.} $\alpha$ runs down to $1/137$ at low energies; nucleosynthesis, recombination and the late universe follow.
\end{enumerate}

% GUID v6, Fig. 4 (label fig:timeline): the timeline of Section sec:split. Input at %FIG:timeline in part1.tex: % GUID v6, Fig. 4 (label fig:timeline): the timeline of Section sec:split. Input at %FIG:timeline in part1.tex: % GUID v6, Fig. 4 (label fig:timeline): the timeline of Section sec:split. Input at %FIG:timeline in part1.tex: \input{figures/fig-timeline}
\begin{figure}[tbp]
\centering
\begin{tikzpicture}[x=1cm,y=1cm,font=\footnotesize,>={Stealth[length=2mm]},
  ev/.style={draw,rounded corners=2pt,fill=white,align=left,text width=2.85cm,inner sep=2.5pt,font=\scriptsize,anchor=north,
    minimum height=1.8cm},
  ttl/.style={align=center,text width=2.95cm,anchor=south,inner sep=1pt},
  tck/.style={draw,circle,fill=black,inner sep=0pt,minimum size=4pt},
  band/.style={fill=white,font=\scriptsize,inner sep=1.5pt},
]
% the start: a shaded region without internal order, closed by a jagged edge
\fill[blue!7] (0,-3.35) rectangle (3.2,1.4);
\draw[blue!45!black] (3.2,1.4) \foreach \k in {1,...,38}{ -- ({3.2+0.05*pow(-1,\k)},{1.4-0.125*\k}) };
% the time axis
\draw[->,thick] (0,0) -- (15.95,0);
\foreach \x in {1.6,4.8,8.0,11.2,14.4}{\node[tck] at (\x,0) {};}
% titles above the axis
\node[ttl] at (1.6,0.15) {\textbf{The start}\\ {\scriptsize the minting (H6, H6$'$)}\\[-1pt] {\scriptsize no internal order}};
\node[ttl] at (4.8,0.15) {\textbf{Electroweak split}\\ {\scriptsize about $10^{-11}$\,s}\\[-1pt] {\scriptsize $T\approx160$\,GeV}};
\node[ttl] at (8.0,0.15) {\textbf{Nucleosynthesis}\\ {\scriptsize 1\,s to minutes}\\[-1pt] {\scriptsize $T\approx1$ to 0.1\,MeV}};
\node[ttl] at (11.2,0.15) {\textbf{Recombination}\\ {\scriptsize about $3.8\times10^5$\,yr}\\[-1pt]
  {\scriptsize $z\approx1100$, $T\approx0.3$\,eV}};
\node[ttl] at (14.4,0.15) {\textbf{Late universe}\\ {\scriptsize accelerating, $z\lesssim0.6$}\\[-1pt]
  {\scriptsize today 13.8\,Gyr}};
% what happens, below the axis
\node[ev] at (1.6,-0.22) {The strong force splits off; inflation dilutes its monopoles; oscillons decay and reheat.};
\node[ev] at (4.8,-0.22) {The Higgs vacuum grants rest mass to W, Z and matter; the photon splits out, with $\alpha\approx1/128$.};
\node[ev] at (8.0,-0.22) {Light nuclei form. Below the electron mass, $\alpha$ has run down to $1/137$.};
\node[ev] at (11.2,-0.22) {Atoms form; the microwave background is released; the dark matter is already there.};
\node[ev] at (14.4,-0.22) {Dark energy dominates; the expansion tends to de Sitter, read as the balanced state.};
% bands: what the axioms say about the geometry, and alpha
\node[font=\scriptsize,align=center,text width=2.95cm,anchor=north,inner sep=0pt] at (1.6,-2.4) {geometry may be in superposition;\\
  the axioms are silent};
\draw[|-|,black!60] (3.2,-2.55) -- node[band] {causality begins; the frame is classical: every geometry is recorded (Lemma~\ref{lem:classical})}
  (15.85,-2.55);
\draw[|-|,black!60] (3.2,-3.05) -- node[band] {no $\alpha$ yet} (4.8,-3.05);
\draw[black!60] (4.8,-3.05) -- node[band] {$\alpha$ runs: $1/128\to1/137$} (8.0,-3.05) -- node[band] {$\alpha\approx1/137$}
  (15.85,-3.05);
\draw[black!60] (15.85,-3.15) -- (15.85,-2.95);
\end{tikzpicture}
\caption{The timeline of Section~\ref{sec:split}, not to scale. The start (shaded; Hypotheses~6 and~6$'$) has no internal order the axioms
can refer to; causality and the classical frame begin as it completes. The electroweak split follows at about $10^{-11}$\,s and
160\,GeV~\cite{donofrio}: the Higgs vacuum grants rest mass, and the photon, and with it $\alpha$, splits out of the electroweak field.
Nucleosynthesis, recombination and the late universe are those of standard cosmology~\cite{navas2024,planck18}.}
\label{fig:timeline}
\end{figure}

\begin{figure}[tbp]
\centering
\begin{tikzpicture}[x=1cm,y=1cm,font=\footnotesize,>={Stealth[length=2mm]},
  ev/.style={draw,rounded corners=2pt,fill=white,align=left,text width=2.85cm,inner sep=2.5pt,font=\scriptsize,anchor=north,
    minimum height=1.8cm},
  ttl/.style={align=center,text width=2.95cm,anchor=south,inner sep=1pt},
  tck/.style={draw,circle,fill=black,inner sep=0pt,minimum size=4pt},
  band/.style={fill=white,font=\scriptsize,inner sep=1.5pt},
]
% the start: a shaded region without internal order, closed by a jagged edge
\fill[blue!7] (0,-3.35) rectangle (3.2,1.4);
\draw[blue!45!black] (3.2,1.4) \foreach \k in {1,...,38}{ -- ({3.2+0.05*pow(-1,\k)},{1.4-0.125*\k}) };
% the time axis
\draw[->,thick] (0,0) -- (15.95,0);
\foreach \x in {1.6,4.8,8.0,11.2,14.4}{\node[tck] at (\x,0) {};}
% titles above the axis
\node[ttl] at (1.6,0.15) {\textbf{The start}\\ {\scriptsize the minting (H6, H6$'$)}\\[-1pt] {\scriptsize no internal order}};
\node[ttl] at (4.8,0.15) {\textbf{Electroweak split}\\ {\scriptsize about $10^{-11}$\,s}\\[-1pt] {\scriptsize $T\approx160$\,GeV}};
\node[ttl] at (8.0,0.15) {\textbf{Nucleosynthesis}\\ {\scriptsize 1\,s to minutes}\\[-1pt] {\scriptsize $T\approx1$ to 0.1\,MeV}};
\node[ttl] at (11.2,0.15) {\textbf{Recombination}\\ {\scriptsize about $3.8\times10^5$\,yr}\\[-1pt]
  {\scriptsize $z\approx1100$, $T\approx0.3$\,eV}};
\node[ttl] at (14.4,0.15) {\textbf{Late universe}\\ {\scriptsize accelerating, $z\lesssim0.6$}\\[-1pt]
  {\scriptsize today 13.8\,Gyr}};
% what happens, below the axis
\node[ev] at (1.6,-0.22) {The strong force splits off; inflation dilutes its monopoles; oscillons decay and reheat.};
\node[ev] at (4.8,-0.22) {The Higgs vacuum grants rest mass to W, Z and matter; the photon splits out, with $\alpha\approx1/128$.};
\node[ev] at (8.0,-0.22) {Light nuclei form. Below the electron mass, $\alpha$ has run down to $1/137$.};
\node[ev] at (11.2,-0.22) {Atoms form; the microwave background is released; the dark matter is already there.};
\node[ev] at (14.4,-0.22) {Dark energy dominates; the expansion tends to de Sitter, read as the balanced state.};
% bands: what the axioms say about the geometry, and alpha
\node[font=\scriptsize,align=center,text width=2.95cm,anchor=north,inner sep=0pt] at (1.6,-2.4) {geometry may be in superposition;\\
  the axioms are silent};
\draw[|-|,black!60] (3.2,-2.55) -- node[band] {causality begins; the frame is classical: every geometry is recorded (Lemma~\ref{lem:classical})}
  (15.85,-2.55);
\draw[|-|,black!60] (3.2,-3.05) -- node[band] {no $\alpha$ yet} (4.8,-3.05);
\draw[black!60] (4.8,-3.05) -- node[band] {$\alpha$ runs: $1/128\to1/137$} (8.0,-3.05) -- node[band] {$\alpha\approx1/137$}
  (15.85,-3.05);
\draw[black!60] (15.85,-3.15) -- (15.85,-2.95);
\end{tikzpicture}
\caption{The timeline of Section~\ref{sec:split}, not to scale. The start (shaded; Hypotheses~6 and~6$'$) has no internal order the axioms
can refer to; causality and the classical frame begin as it completes. The electroweak split follows at about $10^{-11}$\,s and
160\,GeV~\cite{donofrio}: the Higgs vacuum grants rest mass, and the photon, and with it $\alpha$, splits out of the electroweak field.
Nucleosynthesis, recombination and the late universe are those of standard cosmology~\cite{navas2024,planck18}.}
\label{fig:timeline}
\end{figure}

\begin{figure}[tbp]
\centering
\begin{tikzpicture}[x=1cm,y=1cm,font=\footnotesize,>={Stealth[length=2mm]},
  ev/.style={draw,rounded corners=2pt,fill=white,align=left,text width=2.85cm,inner sep=2.5pt,font=\scriptsize,anchor=north,
    minimum height=1.8cm},
  ttl/.style={align=center,text width=2.95cm,anchor=south,inner sep=1pt},
  tck/.style={draw,circle,fill=black,inner sep=0pt,minimum size=4pt},
  band/.style={fill=white,font=\scriptsize,inner sep=1.5pt},
]
% the start: a shaded region without internal order, closed by a jagged edge
\fill[blue!7] (0,-3.35) rectangle (3.2,1.4);
\draw[blue!45!black] (3.2,1.4) \foreach \k in {1,...,38}{ -- ({3.2+0.05*pow(-1,\k)},{1.4-0.125*\k}) };
% the time axis
\draw[->,thick] (0,0) -- (15.95,0);
\foreach \x in {1.6,4.8,8.0,11.2,14.4}{\node[tck] at (\x,0) {};}
% titles above the axis
\node[ttl] at (1.6,0.15) {\textbf{The start}\\ {\scriptsize the minting (H6, H6$'$)}\\[-1pt] {\scriptsize no internal order}};
\node[ttl] at (4.8,0.15) {\textbf{Electroweak split}\\ {\scriptsize about $10^{-11}$\,s}\\[-1pt] {\scriptsize $T\approx160$\,GeV}};
\node[ttl] at (8.0,0.15) {\textbf{Nucleosynthesis}\\ {\scriptsize 1\,s to minutes}\\[-1pt] {\scriptsize $T\approx1$ to 0.1\,MeV}};
\node[ttl] at (11.2,0.15) {\textbf{Recombination}\\ {\scriptsize about $3.8\times10^5$\,yr}\\[-1pt]
  {\scriptsize $z\approx1100$, $T\approx0.3$\,eV}};
\node[ttl] at (14.4,0.15) {\textbf{Late universe}\\ {\scriptsize accelerating, $z\lesssim0.6$}\\[-1pt]
  {\scriptsize today 13.8\,Gyr}};
% what happens, below the axis
\node[ev] at (1.6,-0.22) {The strong force splits off; inflation dilutes its monopoles; oscillons decay and reheat.};
\node[ev] at (4.8,-0.22) {The Higgs vacuum grants rest mass to W, Z and matter; the photon splits out, with $\alpha\approx1/128$.};
\node[ev] at (8.0,-0.22) {Light nuclei form. Below the electron mass, $\alpha$ has run down to $1/137$.};
\node[ev] at (11.2,-0.22) {Atoms form; the microwave background is released; the dark matter is already there.};
\node[ev] at (14.4,-0.22) {Dark energy dominates; the expansion tends to de Sitter, read as the balanced state.};
% bands: what the axioms say about the geometry, and alpha
\node[font=\scriptsize,align=center,text width=2.95cm,anchor=north,inner sep=0pt] at (1.6,-2.4) {geometry may be in superposition;\\
  the axioms are silent};
\draw[|-|,black!60] (3.2,-2.55) -- node[band] {causality begins; the frame is classical: every geometry is recorded (Lemma~\ref{lem:classical})}
  (15.85,-2.55);
\draw[|-|,black!60] (3.2,-3.05) -- node[band] {no $\alpha$ yet} (4.8,-3.05);
\draw[black!60] (4.8,-3.05) -- node[band] {$\alpha$ runs: $1/128\to1/137$} (8.0,-3.05) -- node[band] {$\alpha\approx1/137$}
  (15.85,-3.05);
\draw[black!60] (15.85,-3.15) -- (15.85,-2.95);
\end{tikzpicture}
\caption{The timeline of Section~\ref{sec:split}, not to scale. The start (shaded; Hypotheses~6 and~6$'$) has no internal order the axioms
can refer to; causality and the classical frame begin as it completes. The electroweak split follows at about $10^{-11}$\,s and
160\,GeV~\cite{donofrio}: the Higgs vacuum grants rest mass, and the photon, and with it $\alpha$, splits out of the electroweak field.
Nucleosynthesis, recombination and the late universe are those of standard cosmology~\cite{navas2024,planck18}.}
\label{fig:timeline}
\end{figure}

\subsection{Primordial gravitational waves}\label{sec:tensor}

During inflation the vacuum fluctuations of the metric are stretched into a nearly scale-invariant background of gravitational waves
(Starobinsky~\cite{starobinsky79}). Their power spectrum is
\begin{equation}\label{eq:tensor}
\mathcal P_t=\frac{2}{\pi^2}\left(\frac{\hbar H}{M_{\text{pl}}c^2}\right)^2=\frac{16\,G\hbar H^2}{\pi c^5},
\end{equation}
with $M_{\text{pl}}^2=\hbar c/8\pi G$, linear in $\hbar$ and zero for a classical metric \chk{Q14}; its ratio to the scalar spectrum is
$r=16\epsilon$ in single-field slow roll, with the consistency relation $n_t=-r/8$ for its tilt. Because~\eqref{eq:tensor} vanishes as
$\hbar\to0$, Krauss and Wilczek argued that detecting this background would establish that the gravitational field was quantised during
inflation~\cite{krausswilczek} \tagP. The current bound is $r<0.036$ at 95\% confidence (BICEP/Keck~\cite{bk18}; $r<0.032$ combined with
Planck and baryon acoustic oscillations~\cite{tristram}), which puts the energy scale of inflation below $V^{1/4}=1.4\times10^{16}$\,GeV
and its Hubble rate below $\hbar H=4.7\times10^{13}$\,GeV \chk{Q14}.

\emph{The inflaton is not GUID's.} The action~\eqref{eq:action} contains no inflaton, and Hypothesis~\ref{h:minting} places inflation
within the start, where the axioms place no restriction. GUID therefore takes the inflaton and its potential from physics. The models that
fit the scalar spectrum best are plateau potentials, among them Starobinsky's $R^2$ inflation~\cite{starobinsky80} and Higgs
inflation~\cite{bezrukov}; both predict $n_s=1-2/N$ and $r\approx12/N^2\approx3\times10^{-3}$ at $N=60$ e-folds \chk{Q14}, the target of the
next generation of polarisation surveys~\cite{litebird}. Both are written with a scalar coupled to curvature in one frame and as a
minimally coupled scalar in another. That is the form Proposition~\ref{pr:gamma} excludes for the frame after the start; within the start
the axioms do not constrain it, and afterwards the extra scalar of $R^2$ inflation is heavy, near $10^{13}$\,GeV, so it leaves $\gamma=1$
in the solar system untouched. We score the plateau prediction conditionally on that choice (Section~\ref{sec:data}).

Under Hypothesis~\ref{h:minting} the tensor modes are generated within the start, where the axioms do not constrain the geometry, and
become classical, as the standard account says they do through squeezing and decoherence~\cite{polarski}, as the start completes. GUID
therefore allows the standard spectrum and its consistency relation, and a detection would be the first observation of the quantum
geometry of the start. GUID's own predictions about the geometry begin after the start. \emph{Verdict: holds as a reading}; GUID
predicts nothing about tensor modes beyond the inflation model it borrows.

\subsection{The late universe: the balanced state}

In the flat $\Lambda$CDM model that fits the data ($\Omega_m=0.315$, $\Omega_\Lambda=0.685$~\cite{planck18}), the Hubble rate falls at every
moment from the past into the far future and tends to the de Sitter rate $H_\infty=\sqrt{\Omega_\Lambda}\,H_0\approx0.83\,H_0$; the expansion
has accelerated since $z=(2\Omega_\Lambda/\Omega_m)^{1/3}-1\approx0.63$ \chk{Q13}. A universe settling into de Sitter expansion cools towards
an empty final state. GUID reads this as the cosmological balanced state, the equilibrium in which the need for further existence has
fallen to a constant. \emph{Verdict: holds as a reading} in $\Lambda$CDM. The newest data leave the end open: an evolving dark energy,
the Hubble constant and the curvature are tensions of $\Lambda$CDM that GUID inherits unchanged (Section~\ref{sec:data}).

\section{Dark matter: remnants trapped in the gravitational field}\label{sec:darkmatter}

\begin{hypothesis}[trapped remnants]\label{h:remnant}
The modelling of this paper explicitly allows remnants of interactions to be trapped in the gravitational field: energy and momentum that
interactions have left in the frame, or behind a boundary the frame has closed, continue to gravitate after the interaction.
\tagH
\end{hypothesis}

DMP motivates the hypothesis twice: a record outlives its act (Axiom~3), and a containing boundary keeps what it holds while every outside
reader sees only its counts (T31). Dark matter, as far as the data show, is matter of exactly that kind: it gravitates and does nothing
else that has been seen. Within the classical-frame modelling there are three routes by which remnants are trapped, and each can be
calculated.

\paragraph{What every remnant must satisfy.} Four measurements constrain every route \tagP. (1) It must already be there at recombination:
the acoustic peaks of the microwave background measure the dark matter density at $z\approx1100$, before the first star, as
$\Omega_ch^2=0.120$, the same comoving value later structure requires, with the baryons then already at the share
$\Omega_b/\Omega_m=0.157$~\cite{planck18} \chk{Q15}. So the remnant formed before matter-radiation equality, consistent with formation in the start; the baryon content was fixed early, as nucleosynthesis and the microwave background agree. (2) It must be collisionless and move with the galaxies: in the
merging cluster 1E0657$-$558 the lensing mass follows the galaxies, not the hot gas that holds most of the baryons (Clowe et
al.~\cite{clowe}). (3) It cannot be the known neutrinos: they are less than $0.6\%$ of the dark matter density, since cosmology bounds
their mass sum below $0.064$\,eV~\cite{desi2}; a dwarf galaxy made of them would need each to weigh more than about 170\,eV by the Pauli
principle (Tremaine and Gunn~\cite{tremainegunn}); and relic neutrinos move at about 3200\,km/s and stellar ones at nearly $c$, against a
Galactic escape speed of 533\,km/s~\cite{piffl} \chk{Q15}. (4) It must leave $\alpha$ unchanged to about a part per million over ten
billion years (Section~\ref{sec:alpha}).


\paragraph{The stochastic remnant.} Corollary~\ref{co:diffuse} makes the frame diffuse: every interaction of quantum matter with the
classical frame leaves a random kick in the field. Oppenheim and Russo computed the expectation value of the metric in classical-quantum
gravity from its path-integral formulation, for static, spherically symmetric configurations~\cite{oppenheimrusso}. Around a mass $M$ the
Newtonian potential acquires a linear and a quadratic term,
\begin{equation}\label{eq:stochpot}
\Phi(r)=-\frac{GM}{r}-\frac{\gamma_0}{2}-\frac{\gamma_1}{2}\,r-\frac{\gamma_2}{2}\,r^2 ,
\end{equation}
where the quadratic term is the cosmological constant of Schwarzschild--de Sitter space, $\gamma_2=\Lambda c^2/3>0$, which repels, and
the linear term is anticorrelated with it: its mean is negative, $\gamma_1<0$, so it attracts, with $|\gamma_1|$ of order $c^2\sqrt\Lambda$.
Physically, where accelerations are low the variance of the field's acceleration exceeds the square of its mean and acts as an additional,
entropic attraction. The circular speed is
\begin{equation}\label{eq:vcirc}
v^2(r)=r\,\frac{\dd\Phi}{\dd r}=\frac{GM}{r}+\frac{|\gamma_1|}{2}\,r-\gamma_2\,r^2 .
\end{equation}
The first two terms balance at $r_*=\sqrt{2GM/|\gamma_1|}$, where
\begin{equation}\label{eq:btf}
v_*^4=2GM|\gamma_1| .
\end{equation}
Around $r_*$ the rotation curve is roughly flat, changing by about 12\% between $r_*/2$ and $2r_*$, and its height obeys $v^4\propto M$,
the form of the observed baryonic Tully--Fisher relation. For a galaxy of $10^{11}$ solar masses in baryons and $|\gamma_1|=10^{-10}$\,m/s$^2$,
the order Oppenheim and Russo find, $r_*=17$\,kpc and $v_*=227$\,km/s, with the $\Lambda$ term below $10^{-4}$ of $v_*^2$ there \chk{Q18}.
The acceleration scale is tied to the cosmological constant: $c^2\sqrt{\Lambda/3}/2\pi=8.6\times10^{-11}$\,m/s$^2$, close to Milgrom's
empirical $a_0=1.2\times10^{-10}$\,m/s$^2$~\cite{milgrom} \chk{Q18}. Such a closeness is Milgrom's known coincidence, not support:
GUID supplies no factor $2\pi$ (Section~\ref{sec:remnant} counts how often a random scale comes as close). On this route a single mechanism, the frame's diffusion, would give both a
dark-energy term and a dark-matter-like term. The calculation is preliminary, as its authors stress: it is static and spherically
symmetric, depends on an infrared regularisation, fixes $\gamma_1$ by boundary conditions rather than by a coupling constant, and does not
yet address gravitational lensing or the microwave background, and Hertzberg and Loeb dispute that it can produce flat rotation
curves at all~\cite{hertzberg2024}. A deterministic relative of the same idea exists in tetrad gravity: Hehl
and Mashhoon found that a memory of the frame, a nonlocal dependence on its past, acts in the linear approximation as dark matter given by
an integral transform of the matter~\cite{hehlmashhoon}; such a memory would make the effect depend on each galaxy's history.



\paragraph{Bridge A: the outermost horizon sets the scale.} The first of GUID's two bridges (Section~\ref{sec:bridges}) fixes the
stochastic remnant's scale. It depends on the outer
cut-off $r_{\max}$ of the path integral. Oppenheim and Russo find, in units with $c=1$, the mean
$\gamma_1=-\tfrac92\,\gamma_2\,r_{\max}\,(1-4\beta)/(5-18\beta)$, with $\gamma_2=\Lambda/3$ and a parameter $\beta<0$ of their action to which
the result is insensitive~\cite{oppenheimrusso}. In GUID the outermost containing boundary is the cosmological event horizon, which in
$\Lambda$CDM tends to the de Sitter radius $r_{\text{dS}}=c/H_\infty=1.66\times10^{26}$\,m. With $r_{\max}=r_{\text{dS}}$ and $c$ restored,
\begin{equation}\label{eq:bridgeA}
|\gamma_1|=(0.90\ \text{to}\ 1.00)\;cH_\infty=(4.9\ \text{to}\ 5.4)\times10^{-10}\,\text{m/s}^2,
\end{equation}
with no free parameter \chk{Q19}. It has the right order of magnitude and is too strong. Matching the turnover speed~\eqref{eq:btf} to the
observed baryonic Tully--Fisher relation $v^4=GMa_0$ requires $|\gamma_1|=a_0/2=6\times10^{-11}$\,m/s$^2$; Bridge~A exceeds this by a factor
of 8 in $v^4$, 1.7 in $v$. For a galaxy of $10^{11}$ solar masses it gives 340\,km/s at $r_*=7.6$\,kpc, where the relation gives 200\,km/s,
and the linear term lets the curve rise to 620\,km/s at 50\,kpc, where observed curves stay flat \chk{Q19}. The estimate already rejects
it; Section~\ref{sec:sparc} confronts it with the data. Its near miss is Milgrom's coincidence $a_0\approx cH/2\pi$, which is known.

\paragraph{The contained remnant.} The second route is containment (Hypothesis~\ref{h:horizon}). Energy that collapsed behind a horizon
during the start gravitates through its counts alone and is dark by T31: primordial black holes. The start itself makes them. When the
inflaton fragments into oscillons, which in a broad class of models dominate the universe for a while after inflation~\cite{amin}, enough
of them can collapse into black holes of sublunar mass to make up all of the dark matter (Cotner, Kusenko and Takhistov~\cite{cotner}).
Primordial black holes can be all of the dark matter only in the asteroid-mass window, about $10^{14}$ to $10^{19}$\,kg ($10^{17}$ to
$10^{22}$\,g), which current constraints leave open~\cite{gortongreen}. The window's light edge is set by the Hawking part of their HPI
radiation: a hole of $10^{14}$\,kg glows at $1.2\times10^9$\,K, $k_BT\approx0.1$\,MeV, and its evaporation products (MeV gamma rays, and
positrons that annihilate at 511\,keV) bound how many such holes there can be; holes lighter than a few $10^{11}$\,kg have evaporated
already \chk{Q18}. At the heavy edge, $10^{19}$\,kg with a Schwarzschild radius of 15\,nm and a temperature of $1.2\times10^4$\,K,
microlensing takes over \chk{Q18}. The prediction of this route is specific: if dark matter is the contained remnant of the start, the
Hawking part of HPI radiation, unobservable for astrophysical black holes (Section~\ref{sec:bh}), becomes observable at the window's light
edge. The window is taken from current constraints and the formation from Cotner et al.; this paper does not test either.

\paragraph{The radiated remnant.} Gravitational waves are remnants of interactions that propagate in the frame. Pulsar timing has found a
nanohertz background with characteristic strain $h_c\approx2.4\times10^{-15}$ at the frequency $1/\text{yr}$ (NANOGrav~\cite{nanograv}),
which corresponds to $\Omega_{\text{gw}}=2\pi^2f^2h_c^2/3H_0^2\approx8\times10^{-9}$ per logarithmic frequency \chk{Q18}: more than $10^7$
times too little for the dark matter, and as radiation ($w=1/3$) it does not cluster. It is a real remnant trapped in the field, but not
the dark matter.

\paragraph{What Part III tests.} Section~\ref{sec:sparc} confronts the stochastic remnant and Bridge~A with the rotation curves of the
SPARC galaxies, and Section~\ref{sec:remnant} an accumulated version, cold-dark-matter halos and the change with cosmic time; it also
holds the dark-matter verdict. Section~\ref{sec:data} confronts the other requirements.

\section{The information scalar and matter}\label{sec:matter}

\subsection{The fine-structure constant as the information field}\label{sec:alpha}

In GUID $\alpha$ is set by the information scalar $\phi$ (Hypothesis~\ref{h:info}), so it may vary where $\phi$ does. Two statements
follow, and Sections~\ref{sec:data} and~\ref{sec:sm} test them.

\emph{The forward statement: $\alpha$ does not grow.} With the action~\eqref{eq:action} as written, with no potential for $\phi$, this
is \emph{derived}. The scalar is then driven only by the electrostatic energy of the baryons, as in Bekenstein's theory, so the drift
$\dot\alpha/\alpha$ is negative, proportional to $\lambda^2$ and to the small electrostatic share $f_b\lesssim10^{-3}$ of baryonic mass,
and at the bound~\eqref{eq:lambdabound} far below any clock. A drift large enough to measure needs a potential for $\phi$, and then the
sign is not predicted. A measured positive drift would refute the action as written. \emph{Verdict: derived; untested.}

\emph{Only $\alpha$ varies.} Because the scalar multiplies only the Maxwell term, the proton-to-electron mass ratio $\mu=m_p/m_e$ varies
only through the electromagnetic self-energies of the proton and the electron. The ratio $R=\dd\ln\mu/\dd\ln\alpha$ is then small, at most
a few hundredths in magnitude, against $|R|$ of 20 to 50 in models where all couplings vary together; its exact value depends on the
scale at which the Yukawa couplings are held fixed (Section~\ref{sec:sm}). If $\alpha$ is ever seen to vary, $\mu$ must vary far less.
\emph{Verdict: live.}

\emph{$\alpha$ has a beginning.} Physically $\alpha$ does not exist before the electroweak split: there is no photon then, only the
couplings of the electroweak field. It is born at the split, near $1/128$ at the Z mass, and runs down to $1/137$ at low energies
(Section~\ref{sec:higgsplit}; \chk{Q16}, \chk{Q20}). A value of one at any physical energy is beyond known physics: the Standard Model
gives $1/\alpha\approx105$ even at the Planck energy \chk{Q16}. \emph{Verdict: holds.}

% GUID v6, Primer P9 (label box:particle): particle physics. Background only; it states no result of the paper.
% Place it before "The spin ladder" (marker %PRIMER:particle in part1.tex): % GUID v6, Primer P9 (label box:particle): particle physics. Background only; it states no result of the paper.
% Place it before "The spin ladder" (marker %PRIMER:particle in part1.tex): % GUID v6, Primer P9 (label box:particle): particle physics. Background only; it states no result of the paper.
% Place it before "The spin ladder" (marker %PRIMER:particle in part1.tex): \input{primers/p9-particle}
\begin{primer}{Particle physics}{box:particle}
\setlength{\parskip}{3pt plus 1pt}%
\emph{The Higgs mechanism.} The Standard Model is a gauge theory of $SU(3)\times SU(2)\times U(1)$: colour for the strong force, weak isospin
and hypercharge for the electroweak force~\cite{glashow,weinberg67,salam}. The Higgs field is a doublet whose potential has its minimum
away from zero, at a vacuum value $v=246$\,GeV~\cite{higgs64,be64}. In that vacuum the W and Z bosons acquire masses proportional to $v$,
while one combination of the electroweak fields, the photon, stays massless. Quarks and charged leptons get their masses from
\emph{Yukawa couplings} $y_f$ to the doublet, $m_f=y_fv/\sqrt2$. The Higgs boson therefore couples to each fermion with strength $m_f/v$
and to the W and Z with $2m_V^2/v$.

\emph{The $\kappa$ framework.} Measurements of Higgs production and decay are summarised by \emph{coupling modifiers} $\kappa_i$, which
scale each Standard Model coupling, so that a rate involving particle $i$ scales as $\kappa_i^2$; the Standard Model has every
$\kappa_i=1$~\cite{david2012}.

\emph{Schemes and running.} A coupling or mass in quantum field theory is defined only together with a \emph{renormalisation scheme} and
an energy scale $\mu$. In the $\overline{\text{MS}}$ scheme only the divergent part of each loop is subtracted, and parameters run with
$\mu$. The \emph{on-shell} scheme instead fixes masses at the poles of the propagators and $\alpha$ at zero momentum transfer. The same
quantity can differ between the two by more than its measurement error, so every comparison states its scheme~\cite{erler2017,navas2024}.
The dependence on $\mu$ obeys the renormalisation-group equations (RGEs). At one loop
$\dd\alpha_i^{-1}/\dd\ln\mu=-b_i/2\pi$, with $(b_1,b_2,b_3)=(41/10,-19/6,-7)$ in the Standard Model (hypercharge normalised as in grand
unification). The strong coupling therefore weakens at high energy, and the hypercharge coupling grows. Precision work runs at two or
three loops~\cite{buttazzo2013}.

\emph{Dirac or Majorana.} A \emph{Dirac} mass joins a left-handed neutrino to a right-handed partner and conserves lepton number $L$. A
\emph{Majorana} mass joins the neutrino to its own antiparticle and changes $L$ by two. The lowest-dimension operator of Standard Model
fields that gives neutrinos mass is Weinberg's dimension-five operator, $(c/\Lambda)\,(LH)(LH)$, with $L$ the lepton doublet; in the Higgs
vacuum it gives a Majorana mass of order $cv^2/\Lambda$, small when the new scale $\Lambda$ is large~\cite{weinberg1979}. Majorana masses
can therefore also come from the Higgs vacuum.

\emph{Neutrinoless double-beta decay} ($0\nu\beta\beta$) turns two neutrons of a nucleus into two protons and two electrons with no
neutrinos. It changes $L$ by two and can occur only if neutrinos have a Majorana mass. Through light-neutrino exchange its rate is
proportional to $m_{\beta\beta}^2$, with the effective mass $m_{\beta\beta}=|\sum_iU_{ei}^2m_i|$ set by the mixing matrix $U$ and the
masses $m_i$~\cite{navas2024}. Xenon and germanium experiments search for it~\cite{kamlandzen2025,legend2025}.

\emph{Grand unification.} A grand unified theory (GUT) embeds the three gauge groups in one simple group, the smallest being
$SU(5)$~\cite{georgiglashow}, with a single coupling above a scale $M_X$. With the Standard Model's particles alone, the three running
couplings do not meet at one point~\cite{amaldi}. A GUT has two generic consequences. Its heavy gauge bosons let protons decay, with a
lifetime of order $M_X^4/(\alpha_G^2m_p^5)$; Super-Kamiokande's bound on $p\to e^+\pi^0$ excludes the minimal $SU(5)$
model~\cite{superk}. And its breaking leaves magnetic monopoles of mass near $M_X/\alpha_G$, about one per causal volume, which would
overfill the universe unless inflation dilutes them~\cite{preskill,guth}.
\end{primer}

\begin{primer}{Particle physics}{box:particle}
\setlength{\parskip}{3pt plus 1pt}%
\emph{The Higgs mechanism.} The Standard Model is a gauge theory of $SU(3)\times SU(2)\times U(1)$: colour for the strong force, weak isospin
and hypercharge for the electroweak force~\cite{glashow,weinberg67,salam}. The Higgs field is a doublet whose potential has its minimum
away from zero, at a vacuum value $v=246$\,GeV~\cite{higgs64,be64}. In that vacuum the W and Z bosons acquire masses proportional to $v$,
while one combination of the electroweak fields, the photon, stays massless. Quarks and charged leptons get their masses from
\emph{Yukawa couplings} $y_f$ to the doublet, $m_f=y_fv/\sqrt2$. The Higgs boson therefore couples to each fermion with strength $m_f/v$
and to the W and Z with $2m_V^2/v$.

\emph{The $\kappa$ framework.} Measurements of Higgs production and decay are summarised by \emph{coupling modifiers} $\kappa_i$, which
scale each Standard Model coupling, so that a rate involving particle $i$ scales as $\kappa_i^2$; the Standard Model has every
$\kappa_i=1$~\cite{david2012}.

\emph{Schemes and running.} A coupling or mass in quantum field theory is defined only together with a \emph{renormalisation scheme} and
an energy scale $\mu$. In the $\overline{\text{MS}}$ scheme only the divergent part of each loop is subtracted, and parameters run with
$\mu$. The \emph{on-shell} scheme instead fixes masses at the poles of the propagators and $\alpha$ at zero momentum transfer. The same
quantity can differ between the two by more than its measurement error, so every comparison states its scheme~\cite{erler2017,navas2024}.
The dependence on $\mu$ obeys the renormalisation-group equations (RGEs). At one loop
$\dd\alpha_i^{-1}/\dd\ln\mu=-b_i/2\pi$, with $(b_1,b_2,b_3)=(41/10,-19/6,-7)$ in the Standard Model (hypercharge normalised as in grand
unification). The strong coupling therefore weakens at high energy, and the hypercharge coupling grows. Precision work runs at two or
three loops~\cite{buttazzo2013}.

\emph{Dirac or Majorana.} A \emph{Dirac} mass joins a left-handed neutrino to a right-handed partner and conserves lepton number $L$. A
\emph{Majorana} mass joins the neutrino to its own antiparticle and changes $L$ by two. The lowest-dimension operator of Standard Model
fields that gives neutrinos mass is Weinberg's dimension-five operator, $(c/\Lambda)\,(LH)(LH)$, with $L$ the lepton doublet; in the Higgs
vacuum it gives a Majorana mass of order $cv^2/\Lambda$, small when the new scale $\Lambda$ is large~\cite{weinberg1979}. Majorana masses
can therefore also come from the Higgs vacuum.

\emph{Neutrinoless double-beta decay} ($0\nu\beta\beta$) turns two neutrons of a nucleus into two protons and two electrons with no
neutrinos. It changes $L$ by two and can occur only if neutrinos have a Majorana mass. Through light-neutrino exchange its rate is
proportional to $m_{\beta\beta}^2$, with the effective mass $m_{\beta\beta}=|\sum_iU_{ei}^2m_i|$ set by the mixing matrix $U$ and the
masses $m_i$~\cite{navas2024}. Xenon and germanium experiments search for it~\cite{kamlandzen2025,legend2025}.

\emph{Grand unification.} A grand unified theory (GUT) embeds the three gauge groups in one simple group, the smallest being
$SU(5)$~\cite{georgiglashow}, with a single coupling above a scale $M_X$. With the Standard Model's particles alone, the three running
couplings do not meet at one point~\cite{amaldi}. A GUT has two generic consequences. Its heavy gauge bosons let protons decay, with a
lifetime of order $M_X^4/(\alpha_G^2m_p^5)$; Super-Kamiokande's bound on $p\to e^+\pi^0$ excludes the minimal $SU(5)$
model~\cite{superk}. And its breaking leaves magnetic monopoles of mass near $M_X/\alpha_G$, about one per causal volume, which would
overfill the universe unless inflation dilutes them~\cite{preskill,guth}.
\end{primer}

\begin{primer}{Particle physics}{box:particle}
\setlength{\parskip}{3pt plus 1pt}%
\emph{The Higgs mechanism.} The Standard Model is a gauge theory of $SU(3)\times SU(2)\times U(1)$: colour for the strong force, weak isospin
and hypercharge for the electroweak force~\cite{glashow,weinberg67,salam}. The Higgs field is a doublet whose potential has its minimum
away from zero, at a vacuum value $v=246$\,GeV~\cite{higgs64,be64}. In that vacuum the W and Z bosons acquire masses proportional to $v$,
while one combination of the electroweak fields, the photon, stays massless. Quarks and charged leptons get their masses from
\emph{Yukawa couplings} $y_f$ to the doublet, $m_f=y_fv/\sqrt2$. The Higgs boson therefore couples to each fermion with strength $m_f/v$
and to the W and Z with $2m_V^2/v$.

\emph{The $\kappa$ framework.} Measurements of Higgs production and decay are summarised by \emph{coupling modifiers} $\kappa_i$, which
scale each Standard Model coupling, so that a rate involving particle $i$ scales as $\kappa_i^2$; the Standard Model has every
$\kappa_i=1$~\cite{david2012}.

\emph{Schemes and running.} A coupling or mass in quantum field theory is defined only together with a \emph{renormalisation scheme} and
an energy scale $\mu$. In the $\overline{\text{MS}}$ scheme only the divergent part of each loop is subtracted, and parameters run with
$\mu$. The \emph{on-shell} scheme instead fixes masses at the poles of the propagators and $\alpha$ at zero momentum transfer. The same
quantity can differ between the two by more than its measurement error, so every comparison states its scheme~\cite{erler2017,navas2024}.
The dependence on $\mu$ obeys the renormalisation-group equations (RGEs). At one loop
$\dd\alpha_i^{-1}/\dd\ln\mu=-b_i/2\pi$, with $(b_1,b_2,b_3)=(41/10,-19/6,-7)$ in the Standard Model (hypercharge normalised as in grand
unification). The strong coupling therefore weakens at high energy, and the hypercharge coupling grows. Precision work runs at two or
three loops~\cite{buttazzo2013}.

\emph{Dirac or Majorana.} A \emph{Dirac} mass joins a left-handed neutrino to a right-handed partner and conserves lepton number $L$. A
\emph{Majorana} mass joins the neutrino to its own antiparticle and changes $L$ by two. The lowest-dimension operator of Standard Model
fields that gives neutrinos mass is Weinberg's dimension-five operator, $(c/\Lambda)\,(LH)(LH)$, with $L$ the lepton doublet; in the Higgs
vacuum it gives a Majorana mass of order $cv^2/\Lambda$, small when the new scale $\Lambda$ is large~\cite{weinberg1979}. Majorana masses
can therefore also come from the Higgs vacuum.

\emph{Neutrinoless double-beta decay} ($0\nu\beta\beta$) turns two neutrons of a nucleus into two protons and two electrons with no
neutrinos. It changes $L$ by two and can occur only if neutrinos have a Majorana mass. Through light-neutrino exchange its rate is
proportional to $m_{\beta\beta}^2$, with the effective mass $m_{\beta\beta}=|\sum_iU_{ei}^2m_i|$ set by the mixing matrix $U$ and the
masses $m_i$~\cite{navas2024}. Xenon and germanium experiments search for it~\cite{kamlandzen2025,legend2025}.

\emph{Grand unification.} A grand unified theory (GUT) embeds the three gauge groups in one simple group, the smallest being
$SU(5)$~\cite{georgiglashow}, with a single coupling above a scale $M_X$. With the Standard Model's particles alone, the three running
couplings do not meet at one point~\cite{amaldi}. A GUT has two generic consequences. Its heavy gauge bosons let protons decay, with a
lifetime of order $M_X^4/(\alpha_G^2m_p^5)$; Super-Kamiokande's bound on $p\to e^+\pi^0$ excludes the minimal $SU(5)$
model~\cite{superk}. And its breaking leaves magnetic monopoles of mass near $M_X/\alpha_G$, about one per causal volume, which would
overfill the universe unless inflation dilutes them~\cite{preskill,guth}.
\end{primer}

\subsection{The spin ladder}\label{sec:spin}

GUID reads spin as the rungs of the Duality, and the Standard Model fills each rung \tagP:

\begin{center}\small
\begin{tabularx}{\linewidth}{@{}lLLL@{}}
\toprule
Spin & GUID reading & Particles & Physics \\
\midrule
0 & existence itself & Higgs field; the scalar $\phi$ & the Higgs vacuum value gives rest mass; $\phi$ sets $\alpha$ \\
1/2 & a duality, which can stay in balance & quarks, leptons & a massive Dirac fermion is two chiral halves joined by its mass \\
1 & information itself; delivered, an act & photon, gluons, W, Z & the carriers of every interaction \\
2 & the perfect duality, two in one & the frame's wave & a massless spin-2 field must couple to all energy: general relativity~\cite{weinberg64,deser70} \\
\bottomrule
\end{tabularx}
\end{center}

\noindent The rungs are told apart by rotation: a state of helicity $s$ picks up the phase $e^{is\theta}$ under a turn by $\theta$ about its
momentum, so a spin-1 quantum looks the same again after $360^\circ$, a spin-2 quantum after $180^\circ$, and a spin-1/2 state only after
$720^\circ$, returning with a minus sign after $360^\circ$. The ladder stops at spin~1 for quanta: no elementary spin-3/2 or spin-2 quantum is expected, and searches find none
(Section~\ref{sec:sm}). \emph{Verdict: holds as a reading}; the Standard Model says the same.

\subsection{The Higgs vacuum as the grant of existence}\label{sec:grant}

In GUID existence is granted by the ground. The Higgs field is the one field whose value in empty space is not zero: with potential
$V(H)=-\mu^2H^\dagger H+\lambda_H(H^\dagger H)^2$ it rests at $\langle H^\dagger H\rangle=v^2/2$, $v=\mu/\sqrt{\lambda_H}=246$\,GeV, and an
elementary particle has its rest mass by its coupling to that background~\cite{higgs64,be64} \tagP. GUID reads the background value, which
quantum field theory treats as a classical number, as G0's grant of rest mass, on the existence side and classical, in agreement with
Lemma~\ref{lem:classical}; the Higgs boson, a quantum ripple of the field that decays in about $10^{-22}$\,s, is content. The vacuum at the
minimum of its potential is the physical counterpart of the balanced state $\Phi_{\mathrm V}=0$ of~\eqref{eq:vibe}. Like G0 (``stated, not
derived''), the Higgs potential is an input of the Standard Model, not a derivation.

Two consequences can be tested (Section~\ref{sec:sm}). \emph{The couplings follow the grant.} If the vacuum grants every elementary
rest mass, each massive particle couples to the Higgs boson in proportion to its mass. \emph{Neutrinos are Dirac.} If, in addition,
lepton number is conserved, neutrino masses come from Yukawa couplings and neutrinoless double-beta decay does not occur at any rate. The
addition is needed because the Weinberg operator (Primer~\ref{box:particle}) gives a Majorana mass from the same vacuum. We add one reading, not derived: the difference $B-L$ of baryon and lepton number is one
of the counts every reader sees (Axiom~3), like electric charge, so it changes only when something carrying it crosses a boundary by a
recorded act. Neutrinoless double-beta decay would create two units of lepton number inside a closed nucleus with nothing carrying them
in, so it does not occur; proton decay, which conserves $B-L$, remains allowed. A signal at the next generation of experiments would refute
the statement. \emph{Verdict: live}, given the added reading.

The reading has one sharp limit, and it is gravitational. Gravity weighs all energy, granted and bound, while the Higgs field gives only
about 9\% of a nucleon's mass (Section~\ref{sec:fix}), so the Higgs field is not the gravitational field. And gravity is far too weak to give mass: its coupling to the electron,
$Gm_e^2/\hbar c=1.75\times10^{-45}$, and its share of the binding of hydrogen, $4.4\times10^{-40}$ of the electric one, are negligible
\chk{Q16}. \emph{Verdict: holds as a reading} of the Higgs vacuum as the grant; \emph{refuted} as an identity of the Higgs field with
gravity.

\subsection{The Higgs vacuum as the existence that separated out}\label{sec:higgsplit}

GUID makes the Higgs field the leading candidate for the field that grants existence, the one that separated out when the fields did.
Physics says when \tagP. Above about 160\,GeV the Higgs field rests at zero, and no elementary particle has rest mass; lattice calculations
find that, with the measured Higgs mass, the change is a smooth crossover at $T_c=159.5\pm1.5$\,GeV (D'Onofrio and
Rummukainen~\cite{donofrio}), about $10^{-11}$\,s after the start \chk{Q20}; the measured Higgs mass, 125.1\,GeV, lies well above the end
point of the line of first-order transitions \chk{data:D8.6}. GUID's posit, in its form that holds, has four parts.

\emph{Before the split, existence asserted itself through interaction with existence.} In the hot plasma before the crossover the particles
carry only thermal masses, of order coupling times temperature, which come entirely from their interactions with everything around them.
\emph{Verdict: holds as a reading.}

\emph{That became unstable.} As the universe cooled, the curvature of the Higgs potential at zero field turned negative and the field
rolled to its present value $v=246$\,GeV; with a 125\,GeV Higgs this is gradual, not a sudden phase transition. \emph{Verdict: holds}, as a
crossover.

\emph{The information field split out of the electroweak field.} At the split the photon emerges as the fixed mixture
$A=\sin\theta_W\,W^3+\cos\theta_W\,B$ of the two electroweak gauge fields, and the electric charge is born as $e=g\sin\theta_W$, so that $\alpha$
appears near $1/128$ at the Z mass \chk{Q20} (Section~\ref{sec:sm}). \emph{Verdict: holds}, exactly.

\emph{The Higgs field became needed then.} At that moment the W and Z take their masses from the vacuum, $m_W=gv/2\approx80$\,GeV and
$m_Z=m_W/\cos\theta_W\approx91$\,GeV at tree level \chk{Q20}, within about half a per cent of the measured values, depending on the
scheme (Section~\ref{sec:data}), while the photon stays massless:
the Higgs vacuum grants existence to the bosons that need it at the very moment the information field emerges. \emph{Verdict: holds.} The
Standard Model Higgs belongs to this split, the weak from the electromagnetic; a Higgs-type field for the earlier split of the strong force,
if that split was a unification breaking (Hypothesis~6$'$), would be a second, hypothetical field.

\emph{An open question.} In Higgs inflation the same field also drives the start: with a coupling $\xi H^\dagger HR$ of the Higgs to
curvature, $\xi$ of order $10^4$, it fits the observed spectrum (Bezrukov and Shaposhnikov~\cite{bezrukov}), which would make the field
that grants existence the field of the start itself. As for $R^2$ inflation (Section~\ref{sec:cosmo}), the coupling to curvature is the
form Proposition~\ref{pr:gamma} excludes for the frame after the start, it disappears in the Einstein frame, and within the start the
axioms do not constrain it. Whether Axiom~1 says more is open.

\subsection{The forces as tiers of one interaction}

GUID reads the forces as tiers of one interaction: the strong force binds the cluster, the weak force keeps it stable, electromagnetism
links clusters of clusters. Each does that job \tagP: the strong force binds quarks into nucleons and nucleons into nuclei, the weak force
moves nuclei towards the valley of stability, electromagnetism binds atoms and molecules. The weak and electromagnetic forces are one
electroweak interaction~\cite{glashow,weinberg67,salam}. Whether all three meet at one energy is open: in the Standard Model the couplings
cross pairwise between $10^{13}$ and $3\times10^{16}$\,GeV and do not meet at one point~\cite{amaldi}, \chk{sm:SM4}. Gravity is not among the tiers in GUID: it is the
frame the tiers act in. \emph{Verdict: holds as a reading}; unification \emph{untested}. In GUID's timeline the
strong force splits from the electroweak force at the start (Hypothesis~6$'$), and the electroweak field splits into the weak and the
electromagnetic force about $10^{-11}$\,s later, when the Higgs vacuum appears (Section~\ref{sec:higgsplit}).

% GUID v6, Primer P10 (label box:hadron): hadron spectroscopy. Background only; it states no result of the paper.
% Place it before "Hadrons as trees of dualities" (marker %PRIMER:hadron in part1.tex): % GUID v6, Primer P10 (label box:hadron): hadron spectroscopy. Background only; it states no result of the paper.
% Place it before "Hadrons as trees of dualities" (marker %PRIMER:hadron in part1.tex): % GUID v6, Primer P10 (label box:hadron): hadron spectroscopy. Background only; it states no result of the paper.
% Place it before "Hadrons as trees of dualities" (marker %PRIMER:hadron in part1.tex): \input{primers/p10-hadron}
\begin{primer}{Hadron spectroscopy}{box:hadron}
\emph{Colour.} Quarks carry the colour charge of $SU(3)$: a quark is a triplet $\mathbf3$, an antiquark an antitriplet $\mathbf{\bar3}$.
Only colour \emph{singlets} exist as free particles, the hadrons~\cite{navas2024}. The products
$\mathbf3\otimes\mathbf{\bar3}=\mathbf1\oplus\mathbf8$ and $\mathbf3\otimes\mathbf3=\mathbf{\bar3}\oplus\mathbf6$ give the ordinary
hadrons: a meson $q\bar q$ in the singlet, and a baryon $qqq$, whose first two quarks couple to $\mathbf{\bar3}$ and then with the third to
a singlet. One-gluon exchange attracts in the singlet $q\bar q$ and, half as strongly, in the $\mathbf{\bar3}$ $qq$, and repels in the
$\mathbf8$ and the $\mathbf6$.

\emph{Diquarks.} A \emph{diquark} is a pair of quarks in the attractive $\mathbf{\bar3}$. The ``good'' diquark, with spin 0 and a flavour-
antisymmetric pair such as $[ud]$, binds most; the ``bad'' one has spin 1~\cite{jaffe2005}.

\emph{Compact or molecular.} A hadron with more than three quarks, an \emph{exotic}, can be \emph{compact}, for example a diquark and an
antidiquark bound by colour forces, $[qq][\bar q\bar q]$~\cite{maiani2005}. Or it can be a \emph{hadronic molecule}: two colour-singlet
hadrons held together by the residual strong force, like the deuteron~\cite{guo2018}. Molecules sit just below the threshold
$m_1+m_2$ of their two constituents.

\emph{The colour-spin interaction.} In the constituent-quark model a hadron's mass is the sum of effective quark masses plus a hyperfine
term proportional to $(\boldsymbol\lambda_i\cdot\boldsymbol\lambda_j)(\boldsymbol\sigma_i\cdot\boldsymbol\sigma_j)/m_im_j$ for each pair,
where the $\boldsymbol\lambda$ are colour and the $\boldsymbol\sigma$ spin matrices. This chromomagnetic interaction of De R\'ujula, Georgi
and Glashow explains, for example, why the $\Delta$ is heavier than the nucleon~\cite{derujula1975}. Karliner and Rosner (KR) refined it
with separate effective quark masses in mesons and in baryons, the baryon masses for every constituent of a compact multiquark, and binding
terms for heavy pairs with $B(QQ)=B(Q\bar Q)/2$; with it they predicted the doubly charmed baryon and doubly heavy
tetraquarks~\cite{karliner2014,karliner2017}.

\emph{Thresholds and cusps.} Where a new two-hadron channel opens, at $M=m_1+m_2$, the scattering amplitude has a branch point. A
\emph{cusp} can then appear in a mass spectrum without any state, and a state that couples to the channel is pulled towards the threshold.
Many exotic candidates lie within a few MeV of a threshold~\cite{guo2018}.

\emph{Names.} The Particle Data Group names hadrons by their quantum numbers, not by their structure~\cite{navas2024}; the
$X(3872)$, for example, is listed as $\chi_{c1}(3872)$. LHCb's convention for exotics~\cite{gershon2022} uses $T$ for four-quark and $P$
for five-quark states, with subscripts for heavy flavour, $\psi$ standing for hidden charm $c\bar c$. So $T_{cc}^+$ is $cc\bar u\bar d$,
seen just below the $D^{*+}D^0$ threshold~\cite{lhcb2022h}, and the hidden-charm pentaquark $P_c(4312)^+$~\cite{lhcb2019h} becomes
$P_\psi^N(4312)^+$.

\emph{String junctions.} In a baryon the three colour flux tubes meet at a point, a \emph{string junction}~\cite{rossiveneziano1977}. A
compact tetraquark has two junctions and a compact pentaquark three, and some models add a constant energy per
junction~\cite{karliner2017a}.
\end{primer}

\begin{primer}{Hadron spectroscopy}{box:hadron}
\emph{Colour.} Quarks carry the colour charge of $SU(3)$: a quark is a triplet $\mathbf3$, an antiquark an antitriplet $\mathbf{\bar3}$.
Only colour \emph{singlets} exist as free particles, the hadrons~\cite{navas2024}. The products
$\mathbf3\otimes\mathbf{\bar3}=\mathbf1\oplus\mathbf8$ and $\mathbf3\otimes\mathbf3=\mathbf{\bar3}\oplus\mathbf6$ give the ordinary
hadrons: a meson $q\bar q$ in the singlet, and a baryon $qqq$, whose first two quarks couple to $\mathbf{\bar3}$ and then with the third to
a singlet. One-gluon exchange attracts in the singlet $q\bar q$ and, half as strongly, in the $\mathbf{\bar3}$ $qq$, and repels in the
$\mathbf8$ and the $\mathbf6$.

\emph{Diquarks.} A \emph{diquark} is a pair of quarks in the attractive $\mathbf{\bar3}$. The ``good'' diquark, with spin 0 and a flavour-
antisymmetric pair such as $[ud]$, binds most; the ``bad'' one has spin 1~\cite{jaffe2005}.

\emph{Compact or molecular.} A hadron with more than three quarks, an \emph{exotic}, can be \emph{compact}, for example a diquark and an
antidiquark bound by colour forces, $[qq][\bar q\bar q]$~\cite{maiani2005}. Or it can be a \emph{hadronic molecule}: two colour-singlet
hadrons held together by the residual strong force, like the deuteron~\cite{guo2018}. Molecules sit just below the threshold
$m_1+m_2$ of their two constituents.

\emph{The colour-spin interaction.} In the constituent-quark model a hadron's mass is the sum of effective quark masses plus a hyperfine
term proportional to $(\boldsymbol\lambda_i\cdot\boldsymbol\lambda_j)(\boldsymbol\sigma_i\cdot\boldsymbol\sigma_j)/m_im_j$ for each pair,
where the $\boldsymbol\lambda$ are colour and the $\boldsymbol\sigma$ spin matrices. This chromomagnetic interaction of De R\'ujula, Georgi
and Glashow explains, for example, why the $\Delta$ is heavier than the nucleon~\cite{derujula1975}. Karliner and Rosner (KR) refined it
with separate effective quark masses in mesons and in baryons, the baryon masses for every constituent of a compact multiquark, and binding
terms for heavy pairs with $B(QQ)=B(Q\bar Q)/2$; with it they predicted the doubly charmed baryon and doubly heavy
tetraquarks~\cite{karliner2014,karliner2017}.

\emph{Thresholds and cusps.} Where a new two-hadron channel opens, at $M=m_1+m_2$, the scattering amplitude has a branch point. A
\emph{cusp} can then appear in a mass spectrum without any state, and a state that couples to the channel is pulled towards the threshold.
Many exotic candidates lie within a few MeV of a threshold~\cite{guo2018}.

\emph{Names.} The Particle Data Group names hadrons by their quantum numbers, not by their structure~\cite{navas2024}; the
$X(3872)$, for example, is listed as $\chi_{c1}(3872)$. LHCb's convention for exotics~\cite{gershon2022} uses $T$ for four-quark and $P$
for five-quark states, with subscripts for heavy flavour, $\psi$ standing for hidden charm $c\bar c$. So $T_{cc}^+$ is $cc\bar u\bar d$,
seen just below the $D^{*+}D^0$ threshold~\cite{lhcb2022h}, and the hidden-charm pentaquark $P_c(4312)^+$~\cite{lhcb2019h} becomes
$P_\psi^N(4312)^+$.

\emph{String junctions.} In a baryon the three colour flux tubes meet at a point, a \emph{string junction}~\cite{rossiveneziano1977}. A
compact tetraquark has two junctions and a compact pentaquark three, and some models add a constant energy per
junction~\cite{karliner2017a}.
\end{primer}

\begin{primer}{Hadron spectroscopy}{box:hadron}
\emph{Colour.} Quarks carry the colour charge of $SU(3)$: a quark is a triplet $\mathbf3$, an antiquark an antitriplet $\mathbf{\bar3}$.
Only colour \emph{singlets} exist as free particles, the hadrons~\cite{navas2024}. The products
$\mathbf3\otimes\mathbf{\bar3}=\mathbf1\oplus\mathbf8$ and $\mathbf3\otimes\mathbf3=\mathbf{\bar3}\oplus\mathbf6$ give the ordinary
hadrons: a meson $q\bar q$ in the singlet, and a baryon $qqq$, whose first two quarks couple to $\mathbf{\bar3}$ and then with the third to
a singlet. One-gluon exchange attracts in the singlet $q\bar q$ and, half as strongly, in the $\mathbf{\bar3}$ $qq$, and repels in the
$\mathbf8$ and the $\mathbf6$.

\emph{Diquarks.} A \emph{diquark} is a pair of quarks in the attractive $\mathbf{\bar3}$. The ``good'' diquark, with spin 0 and a flavour-
antisymmetric pair such as $[ud]$, binds most; the ``bad'' one has spin 1~\cite{jaffe2005}.

\emph{Compact or molecular.} A hadron with more than three quarks, an \emph{exotic}, can be \emph{compact}, for example a diquark and an
antidiquark bound by colour forces, $[qq][\bar q\bar q]$~\cite{maiani2005}. Or it can be a \emph{hadronic molecule}: two colour-singlet
hadrons held together by the residual strong force, like the deuteron~\cite{guo2018}. Molecules sit just below the threshold
$m_1+m_2$ of their two constituents.

\emph{The colour-spin interaction.} In the constituent-quark model a hadron's mass is the sum of effective quark masses plus a hyperfine
term proportional to $(\boldsymbol\lambda_i\cdot\boldsymbol\lambda_j)(\boldsymbol\sigma_i\cdot\boldsymbol\sigma_j)/m_im_j$ for each pair,
where the $\boldsymbol\lambda$ are colour and the $\boldsymbol\sigma$ spin matrices. This chromomagnetic interaction of De R\'ujula, Georgi
and Glashow explains, for example, why the $\Delta$ is heavier than the nucleon~\cite{derujula1975}. Karliner and Rosner (KR) refined it
with separate effective quark masses in mesons and in baryons, the baryon masses for every constituent of a compact multiquark, and binding
terms for heavy pairs with $B(QQ)=B(Q\bar Q)/2$; with it they predicted the doubly charmed baryon and doubly heavy
tetraquarks~\cite{karliner2014,karliner2017}.

\emph{Thresholds and cusps.} Where a new two-hadron channel opens, at $M=m_1+m_2$, the scattering amplitude has a branch point. A
\emph{cusp} can then appear in a mass spectrum without any state, and a state that couples to the channel is pulled towards the threshold.
Many exotic candidates lie within a few MeV of a threshold~\cite{guo2018}.

\emph{Names.} The Particle Data Group names hadrons by their quantum numbers, not by their structure~\cite{navas2024}; the
$X(3872)$, for example, is listed as $\chi_{c1}(3872)$. LHCb's convention for exotics~\cite{gershon2022} uses $T$ for four-quark and $P$
for five-quark states, with subscripts for heavy flavour, $\psi$ standing for hidden charm $c\bar c$. So $T_{cc}^+$ is $cc\bar u\bar d$,
seen just below the $D^{*+}D^0$ threshold~\cite{lhcb2022h}, and the hidden-charm pentaquark $P_c(4312)^+$~\cite{lhcb2019h} becomes
$P_\psi^N(4312)^+$.

\emph{String junctions.} In a baryon the three colour flux tubes meet at a point, a \emph{string junction}~\cite{rossiveneziano1977}. A
compact tetraquark has two junctions and a compact pentaquark three, and some models add a constant energy per
junction~\cite{karliner2017a}.
\end{primer}

\subsection{Hadrons as trees of dualities}\label{sec:trees}

Axiom~1 makes every part a Duality of two halves, and Axiom~3 makes every layer a tree rooted at its own Duality \tagD. Read for the
strong force, which binds the cluster, the two give one more hypothesis.

\begin{hypothesis}[hadrons as trees of dualities]\label{h:tree}
Every hadron is a binary tree of two colour-conjugate halves, each internal node a bound duality: a meson is $(q)(\bar q)$, a baryon
$(q)[qq]$, a tetraquark $[qq][\bar q\bar q]$ (compact) or $(q\bar q)(q\bar q)$ (molecular), a pentaquark $[qq][qq]\bar q$ or
$(qqq)(q\bar q)$. \tagH
\end{hypothesis}

\noindent We call it H8 for short. QCD's colour algebra already allows every colour singlet to be written as such a tree, so the statement as such holds as a
reading. Its content lies in the masses a colour-spin model gives the compact trees, in how close the molecular ones lie to two-hadron
thresholds, and in the partner states a tree implies. Section~\ref{sec:exotics} tests all three with a standard colour-spin model on
the compact trees, fitted to ordinary hadrons only. \emph{Verdict: holds as a reading}; the tests of its content are reported in
Section~\ref{sec:exotics}.

G0 gives a sharper reading of the five-quark hadrons. The dual gate cannot witness its own first instance: ``somebody has to trust the
first two witnesses'', and the minting makes at least five parts the holders of the root keys (Section~\ref{sec:dmp}).

\begin{hypothesis}[the minting in hadrons]\label{h:mint}
A pentaquark is a minting: a bound triple of colour-bearers that together carry all three colours, and the first two witnesses, a quark
and its own antiquark, so that the pentaquark's antiquark has a partner quark of its own flavour. \tagH
\end{hypothesis}

\noindent We call it H9. It reads the structure of G0 at the scale of one hadron; it does not identify a pentaquark with the start of
Hypothesis~\ref{h:minting}, which is G0 for the universe. QCD requires the three colours and some quark-antiquark pair in every pentaquark; it does not require the pair
to be of one flavour. H9 does, so it forbids pentaquarks whose antiquark has no partner, such as the $\Theta^+$ ($uudd\bar s$) or
$ccuu\bar d$, but not a doubly charmed state such as $ccud\bar u$, and it predicts hidden-bottom pentaquarks ($uudb\bar b$) as partners of
the observed hidden-charm ones. It was conceived after seeing that every established
pentaquark carries a $c\bar c$ pair, so those states are not evidence for it; the test lies in the states not yet established
(Section~\ref{sec:exotics}).

\section{Readings that add no prediction}\label{sec:readings}

Several of GUID's identifications reproduce established physics without adding anything a measurement could tell apart. They show that
the reading is consistent, not that it is right, and we collect them in one table; Appendix~\ref{app:readings} gives each in full.

\begin{table}[tb]\centering\small
\caption{Readings that agree with established physics and add no prediction. Verdict: \emph{holds as a reading} for each.}
\label{tab:readings}
\begin{tabularx}{\linewidth}{@{}lLl@{}}
\toprule
Reading & What it reproduces & Check \\
\midrule
The existence a boundary asserts, $gh/c$ & $4\pi^2$ times the Unruh temperature; the Hawking temperature at a horizon & \chk{Q1} \\
Time slows where the frame is strong & the lapse; the Pound--Rebka and clock redshifts & \chk{Q2} \\
Einstein's equation as an equation of state & Jacobson's derivation from the Clausius relation; the constant is open & \eqref{eq:jacobson} \\
Entropy as the cost of converting existence and information & Landauer's principle; DMP's entropy floor (T31) & \\
The containment split of a horizon & Christodoulou's mass formula for the Kerr hole & \eqref{eq:christo}, \chk{Q10} \\
The early universe a Kugelblitz & the Hubble mass at its own Schwarzschild radius, at every epoch & \chk{Q12} \\
Inflation as the possibility space of existence & the flattening of curvature over about 60 e-folds & \chk{Q13} \\
$\alpha$ as the VIBE coefficient & $\alpha$ as the self-interaction strength of charge; $a_e=\alpha/2\pi$ & \chk{Q16} \\
\bottomrule
\end{tabularx}
\end{table}


\section{Relation to the approaches to quantum gravity}\label{sec:landscape}

\begin{center}\small
\begin{longtable}{@{}P{0.18\linewidth}P{0.22\linewidth}P{0.28\linewidth}P{0.22\linewidth}@{}}
\caption{The approaches to quantum gravity and GUID.}\label{tab:landscape}\\
\toprule
Approach & What is quantum & Relation to GUID & What tells them apart \\
\midrule
\endfirsthead
\toprule
Approach & What is quantum & Relation to GUID & What tells them apart \\
\midrule
\endhead
Effective field theory~\cite{donoghue,bdh} & metric fluctuations below $E_P$ & GUID keeps the classical part; the graviton-loop term
of~\eqref{eq:eft} is absent & a term $10^{-70}$ of Newton's at 1\,m \chk{Q1}: unmeasurable \\
String theory, AdS/CFT~\cite{maldacena} & everything; geometry emerges from a quantum boundary & hostable as content; the code kinematics
of Section~\ref{sec:holo} are shared, with a classical centre & entanglement of masses, if decoherence $<\omega/2$ \\
Loop quantum gravity~\cite{ashtekar,rovellismolin} & the triad/tetrad; area has a discrete spectrum & the frame is LQG's own variable, not
quantised; area a number & area spectrum (not accessible); BMV if decoherence $<\omega/2$ \\
Asymptotic safety~\cite{weinberg79,reuter} & the metric, with a UV fixed point & hostable & BMV if decoherence $<\omega/2$ \\
Causal dynamical triangulations~\cite{cdt} & sum over geometries & allowed within the start only & BMV if decoherence $<\omega/2$ \\
Causal sets~\cite{bombelli} & discrete order; $\Lambda\sim1/\sqrt N$ & DMP's order is a tree, not a sprinkled causal set & Lorentz
invariance needs random sprinkling \\
Thermodynamic gravity~\cite{jacobson95} & the microscopic entropy only & GUID's thermal reading, Section~\ref{sec:thermo} & none
beyond GR \\
Classical-quantum gravity~\cite{oppenheim,osswd} & matter only & GUID's dynamics, Corollary~\ref{co:diffuse} & diffusion against
decoherence \\
Channel models~\cite{ktm,tilloydiosi} & matter only & GUID's architecture, Proposition~\ref{pr:ktm} & noise and decoherence rates \\
Gravity-related collapse~\cite{diosi87,penrose96} & matter, collapsing & Bridge~B lands on its rate; nuclear-scale smearing excluded,
$R_0\ge4\times10^{-10}$\,m allowed~\cite{donadi,figurato} & decoherence of a $10^{-14}$\,kg superposition (Section~\ref{sec:bmv}) \\
Semiclassical gravity~\cite{moller,rosenfeld63} & matter; metric sourced by $\langle T\rangle$ & rejected: nonlinear; Page--Geilker &
done~\cite{pagegeilker} \\
Modified dynamics (MOND)~\cite{milgrom} & nothing; a law at low acceleration & the stochastic remnant gives a term of the size $a_0$; tested in Sections~\ref{sec:sparc} and~\ref{sec:remnant} & lensing,
acoustic peaks, merging clusters, planetary ephemerides \\
Spacetime from entanglement~\cite{vanraamsdonk,ermepr} & content; geometry emergent & GUID's frame as the classical centre of such a geometry is
open & unclear \\
\bottomrule
\end{longtable}
\end{center}


\section{What Part I predicts and where it is tested}\label{sec:plan}

Table~\ref{tab:plan} closes Part I with every prediction that an experiment or an observation can reach, the alternative it must be told
apart from, and where Parts II and III test it. Part IV collects the verdicts.

\begin{center}\small
\begin{longtable}{@{}P{0.3\linewidth}P{0.3\linewidth}P{0.22\linewidth}P{0.1\linewidth}@{}}
\caption{The test plan: predictions of Part I, their alternatives and where they are tested.}\label{tab:plan}\\
\toprule
Observable & GUID & Alternative & Tested in \\
\midrule
\endfirsthead
\toprule
Observable & GUID & Alternative & Tested in \\
\midrule
\endhead
\bottomrule
\endfoot
\multicolumn{4}{@{}l}{\emph{The frame}} \\
Free fall; the information scalar & universal; $\eta=2\lambda^2f_\oplus\Delta f$, $|\lambda|$ bounded & universal & Sec.~\ref{sec:data} \\
Post-Newtonian $\gamma$, $\beta$, Nordtvedt $\eta$ & as in general relativity & as GR & Sec.~\ref{sec:data} \\
Gravitational waves & speed $c$, massless, two tensor modes & as GR & Sec.~\ref{sec:data} \\
Entanglement of two masses through gravity & none, $\mathcal N=0$ & quantised metric: $\mathcal N=\tfrac12|\sin(\Delta\varphi/2)|$ & Sec.~\ref{sec:bmv} \\
Decoherence of mesoscopic superpositions (Bridge~B) & rate $GM^2/\hbar R$ & none required & Secs.~\ref{sec:bmv}, \ref{sec:data} \\
LOCC inequality, massive oscillators & satisfied & can be violated & proposed \\
Decoherence against diffusion & both non-zero, trade-off at equality & neither required & Sec.~\ref{sec:gravsim} \\
\multicolumn{4}{@{}l}{\emph{Black holes and holography}} \\
Kerr holes & efficiency $\le42.3\%$; $|a|\le1$; area never shrinks & as GR & Sec.~\ref{sec:data} \\
Holographic assembly & minimal cuts, reconstruction, classical centre & quantum code kinematics & Sec.~\ref{sec:holocode} \\
\multicolumn{4}{@{}l}{\emph{Cosmology}} \\
Late universe & de Sitter floor, $w=-1$ & evolving dark energy & Sec.~\ref{sec:data} \\
Scalar tilt (plateau inflation, borrowed) & $n_s=1-2/N$ & other potentials & Sec.~\ref{sec:data} \\
Primordial tensor modes & standard spectrum, made within the start & $r=16\epsilon$ & Sec.~\ref{sec:data} \\
Relic monopoles; proton decay & none; possible from a unification before the start & --- & Secs.~\ref{sec:data}, \ref{sec:sm} \\
\multicolumn{4}{@{}l}{\emph{Dark matter}} \\
Stochastic remnant; Bridge~A & flat curves, $v^4\propto M$, scale from the horizon & halos; RAR & Secs.~\ref{sec:sparc}, \ref{sec:remnant} \\
Contained remnant & primordial black holes, asteroid-mass window & particle dark matter & Sec.~\ref{sec:remnant} \\
\multicolumn{4}{@{}l}{\emph{The information scalar and matter}} \\
Drift of $\alpha$ & zero or negative & none predicted & Sec.~\ref{sec:data} \\
Variation of $\mu=m_p/m_e$ with $\alpha$ & small, $|R|$ below a few hundredths & unification: $|R|=20$ to 50 & Sec.~\ref{sec:sm} \\
Higgs couplings & proportional to mass & Standard Model & Sec.~\ref{sec:sm} \\
Neutrinoless double-beta decay & none (with $B-L$ as a count) & Majorana masses & Sec.~\ref{sec:sm} \\
Unification of the gauge couplings & none within the Standard Model's running & grand unification & Sec.~\ref{sec:sm} \\
Elementary spin & 0, 1/2 and 1 only & Standard Model & Sec.~\ref{sec:sm} \\
Exotic hadrons (H8) & compact-tree masses; molecules near thresholds; partners & QCD models & Sec.~\ref{sec:exotics} \\
Pentaquarks (H9) & hidden flavour only; hidden-bottom states exist & open flavour allowed & Sec.~\ref{sec:exotics} \\
Hawking radiation of real holes & thermal; return of content open & Page curve & not observable \\
\end{longtable}
\end{center}

